%% titre: Introduction à la programmation linéaire
%% matiere: Optimisation
%% niveau: Licence 3
%% annee: 2025-2026
%% ordre: 1
%% description: Cours d'optimisation licence 3 2025-2026.
\documentclass[a4paper,12pt,openany]{book}

% ========================================
% ENCODAGE ET LANGUE
% ========================================
\usepackage[utf8]{inputenc}
\usepackage[T1]{fontenc}
\usepackage[french]{babel}
\usepackage{lmodern}

% ========================================
% MISE EN PAGE ET GÉOMÉTRIE
% ========================================
\usepackage{geometry}
\geometry{top=2cm, bottom=2cm, left=1cm, right=1cm}
\usepackage{fancyhdr}
\usepackage{multicol}

% ========================================
% COULEURS
% ========================================
\usepackage[table]{xcolor}
\usepackage{colortbl}

% Couleurs personnalisées
\definecolor{maincolor}{RGB}{41, 128, 185}
\definecolor{accentcolor}{RGB}{236, 240, 241}
\definecolor{darkcolor}{RGB}{44, 62, 80}
\definecolor{titlecolor}{RGB}{0, 51, 102}
\definecolor{linecolor}{RGB}{70, 130, 180}
\definecolor{titleblue}{RGB}{0, 70, 140}
\definecolor{bggray}{RGB}{245, 245, 245}
\definecolor{keywordred}{RGB}{180, 0, 0}
\definecolor{vargreen}{RGB}{0, 120, 0}
\definecolor{mathblue}{RGB}{0, 0, 120}
\definecolor{headerbg}{RGB}{0, 84, 159}
\definecolor{rowbg}{RGB}{240, 247, 255}

% ========================================
% GRAPHIQUES ET IMAGES
% ========================================
\usepackage{graphicx}
\usepackage[export]{adjustbox}
\usepackage{float}
\usepackage{caption}
\graphicspath{{./images/}}

% ========================================
% MATHÉMATIQUES
% ========================================
\usepackage{amsmath,amssymb,amsthm,amsfonts}
\usepackage{dsfont, bm, esint, pifont}
\usepackage{stmaryrd}
\usepackage{bbold}
\usepackage[version=4]{mhchem}

% ========================================
% TIKZ ET GRAPHIQUES
% ========================================
\usepackage{tikz}
\usetikzlibrary{patterns,tikzmark}
\usepackage{tkz-tab}
\usepackage{pgfplots}
\pgfplotsset{compat=1.17}

% ========================================
% TABLEAUX
% ========================================
\usepackage{array,booktabs,multirow}

% ========================================
% BOÎTES ET CADRES
% ========================================
\usepackage[most]{tcolorbox}
\tcbuselibrary{theorems,skins,breakable}
\usepackage{mdframed,framed}

% Style de l'encadré coloré
\newenvironment{colorframed}{%
	\def\FrameCommand{\fboxsep=10pt \fcolorbox{black}{bggray}}%
	\MakeFramed{\advance\hsize-\width \FrameRestore}%
}{\endMakeFramed}

% ========================================
% LISTES
% ========================================
\usepackage{enumitem}

% ========================================
% FILIGRANE
% ========================================
\usepackage{draftwatermark}
\SetWatermarkText{Dr. TOUGMA}
\SetWatermarkScale{0.8}
\SetWatermarkColor[gray]{0.9}

% ========================================
% HYPERLIENS
% ========================================
\usepackage{hyperref}
\hypersetup{
	colorlinks=true,
	linkcolor=blue,
	filecolor=magenta,
	urlcolor=cyan,
	pdftitle={Livre d'Optimisation},
	pdfpagemode=FullScreen,
}

% ========================================
% STYLE DES TITRES DE CHAPITRES
% ========================================
\makeatletter
\newcommand{\thechapterwords}{%
	\ifcase\thechapter\or 1\or 2\or 3\or 4\or
	5\or 6\or 7\or 8\or 9\or 10\or 11\fi}
\def\thickhrulefill{\leavevmode \leaders \hrule height 1ex \hfill \kern \z@}
\def\@makechapterhead#1{%
	\vspace*{10\p@}%
	{\parindent \z@ \centering \reset@font
		\thickhrulefill\quad
		\scshape \@chapapp{} \thechapterwords
		\quad \thickhrulefill
		\par\nobreak
		\vspace*{10\p@}%
		\interlinepenalty\@M
		\hrule
		\vspace*{10\p@}%
		\Huge \bfseries #1\par\nobreak
		\vspace*{10\p@}%
		\hrule
		\vskip 35\p@}}
\makeatother

% ========================================
% EN-TÊTES ET PIEDS DE PAGE
% ========================================
\pagestyle{fancy}
\setlength{\headheight}{15pt}
\fancyhf{}
\renewcommand{\headrulewidth}{0.4pt}
\renewcommand{\footrulewidth}{0.4pt}
\fancyhead[L]{ULBO}
\fancyhead[C]{Optimisation}
\fancyhead[R]{L3S5 Mathématiques}
\fancyfoot[L]{Dr TOUGMA}
\fancyfoot[C]{Page \thepage}
\fancyfoot[R]{\copyright\, 2025-2026}

% ========================================
% ENVIRONNEMENTS THÉORÈMES
% ========================================
\theoremstyle{plain}
\newtheorem{theorem}{Théorème}[section]
\newtheorem{proposition}[theorem]{Proposition}
\newtheorem{lemma}[theorem]{Lemme}
\newtheorem{corollary}[theorem]{Corollaire}

\theoremstyle{definition}
\newtheorem{definition}[theorem]{Définition}
\newtheorem{example}[theorem]{Exemple}
\newtheorem{exemple}[theorem]{Exemple}
\newtheorem{hypothese}[theorem]{Hypothèse}

\theoremstyle{remark}
\newtheorem{remark}[theorem]{Remarque}
\newtheorem{remarque}[theorem]{Remarque}
\newtheorem{exercise}{Exercice}[section]
\newtheorem{notation}[theorem]{Notation}
\newtheorem{solution}{Solution}[section]
%\newtheorem{prooff}{Démonstration}[section]

% Style NB
\makeatletter
\newtheoremstyle{nbstyle}
{\topsep}{\topsep}{}{}{\bfseries}{}{ }{\thmname{#1}}
\makeatother
\theoremstyle{nbstyle}
\newtheorem*{NB}{NB:}

% ========================================
% HABILLAGE COLORÉ AVEC TCOLORBOX
% ========================================
\tcolorboxenvironment{theorem}{
	colback=red!5!white, colframe=red!50!black,
	fonttitle=\bfseries, boxrule=0.8pt, arc=3pt,
	before skip=10pt, after skip=10pt, breakable
}
\tcolorboxenvironment{proposition}{
	colback=red!5!white, colframe=red!50!black,
	fonttitle=\bfseries, boxrule=0.8pt, arc=3pt,
	before skip=10pt, after skip=10pt, breakable
}
\tcolorboxenvironment{lemma}{
	colback=red!5!white, colframe=red!50!black,
	fonttitle=\bfseries, boxrule=0.8pt, arc=3pt,
	before skip=10pt, after skip=10pt, breakable
}
\tcolorboxenvironment{corollary}{
	colback=red!5!white, colframe=red!50!black,
	fonttitle=\bfseries, boxrule=0.8pt, arc=3pt,
	before skip=10pt, after skip=10pt, breakable
}
\tcolorboxenvironment{definition}{
	colback=blue!5!white, colframe=blue!50!black,
	fonttitle=\bfseries, boxrule=0.8pt, arc=3pt,
	before skip=10pt, after skip=10pt, breakable
}
\tcolorboxenvironment{example}{
	colback=green!5!white, colframe=green!50!black,
	fonttitle=\bfseries, boxrule=0.8pt, arc=3pt,
	before skip=10pt, after skip=10pt, breakable
}
\tcolorboxenvironment{exemple}{
	colback=green!5!white, colframe=green!50!black,
	fonttitle=\bfseries, boxrule=0.8pt, arc=3pt,
	before skip=10pt, after skip=10pt, breakable
}
\tcolorboxenvironment{remark}{
	colback=yellow!5!white, colframe=orange!50!black,
	fonttitle=\bfseries, boxrule=0.8pt, arc=3pt,
	before skip=10pt, after skip=10pt, breakable
}
\tcolorboxenvironment{remarque}{
	colback=yellow!5!white, colframe=orange!50!black,
	fonttitle=\bfseries, boxrule=0.8pt, arc=3pt,
	before skip=10pt, after skip=10pt, breakable
}
\tcolorboxenvironment{exercise}{
	colback=purple!5!white, colframe=purple!70!black,
	fonttitle=\bfseries, boxrule=0.8pt, arc=3pt,
	before skip=10pt, after skip=10pt, breakable
}

% ========================================
% COMMANDES PERSONNALISÉES
% ========================================
\newcommand{\R}{\mathbb{R}}
\newcommand{\N}{\mathbb{N}}
\newcommand{\norm}[1]{\left\lVert#1\right\rVert}
\newcommand{\abs}[1]{\left\lvert#1\right\rvert}
\newcommand{\grad}{\nabla}
\newcommand{\hess}{\nabla^2}
\newcommand{\T}{^T}

% ========================================
% NOMS FRANÇAIS
% ========================================
\addto\captionsfrench{\renewcommand{\chaptername}{Chapitre}}
\addto\captionsfrench{\renewcommand{\tablename}{Tableau}}
\addto\captionsfrench{\renewcommand{\contentsname}{Table des matières}}
\addto\captionsfrench{\renewcommand{\figurename}{Figure}}
\begin{document}
	
	% Suppression des numéros de page pour la page de garde
	\thispagestyle{empty}
	
	\begin{center}
		% Logo de l'université (à remplacer par votre chemin)
		\vspace*{1cm}
		\includegraphics[width=0.4\textwidth]{ULBO}
		
		\vspace{2cm}
		
		% Titre principal
		{\color{titlecolor}\Huge\textbf{Optimisation}\par}
		
		\vspace{0.5cm}
		
		% Ligne décorative
		\begin{tikzpicture}
		\draw[thick, linecolor] (0,0) -- (12,0);
		\end{tikzpicture}
		
		\vspace{1cm}
		
		% Sous-titre
		{\Large\textsc{UNIVERSITÉ LÉDÉA BERNARD OUÉDRAOGO}\par}
		
		\vspace{2cm}
		
		% Informations sur le cours
		{\Large
			\textbf{Licence 3 Mathématiques}\\
			\vspace{0.3cm}
			Semestre 5 - 2025/2026\\
		}
		{\large
			\vspace{0.9cm}
			\textbf{ CM: 20h - TD: 06h}
		}
		
		\vspace{2.5cm}
		
		% Informations sur l'étudiant/professeur
		{\large
			\textbf{Dr. A. TOUGMA}
		}
		
		\vfill
		
		% Département/Institution
		{\large
			Enseignants-Chercheurs; Assistants en mathématiques appliquées\\
			Université Lédéa Bernard Ouédraogo
		}
		
		\vspace{1cm}
	\end{center}
	%============================================================================================================
	
	% --- TABLE DES MATIÈRES ---
	\tableofcontents
	
	% --- CORPS DU LIVRE ---
	\mainmatter
	
	\chapter{Rappels et compléments de calculs différentiels}
	\label{chap:calcul_diff}
	
	\section{Différentielle et gradient}
	
	L'optimisation repose sur l'étude des variations des fonctions. Le calcul différentiel est l'outil fondamental pour quantifier ces variations. Dans cette section, nous formalisons la notion de "meilleure approximation linéaire locale" d'une fonction, ce qui nous conduit aux concepts de différentielle et de gradient.
	
	\begin{definition}[Différentiabilité]
		Soit \(U \subset \mathbb{R}^n\) un ouvert et \(f : U \to \mathbb{R}^m\) une fonction. On dit que \(f\) est \textbf{différentiable} en un point \(x_0 \in U\) s'il existe une application linéaire \(L : \mathbb{R}^n \to \mathbb{R}^m\) telle que :
		\[
		f(x_0 + h) = f(x_0) + L(h) + o(\lVert h \rVert) \quad \text{lorsque } h \to 0,
		\]
		où \(o(\lVert h \rVert)\) désigne un terme \(\varepsilon(h) \in \mathbb{R}^m\) tel que \(\lVert \varepsilon(h) \rVert = o(\lVert h \rVert)\). Cela équivaut à dire que :
		\[
		\lim_{h \to 0} \frac{\lVert f(x_0 + h) - f(x_0) - L(h) \rVert}{\lVert h \rVert} = 0.
		\]
		L'application linéaire \(L\) est unique et est appelée la \textbf{différentielle} de \(f\) en \(x_0\), notée \(df(x_0)\).
	\end{definition}
	Si \(f\) est différentiable en \(x_0\), alors \(f\) admet des dérivées partielles par rapport à chacune de ses variables pour chacune de ses composantes en \(x_0\). Si \(f = (f_1, \dots, f_m)^T\), la \(i\)-ème dérivée partielle de la composante \(f_k\) est définie par :
	\[
	\frac{\partial f_k}{\partial x_i}(x_0) = \lim_{t \to 0} \frac{f_k(x_0 + t e_i) - f_k(x_0)}{t},
	\]
	où \((e_1, \dots, e_n)\) est la base canonique de \(\mathbb{R}^n\).
	
	L'application linéaire \(df(x_0)\) est alors représentée dans la base canonique par la \textbf{matrice jacobienne} :
	\[
	J_f(x_0) = \begin{pmatrix}
	\frac{\partial f_1}{\partial x_1}(x_0) & \cdots & \frac{\partial f_1}{\partial x_n}(x_0) \\
	\vdots & \ddots & \vdots \\
	\frac{\partial f_m}{\partial x_1}(x_0) & \cdots & \frac{\partial f_m}{\partial x_n}(x_0)
	\end{pmatrix}.
	\]
	L'action de la différentielle sur un vecteur \(h \in \mathbb{R}^n\) est donnée par le produit matriciel :
	\[
	df(x_0)(h) = J_f(x_0) h.
	\]
	
	
	\begin{exemple}
		Considérons la fonction \( f : \mathbb{R}^2 \to \mathbb{R}^2 \) définie par :
		\[
		f(x, y) = \begin{pmatrix}
		x^2 + y \\
		xy
		\end{pmatrix}.
		\]
		Vérifions que \( f \) est différentiable au point \( x_0 = (1, 1) \) en montrant que :
		\[
		\lim_{h \to 0} \frac{\lVert f(x_0 + h) - f(x_0) - L(h) \rVert}{\lVert h \rVert} = 0,
		\]
		où \( L : \mathbb{R}^2 \to \mathbb{R}^2 \) est une application linéaire.
		
		\textbf{Étape 1 : Calcul de la matrice jacobienne}\\
		La fonction \( f \) a deux composantes : \( f_1(x, y) = x^2 + y \) et \( f_2(x, y) = xy \). La matrice jacobienne est donnée par :
		\[
		J_f(x, y) = \begin{pmatrix}
		\frac{\partial f_1}{\partial x} & \frac{\partial f_1}{\partial y} \\
		\frac{\partial f_2}{\partial x} & \frac{\partial f_2}{\partial y}
		\end{pmatrix} = \begin{pmatrix}
		2x & 1 \\
		y & x
		\end{pmatrix}.
		\]
		Au point \( x_0 = (1, 1) \), on a :
		\[
		J_f(1, 1) = \begin{pmatrix}
		2 & 1 \\
		1 & 1
		\end{pmatrix}.
		\]
		L'application linéaire \( L = df(x_0) \) est donc :
		\[
		L(h_1, h_2) = \begin{pmatrix}
		2 & 1 \\
		1 & 1
		\end{pmatrix} \begin{pmatrix}
		h_1 \\
		h_2
		\end{pmatrix} = \begin{pmatrix}
		2h_1 + h_2 \\
		h_1 + h_2
		\end{pmatrix}.
		\]
		
		\textbf{Étape 2 : Calcul de l'erreur}\\
		Évaluons \( f(x_0) = f(1, 1) \) :
		\[
		f(1, 1) = \begin{pmatrix}
		1^2 + 1 \\
		1 \cdot 1
		\end{pmatrix} = \begin{pmatrix}
		2 \\
		1
		\end{pmatrix}.
		\]
		Pour \( h = (h_1, h_2) \), calculons \( f(x_0 + h) = f(1 + h_1, 1 + h_2) \) :
		\[
		f(1 + h_1, 1 + h_2) = \begin{pmatrix}
		(1 + h_1)^2 + (1 + h_2) \\
		(1 + h_1)(1 + h_2)
		\end{pmatrix} = \begin{pmatrix}
		2 + 2h_1 + h_2 + h_1^2 \\
		1 + h_1 + h_2 + h_1 h_2
		\end{pmatrix}.
		\]
		Ainsi :
		\[
		f(x_0 + h) - f(x_0) = \begin{pmatrix}
		2h_1 + h_2 + h_1^2 \\
		h_1 + h_2 + h_1 h_2
		\end{pmatrix}.
		\]
		Soustrayons \( L(h) \) :
		\[
		f(x_0 + h) - f(x_0) - L(h) = \begin{pmatrix}
		2h_1 + h_2 + h_1^2 \\
		h_1 + h_2 + h_1 h_2
		\end{pmatrix} - \begin{pmatrix}
		2h_1 + h_2 \\
		h_1 + h_2
		\end{pmatrix} = \begin{pmatrix}
		h_1^2 \\
		h_1 h_2
		\end{pmatrix}.
		\]
		
		\textbf{Étape 3 : Vérification de la limite}\\
		Calculons la norme de l'erreur :
		\[
		\lVert f(x_0 + h) - f(x_0) - L(h) \rVert = \sqrt{(h_1^2)^2 + (h_1 h_2)^2} = \sqrt{h_1^4 + h_1^2 h_2^2} = |h_1| \sqrt{h_1^2 + h_2^2}.
		\]
		La norme de \( h \) est :
		\[
		\lVert h \rVert = \sqrt{h_1^2 + h_2^2}.
		\]
		Ainsi :
		\[
		\frac{\lVert f(x_0 + h) - f(x_0) - L(h) \rVert}{\lVert h \rVert} = \frac{|h_1| \sqrt{h_1^2 + h_2^2}}{\sqrt{h_1^2 + h_2^2}} = |h_1|.
		\]
		Lorsque \( h \to 0 \), \( |h_1| \to 0 \), donc :
		\[
		\lim_{h \to 0} \frac{\lVert f(x_0 + h) - f(x_0) - L(h) \rVert}{\lVert h \rVert} = 0.
		\]
		La fonction \( f \) est donc différentiable en \( (1, 1) \), avec la différentielle \( df(1, 1) \) donnée par la matrice jacobienne \( J_f(1, 1) \).
	\end{exemple}
	
	\begin{definition}[Gradient]
		Si $f: U \to \R$ est différentiable en $x \in U$, on appelle \textbf{gradient} de $f$ en $x$ le vecteur unique de $\R^n$, noté $\grad f(x)$, qui vérifie pour tout $h \in \R^n$ :
		$$ df(x)(h) = \langle \grad f(x), h \rangle $$
		où $\langle \cdot, \cdot \rangle$ est le produit scalaire euclidien.
		Les composantes du gradient dans la base canonique sont les dérivées partielles :
		$$ \grad f(x) = \begin{pmatrix} \frac{\partial f}{\partial x_1}(x) \\ \vdots \\ \frac{\partial f}{\partial x_n}(x) \end{pmatrix} $$
	\end{definition}
	
	\begin{remark}[Interprétation géométrique du gradient]
		Le gradient a une importance capitale en optimisation.
		\begin{enumerate}
			\item \textbf{Direction de plus grande pente} : Le gradient $\grad f(x)$ pointe dans la direction où la fonction $f$ augmente le plus rapidement au point $x$. La norme $\norm{\grad f(x)}$ donne la valeur de cette pente maximale.
			\item \textbf{Orthogonalité aux lignes de niveau} : Le gradient $\grad f(x)$ est orthogonal à la ligne de niveau (ou surface de niveau) de $f$ passant par $x$, c'est-à-dire l'ensemble $\{y \in U \mid f(y) = f(x)\}$.
		\end{enumerate}
	\end{remark}
	
	\begin{example}
		Soit la fonction $f(x, y) = x^2 + 2y^2$ définie sur $\R^2$.
		Les dérivées partielles sont $\frac{\partial f}{\partial x} = 2x$ et $\frac{\partial f}{\partial y} = 4y$.
		Le gradient est $\grad f(x, y) = \begin{pmatrix} 2x \\ 4y \end{pmatrix}$.
		Au point $(2, 1)$, $f(2, 1) = 6$. Le gradient est $\grad f(2, 1) = \begin{pmatrix} 4 \\ 4 \end{pmatrix}$.
		Ce vecteur indique la direction $(4,4)$ comme étant celle de la plus forte croissance de $f$ à partir du point $(2,1)$. Ce vecteur est bien orthogonal à la tangente de l'ellipse $x^2 + 2y^2 = 6$ au point $(2,1)$.
	\end{example}
	
	
	\section{Dérivation des fonctions composées}
	
	La règle de dérivation en chaîne est essentielle pour calculer les dérivées de fonctions complexes, ce qui est omniprésent en optimisation, notamment dans les algorithmes de descente de gradient appliqués aux réseaux de neurones.
	
	\begin{theorem}[Règle de la chaîne - The Chain Rule]
		Soient $U \subset \R^n$ et $V \subset \R^m$ des ouverts. Soit $g: U \to V$ une fonction différentiable en $x_0 \in U$ et $f: V \to \R^p$ une fonction différentiable en $y_0 = g(x_0)$.
		Alors la fonction composée $h = f \circ g : U \to \R^p$ est différentiable en $x_0$ et sa différentielle est la composition des différentielles :
		$$ dh(x_0) = df(g(x_0)) \circ dg(x_0) $$
		En termes de matrices jacobiennes, cela se traduit par le produit matriciel :
		$$ J_h(x_0) = J_f(g(x_0)) \cdot J_g(x_0) $$
	\end{theorem}
	
	\begin{proof}
		Par définition de la différentiabilité de $g$ en $x_0$ et de $f$ en $y_0 = g(x_0)$:
		\begin{align*}
		g(x_0+k) - g(x_0) &= J_g(x_0)k + \epsilon_g(k) \norm{k} \\
		f(y_0+h) - f(y_0) &= J_f(y_0)h + \epsilon_f(h) \norm{h}
		\end{align*}
		où $\epsilon_g(k) \to 0$ quand $k \to 0$ et $\epsilon_f(h) \to 0$ quand $h \to 0$.
		On veut analyser $h(x_0+k) - h(x_0) = f(g(x_0+k)) - f(g(x_0))$.
		Posons $h(k) = g(x_0+k) - g(x_0)$. On a alors :
		\begin{align*}
		f(g(x_0+k)) - f(g(x_0)) &= f(y_0 + h(k)) - f(y_0) \\
		&= J_f(y_0)h(k) + \epsilon_f(h(k)) \norm{h(k)} \\
		&= J_f(y_0)(J_g(x_0)k + \epsilon_g(k)\norm{k}) + \epsilon_f(h(k))\norm{h(k)} \\
		&= J_f(y_0)J_g(x_0)k + J_f(y_0)\epsilon_g(k)\norm{k} + \epsilon_f(h(k))\norm{h(k)}
		\end{align*}
		Le terme principal est $J_f(g(x_0))J_g(x_0)k$. Il reste à montrer que les autres termes sont un $o(\norm{k})$. C'est le cas car $\norm{h(k)}$ est de l'ordre de $\norm{k}$ (car $g$ est différentiable).
	\end{proof}
	
	\begin{corollary}[Cas particulier pour le gradient]
		Si $f: \R^m \to \R$ (donc $p=1$) et $g: \R^n \to \R^m$. Alors $h = f \circ g : \R^n \to \R$.
		La jacobienne de $f$ est $(\grad f)\T$. La jacobienne de $h$ est $(\grad h)\T$.
		La règle de la chaîne donne :
		$$ (\grad h(x))\T = (\grad f(g(x)))\T \cdot J_g(x) $$
		En transposant, on obtient la formule pour le gradient :
		$$ \grad h(x) = (J_g(x))\T \grad f(g(x)) $$
	\end{corollary}
	
	\begin{example}
		Soit $f(y_1, y_2) = y_1^2 + y_2^3$ et $g(x) = \begin{pmatrix} \sin(x) \\ \cos(x) \end{pmatrix}$. Calculons la dérivée de $h(x) = f(g(x)) = \sin^2(x) + \cos^3(x)$.
		On a $\grad f(y_1, y_2) = \begin{pmatrix} 2y_1 \\ 3y_2^2 \end{pmatrix}$ et $J_g(x) = \begin{pmatrix} \cos(x) \\ -\sin(x) \end{pmatrix}$.
		Le gradient de $h$ (ici, juste une dérivée car $n=1$) est :
		\begin{align*}
		h'(x) = \grad h(x) &= (J_g(x))\T \grad f(g(x)) \\
		&= \begin{pmatrix} \cos(x) & -\sin(x) \end{pmatrix} \begin{pmatrix} 2\sin(x) \\ 3\cos^2(x) \end{pmatrix} \\
		&= 2\sin(x)\cos(x) - 3\sin(x)\cos^2(x)
		\end{align*}
		Ce qui correspond bien à la dérivée directe.
	\end{example}
	
	
	\section{Différentielle d'ordre deux et matrice hessienne}
	
	Pour comprendre la nature d'un point critique (minimum, maximum, ou selle), l'approximation linéaire du gradient ne suffit pas. Il faut examiner la courbure locale de la fonction, qui est capturée par les dérivées secondes.
	
	\begin{definition}[Dérivées partielles secondes]
		Si les dérivées partielles $\frac{\partial f}{\partial x_i}$ d'une fonction $f: U \to \R$ sont elles-mêmes dérivables, on peut définir les dérivées partielles secondes :
		$$ \frac{\partial^2 f}{\partial x_j \partial x_i}(x) = \frac{\partial}{\partial x_j}\left(\frac{\partial f}{\partial x_i}\right)(x) $$
	\end{definition}
	
	\begin{theorem}[Théorème de Schwarz]
		Si les dérivées partielles secondes $\frac{\partial^2 f}{\partial x_j \partial x_i}$ et $\frac{\partial^2 f}{\partial x_i \partial x_j}$ existent et sont continues dans un voisinage d'un point $x_0$, alors elles sont égales en ce point :
		$$ \frac{\partial^2 f}{\partial x_j \partial x_i}(x_0) = \frac{\partial^2 f}{\partial x_i \partial x_j}(x_0) $$
	\end{theorem}
	Ce théorème est fondamental car il garantit la symétrie de la matrice des dérivées secondes.
	
	\begin{definition}[Matrice Hessienne]
		Soit $f: U \to \R$ une fonction de classe $\mathcal{C}^2$ (deux fois continûment différentiable). La \textbf{matrice hessienne} de $f$ en un point $x \in U$, notée $\hess f(x)$ ou $H_f(x)$, est la matrice carrée $n \times n$ des dérivées partielles secondes :
		$$ H_f(x) = \begin{pmatrix}
		\frac{\partial^2 f}{\partial x_1^2}(x) & \frac{\partial^2 f}{\partial x_1 \partial x_2}(x) & \cdots & \frac{\partial^2 f}{\partial x_1 \partial x_n}(x) \\
		\frac{\partial^2 f}{\partial x_2 \partial x_1}(x) & \frac{\partial^2 f}{\partial x_2^2}(x) & \cdots & \frac{\partial^2 f}{\partial x_2 \partial x_n}(x) \\
		\vdots & \vdots & \ddots & \vdots \\
		\frac{\partial^2 f}{\partial x_n \partial x_1}(x) & \frac{\partial^2 f}{\partial x_n \partial x_2}(x) & \cdots & \frac{\partial^2 f}{\partial x_n^2}(x)
		\end{pmatrix} $$
		En vertu du théorème de Schwarz, la matrice hessienne d'une fonction de classe $\mathcal{C}^2$ est toujours symétrique.
	\end{definition}
	La hessienne représente la partie quadratique de l'approximation locale de la fonction.
	
	\section{Formules de Taylor}
	
	Les formules de Taylor sont des généralisations du théorème des accroissements finis. Elles permettent d'approximer une fonction au voisinage d'un point par un polynôme dont les coefficients dépendent des dérivées de la fonction en ce point.
	
	\begin{theorem}[Formule de Taylor-Young à l'ordre 2]
		Soit $f: U \to \R$ une fonction de classe $\mathcal{C}^2$ sur un ouvert $U \subset \R^n$. Pour tout $x \in U$, on a, pour $h \to 0$ :
		$$ f(x+h) = f(x) + \langle \grad f(x), h \rangle + \frac{1}{2} \langle h, \hess f(x) h \rangle + o(\norm{h}^2) $$
		$$ f(x+h) = f(x) + \grad f(x)\T h + \frac{1}{2} h\T \hess f(x) h + o(\norm{h}^2) $$
	\end{theorem}
	
	\begin{remark}[Interprétation]
		Cette formule est cruciale pour l'analyse des optima locaux. Au voisinage d'un point $x$:
		\begin{itemize}
			\item $f(x)$ est la valeur de base.
			\item $\grad f(x)\T h$ est l'approximation linéaire (le plan tangent).
			\item $\frac{1}{2} h\T \hess f(x) h$ est l'approximation quadratique qui décrit la courbure.
		\end{itemize}
		Si $x^*$ est un point critique, $\grad f(x^*) = 0$. Le comportement local de $f$ autour de $x^*$ est alors entièrement déterminé par le signe de la forme quadratique associée à la hessienne $\hess f(x^*)$.
	\end{remark}
	
	\begin{theorem}[Formule de Taylor avec reste intégral]
		Soit $f$ de classe $\mathcal{C}^{k+1}$ sur $U$. Alors :
		$$ f(x+h) = \sum_{j=0}^k \frac{1}{j!} d^j f(x)(h, \dots, h) + R_k(h) $$
		où le reste peut s'écrire sous forme intégrale. Pour $k=1$:
		$$ f(x+h) = f(x) + \grad f(x)\T h + \int_0^1 (1-t) h\T \hess f(x+th) h \, dt $$
	\end{theorem}
	
	\begin{theorem}[Formule de Taylor-Lagrange]
		Si $f$ est de classe $\mathcal{C}^2$, il existe $c \in ]0, 1[$ tel que :
		$$ f(x+h) = f(x) + \grad f(x)\T h + \frac{1}{2} h\T \hess f(x+ch) h $$
		Cette formule est une généralisation du théorème des accroissements finis et est très utile pour les preuves théoriques.
	\end{theorem}
	
	\section{Rappels d'analyse Convexe}
	La convexité est une propriété structurante en optimisation. Lorsqu'un problème est convexe, on peut affirmer que tout minimum local est également un minimum global, ce qui simplifie considérablement la recherche de solutions.
	
	\begin{definition}[Ensemble convexe]
		Un sous-ensemble $C \subseteq \R^n$ est dit \textbf{convexe} si pour tout couple de points $(x, y) \in C \times C$ et pour tout $\lambda \in [0, 1]$, le segment qui les relie est entièrement contenu dans $C$.
		$$ \forall x,y \in C, \forall \lambda \in [0,1], \quad \lambda x + (1-\lambda)y \in C $$
	\end{definition}
	
	\begin{definition}[Fonction convexe]
		Soit $C \subseteq \R^n$ un ensemble convexe. Une fonction $f: C \to \R$ est dite \textbf{convexe} si son épigraphe, l'ensemble $\{(x,t) \in C \times \R \mid t \ge f(x)\}$, est un ensemble convexe.
		Une caractérisation équivalente (et plus commune) est l'inégalité de Jensen :
		$$ \forall x,y \in C, \forall \lambda \in [0,1], \quad f(\lambda x + (1-\lambda)y) \le \lambda f(x) + (1-\lambda)f(y) $$
		La fonction est dite \textbf{strictement convexe} si l'inégalité est stricte pour $x \ne y$ et $\lambda \in ]0, 1[$.
	\end{definition}
	
	\begin{proposition}[Caractérisation au premier ordre]
		Soit $f: C \to \R$ une fonction différentiable sur un ensemble convexe ouvert $C$.
		$f$ est convexe si et seulement si pour tous $x, y \in C$:
		$$ f(y) \ge f(x) + \grad f(x)\T (y-x) $$
		Géométriquement, le graphe de la fonction est toujours au-dessus de sa tangente en n'importe quel point.
	\end{proposition}
	
	\begin{proof}
		($\Rightarrow$) Supposons $f$ convexe. Par définition, pour $\lambda \in ]0,1]$:
		$$ f(x + \lambda(y-x)) \le (1-\lambda)f(x) + \lambda f(y) $$
		En réarrangeant, on a :
		$$ \frac{f(x + \lambda(y-x)) - f(x)}{\lambda} \le f(y) - f(x) $$
		En faisant tendre $\lambda \to 0^+$, le terme de gauche devient la dérivée directionnelle de $f$ en $x$ dans la direction $y-x$, qui est $\grad f(x)\T(y-x)$. D'où le résultat.
		($\Leftarrow$) La preuve de la réciproque est un peu plus technique et utilise le théorème des accroissements finis.
	\end{proof}
	
	\begin{proposition}[Caractérisation au second ordre]
		Soit $f: C \to \R$ une fonction de classe $\mathcal{C}^2$ sur un ensemble convexe ouvert $C$.
		\begin{enumerate}
			\item $f$ est convexe si et seulement si sa matrice hessienne $\hess f(x)$ est semi-définie positive pour tout $x \in C$. ($\forall h \in \R^n, h\T \hess f(x) h \ge 0$).
			\item Si $\hess f(x)$ est définie positive pour tout $x \in C$, alors $f$ est strictement convexe.
		\end{enumerate}
	\end{proposition}
	\begin{proof}
		Utilisons la formule de Taylor-Lagrange entre $x$ et $y=x+h$:
		$$ f(y) = f(x) + \grad f(x)\T(y-x) + \frac{1}{2}(y-x)\T \hess f(x+c(y-x)) (y-x) $$
		pour un certain $c \in ]0,1[$.
		Si la hessienne est semi-définie positive partout, alors le dernier terme est $\ge 0$. Donc $f(y) \ge f(x) + \grad f(x)\T(y-x)$, ce qui implique la convexité par la caractérisation du premier ordre. La réciproque se fait par l'absurde.
	\end{proof}
	
	\begin{theorem}[Point critique d'une fonction convexe]
		Soit $f: C \to \R$ une fonction convexe et différentiable. Tout point critique $x^*$ ($\grad f(x^*) = 0$) est un \textbf{minimum global} de $f$ sur $C$. S'il $f$ est strictement convexe, ce minimum est unique.
	\end{theorem}
	\begin{proof}
		Par la caractérisation du premier ordre, pour tout $y \in C$ :
		$$ f(y) \ge f(x^*) + \grad f(x^*)\T(y-x^*) $$
		Comme $\grad f(x^*) = 0$, on a $f(y) \ge f(x^*)$ pour tout $y \in C$. $x^*$ est donc un minimum global. L'unicité pour le cas strictement convexe découle de la définition.
	\end{proof}
	
	% ... (Suite du document LaTeX précédent) ...
	
	\section{Etude des fonctionnelles quadratiques}
	
	Les fonctionnelles quadratiques sont fondamentales en optimisation car elles constituent l'approximation d'ordre 2 de n'importe quelle fonction régulière (via la formule de Taylor). De plus, de nombreux problèmes se modélisent directement sous cette forme.
	
	\begin{definition}[Fonctionnelle quadratique]
		Une fonction $f: \R^n \to \R$ est une \textbf{fonctionnelle quadratique} si elle peut s'écrire sous la forme :
		$$ f(x) = \frac{1}{2} x\T A x - b\T x + c $$
		où $A$ est une matrice carrée de taille $n \times n$, $b \in \R^n$ et $c \in \R$.
		Sans perte de généralité, on peut toujours supposer que la matrice $A$ est \textbf{symétrique}. En effet, si elle ne l'est pas, on peut la remplacer par sa partie symétrique $A_{sym} = \frac{1}{2}(A + A\T)$ sans changer la valeur de la fonction, car $x\T A x = x\T A_{sym} x$.
	\end{definition}
	
	Calculons le gradient et la hessienne de cette fonction.
	Pour le gradient, en développant $f(x) = \frac{1}{2} \sum_{i,j} A_{ij}x_i x_j - \sum_k b_k x_k + c$, on dérive par rapport à $x_k$:
	$$ \frac{\partial f}{\partial x_k} = \frac{1}{2} \sum_j A_{kj}x_j + \frac{1}{2} \sum_i A_{ik}x_i - b_k $$
	Comme $A$ est symétrique ($A_{ik} = A_{ki}$), cela devient $\sum_j A_{kj}x_j - b_k$, qui est la $k$-ième composante du vecteur $Ax - b$.
	Pour la hessienne, on dérive une seconde fois : $\frac{\partial^2 f}{\partial x_l \partial x_k} = A_{kl}$.
	
	\begin{proposition}[Gradient et Hessienne d'une forme quadratique]
		Soit $f(x) = \frac{1}{2} x\T A x - b\T x + c$ avec $A$ symétrique.
		Le gradient de $f$ est :
		$$ \grad f(x) = Ax - b $$
		La matrice hessienne de $f$ est constante et égale à $A$ :
		$$ \hess f(x) = A $$
	\end{proposition}
	
	La convexité de $f$ est donc directement liée aux propriétés de la matrice $A$.
	\begin{corollary}[Convexité d'une forme quadratique]
		La fonctionnelle quadratique $f(x) = \frac{1}{2} x\T A x - b\T x + c$ est :
		\begin{enumerate}
			\item \textbf{convexe} sur $\R^n$ si et seulement si $A$ est semi-définie positive.
			\item \textbf{strictement convexe} sur $\R^n$ si et seulement si $A$ est définie positive.
		\end{enumerate}
	\end{corollary}
	
	Un point critique (ou stationnaire) $x^*$ est un point qui annule le gradient : $\grad f(x^*) = Ax^* - b = 0$.
	\begin{itemize}
		\item Si $A$ est inversible, il existe un unique point critique $x^* = A^{-1}b$.
		\item Si de plus $A$ est définie positive, alors $f$ est strictement convexe. D'après le théorème précédent, cet unique point critique est l'unique \textbf{minimum global} de $f$.
		\item Si $A$ n'est pas semi-définie positive, le point critique (s'il existe) ne sera pas un minimum.
	\end{itemize}
	
	\section{Mineurs principaux et critère de Sylvester}
	
	Pour appliquer les résultats précédents, nous avons besoin d'un outil pratique pour déterminer si une matrice symétrique est définie positive.
	
	\begin{definition}[Mineurs principaux dominants]
		Soit $A$ une matrice carrée de taille $n \times n$. Pour $k \in \{1, \dots, n\}$, le \textbf{$k$-ième mineur principal dominant} de $A$, noté $\Delta_k$, est le déterminant de la sous-matrice carrée de taille $k \times k$ extraite du coin supérieur gauche de $A$.
		$$ \Delta_1 = A_{11}, \quad \Delta_2 = \det \begin{pmatrix} A_{11} & A_{12} \\ A_{21} & A_{22} \end{pmatrix}, \quad \dots, \quad \Delta_n = \det(A) $$
	\end{definition}
	
	\begin{theorem}[Critère de Sylvester]
		Soit $A$ une matrice symétrique.
		\begin{enumerate}
			\item $A$ est \textbf{définie positive} si et seulement si tous ses mineurs principaux dominants sont strictement positifs.
			$$ A \succ 0 \iff \forall k \in \{1, \dots, n\}, \quad \Delta_k > 0 $$
			\item $A$ est \textbf{définie négative} si et seulement si ses mineurs principaux dominants alternent en signe, en commençant par un signe négatif.
			$$ A \prec 0 \iff \forall k \in \{1, \dots, n\}, \quad (-1)^k \Delta_k > 0 $$
		\end{enumerate}
	\end{theorem}
	
	\begin{remark}
		Attention, si certains $\Delta_k$ sont nuls et d'autres positifs, on ne peut pas conclure directement sur le caractère semi-défini positif. Il faut alors examiner tous les mineurs principaux (pas seulement les dominants), ou bien calculer les valeurs propres de la matrice. Une matrice symétrique est semi-définie positive si et seulement si toutes ses valeurs propres sont $\ge 0$.
	\end{remark}
	
	\begin{example}
		Considérons la fonction $f(x,y,z) = x^2 + 2y^2 + 3z^2 + 2xy + 2yz$.
		La forme quadratique associée est $f(x,y,z) = \frac{1}{2} \begin{pmatrix} x & y & z \end{pmatrix} A \begin{pmatrix} x \\ y \\ z \end{pmatrix}$ avec :
		$$ A = \begin{pmatrix} 2 & 2 & 0 \\ 2 & 4 & 2 \\ 0 & 2 & 6 \end{pmatrix} $$
		Calculons les mineurs principaux dominants :
		\begin{itemize}
			\item $\Delta_1 = 2 > 0$
			\item $\Delta_2 = \det \begin{pmatrix} 2 & 2 \\ 2 & 4 \end{pmatrix} = 8 - 4 = 4 > 0$
			\item $\Delta_3 = \det(A) = 2(24-4) - 2(12-0) = 40 - 24 = 16 > 0$
		\end{itemize}
		Tous les mineurs principaux dominants sont strictement positifs. Par le critère de Sylvester, la matrice $A$ est définie positive. La fonction $f$ est donc strictement convexe et admet un unique minimum global (qui est en $(0,0,0)$ car il n'y a pas de terme linéaire).
	\end{example}
	
	
	\section{Résultats d'existence et d'unicité}
	
	Un problème d'optimisation peut-il ne pas avoir de solution ? Oui, par exemple si la fonction tend vers $-\infty$ ou si le minimum est "à la limite" d'un ensemble ouvert. Les théorèmes suivants donnent des conditions suffisantes pour garantir l'existence d'une solution.
	
	\subsection{Théorème de Weierstrass}
	Ce théorème est le pilier de l'existence d'optima sur des domaines contraints.
	
	\begin{definition}[Ensemble compact]
		Un sous-ensemble $K \subset \R^n$ est dit \textbf{compact} s'il est fermé et borné.
		\begin{itemize}
			\item \textbf{Fermé} signifie que $K$ contient la limite de toute suite convergente d'éléments de $K$. Typiquement, les ensembles définis par des inégalités larges ($\le, \ge$) ou des égalités ($=$) de fonctions continues sont fermés.
			\item \textbf{Borné} signifie qu'il existe un réel $M > 0$ tel que pour tout $x \in K$, $\norm{x} \le M$.
		\end{itemize}
	\end{definition}
	
	\begin{theorem}[Théorème de Weierstrass ou des bornes atteintes]
		Soit $K \subset \R^n$ un ensemble compact non vide, et soit $f: K \to \R$ une fonction continue sur $K$.
		Alors $f$ est bornée sur $K$ et atteint ses bornes. C'est-à-dire qu'il existe au moins un point $x_{min} \in K$ (minimum global) et au moins un point $x_{max} \in K$ (maximum global) tels que :
		$$ \forall x \in K, \quad f(x_{min}) \le f(x) \le f(x_{max}) $$
	\end{theorem}
	Ce théorème garantit l'existence d'une solution pour un très grand nombre de problèmes d'optimisation sous contraintes, à condition que l'ensemble des solutions admissibles soit compact.
	
	\subsection{Existence de minimum sur un ensemble non-compact}
	
	Que se passe-t-il si le domaine n'est pas borné, comme $\R^n$ entier ? L'existence n'est plus garantie (e.g., $f(x)=x$ n'a pas de minimum sur $\R$). Il faut une condition sur le comportement de la fonction "à l'infini".
	
	\begin{definition}[Fonction coercive]
		Une fonction $f: \R^n \to \R$ est dite \textbf{coercive} si elle tend vers $+\infty$ lorsque la norme de son argument tend vers l'infini.
		$$ \lim_{\norm{x} \to +\infty} f(x) = +\infty $$
	\end{definition}
	Intuitivement, une fonction coercive "remonte" sur les bords de son domaine, créant une sorte de "cuvette" qui "emprisonne" le minimum.
	
	\begin{theorem}[Existence d'un minimum global]
		Soit $f: \R^n \to \R$ une fonction continue et coercive. Alors $f$ admet au moins un minimum global sur $\R^n$.
	\end{theorem}
	
	\begin{proof}[Idée de la preuve]
		Soit $x_0 \in \R^n$ un point quelconque. Puisque $f$ est coercive, il existe $R > 0$ tel que pour tout $x$ avec $\norm{x} > R$, on a $f(x) > f(x_0)$.
		Considérons la restriction de $f$ à la boule fermée $B = \{x \in \R^n \mid \norm{x} \le R\}$.
		$B$ est un ensemble compact. D'après le théorème de Weierstrass, $f$ admet un minimum sur $B$, disons en un point $x^* \in B$.
		On a donc $f(x^*) \le f(x)$ pour tout $x \in B$. En particulier, comme $x_0 \in B$, $f(x^*) \le f(x_0)$.
		Maintenant, si l'on prend un point $y \notin B$, on a $\norm{y} > R$, donc $f(y) > f(x_0) \ge f(x^*)$.
		Ainsi, $f(x^*) \le f(z)$ pour tout $z \in \R^n$, que $z$ soit dans $B$ ou non. $x^*$ est donc un minimum global.
	\end{proof}
	
	En combinant les notions de convexité et de coercivité, on obtient un résultat très puissant d'existence et d'unicité, qui est le Graal de nombreux problèmes d'optimisation.
	
	\begin{theorem}[Existence et Unicité]
		Soit $f: \R^n \to \R$ une fonction continue.
		\begin{enumerate}
			\item Si $f$ est convexe et coercive, elle admet au moins un minimum global.
			\item Si $f$ est strictement convexe et coercive, elle admet un \textbf{unique} minimum global.
		\end{enumerate}
	\end{theorem}
	
	% ... (Fin du Chapitre 1) ...
	
	\chapter{Optimisation sans contraintes}
	\section{Définitions et notations :}
	\begin{definition}[Définition d'un problème d'optimisation] Soit $D \subset \Re^{n}$ et $f: \Re^{n} \longrightarrow \Re$. Un problème d'optimisation consiste à chercher un élément $x^{*} \in D$ ( s'il existe ) pour lequel $f\left(x^{*}\right)$ est la plus petite ( resp la plus grande ) valeur de $f$ sur $D$ et on écrit : $\displaystyle\min _{x \in D} f(x)=f\left(x^{*}\right),\left(\right.$ resp $\left.\max _{x \in D} f(x)=f\left(x^{*}\right)\right)$.\\
	\end{definition}
	
	$D$ est appelé ensemble des contraintes ( ou des solutions admissibles, ou réalisables ), dit aussi domaine de faisabilité.\\
	La fonction $f$ est appelée fonction coût, objectif ou économique.\\
	Considérons le problème d'optimisation suivant :
	
	$$
	(P):\left\{\begin{array}{c}
	\displaystyle\min f(x) \\
	x \in D
	\end{array}\right.
	$$
	
	\begin{remark}
		\begin{enumerate}
			\item Si $D=\Re^{n}$, on dit que ( $P$ ) est un problème sans contraintes.
			\item Un problème de maximisation se ramène à un problème de minimisation par la relation : $\max f(x)=-\displaystyle\min (-f(x))$
			\item L'ensemble $D$ est donné souvent comme suit :
		\end{enumerate}
		
		$$
		D=\left\{x \in \Re^{n} / g_{i}(x) \leq 0, i=1, \ldots, m, \text { et }, h_{j}(x)=0, j=1, \ldots, p\right\}
		$$
		
		(P) est appelé dans ce cas un programme mathématique.
	\end{remark}
	
	
	\begin{notation}
		\begin{enumerate}
			\item Le problème ( $P$ ) est dit réalisable si : $D \neq \emptyset$, un point $x \in D$ est appelé solution réalisable de $(P)$.
			\item Une solution réalisable qui minimise ( resp maximise) $f$ sur $D$ est appelée solution optimale.
			\item La valeur de l'objectif au point $x^{*}$ est appelée valeur optimale, c'est la valeur minimale ou la borne inférieure (resp la valeur maximale ou la borne supérieure ) de $f$ sur $D$.
		\end{enumerate}
	\end{notation}
	
	\subsection{Classification des problèmes d'optimisation :}
	Les problèmes d'optimisation sont classifiés selon les caractéristiques des fonctions $f, g_{i}$, et $h_{j}$.
	
	\begin{enumerate}
		\item Si $f$ et $g_{i}$ sont convexes et $h_{j}$ sont affines, on dit qu'il s'agit d'un programme convexe ( dans ce cas l'ensemble $D$ est convexe ).
		\item Si les fonctions $f$, $g_{i}$ et $h_{j}$ sont toutes différentiables, on parle de programmation mathématique différentiable.
	\end{enumerate}
	
	\begin{definition}[ Minimum local/Maximum global]
		\begin{enumerate}
			\item Un point $x^{*} \in D$ est un minimum local ( resp maximum local ) de $f$ s'il existe une boule ouverte $B$ centrée en $x^{*}$ telle que :
		\end{enumerate}
		
		$$
		\forall x \in B \cap D, f\left(x^{*}\right) \leq f(x), \quad\left(\operatorname{resp} f\left(x^{*}\right) \geq f(x)\right) .
		$$
		
		\begin{enumerate}
			\setcounter{enumi}{1}
			\item Un point $x^{*} \in D$ est un minimum global (resp maximum global) de $f$ si :
		\end{enumerate}
		
		$$
		\forall x \in D, f\left(x^{*}\right) \leq f(x), \quad\left(\operatorname{resp} f\left(x^{*}\right) \geq f(x)\right) .
		$$
	\end{definition}
	
	
	\section{Optimisation sans contraintes :}
	Considérons le problème d'optimisation sans contraintes suivant :
	
	$$
	(P):\left\{\begin{array}{c}
	\displaystyle\min f(x) \\
	x \in \Re^{n}
	\end{array}\right.
	$$
	
	
	Avant d'étudier la solution (ou les solutions) d'un problème d'optimisation avec ou sans contrainte, il faut s'assurer de leur existence.
	\subsection{Résultats d'existence et d'unicité de la solution :}
	\begin{definition} Une fonction $f: \Re^{n} \longrightarrow \Re$ est dite coercive si et seulement si :
		
		$$
		\begin{gathered}
		\lim _{\|x\| \longrightarrow+\infty} f(x)=+\infty . \\
		\left(c-\grave{a}-d: \forall A \in \Re_{+}^{*} \exists N>0 / \forall x \in \Re^{n},\|x\| \geq N \Longrightarrow f(x) \geq A\right) .
		\end{gathered}
		$$
	\end{definition}
	Pour montrer que $f$ est coercive on utilise souvent la proposition suivante :\\
	\begin{proposition} Soit $f: \Re^{n} \longrightarrow \Re$ une fonction et $g: \Re \longrightarrow \Re$ vérifiant :
		
		$$
		f(x) \geq g(\|x\|) \text { avec } \lim _{t \longrightarrow+\infty} g(t)=+\infty
		$$
		
		Alors f est coercive.\\
	\end{proposition}
	
	\begin{theorem}[Existence] Soit $f: \Re^{n} \longrightarrow \Re$ une fonction propre, continue, et, coercive, alors il existe un point $x^{*} \in \Re^{n}$ qui réalise le minimum de $f$ sur $\Re^{n}$.
		
		Autrement dit $\exists x^{*} \in \Re^{n}$ tel que :
		
		$$
		f\left(x^{*}\right) \leq f(x) \quad, \forall x \in \Re^{n}
		$$
	\end{theorem}
	\begin{proof}
		Soit : $m=\inf _{x \in \Re^{n}}(f(x))<+\infty$, soit $\left(x_{n}\right)_{n \in N}$ une suite minimisante, $c-\grave{a}-d$
		
		
		\begin{equation*}
		\lim _{n \longrightarrow+\infty} f\left(x_{n}\right)=m<+\infty \tag{1}
		\end{equation*}
		
		
		Supposons que $\left(x_{n}\right)_{n \in N}$ n'est pas bornée, donc il existe une sous suite $\left(x_{\psi(n)}\right)_{n \in N}$ de $\left(x_{n}\right)_{n \in N}$ telle que : $\lim _{n \longrightarrow+\infty}\left\|x_{\psi(n)}\right\|=+\infty$.
		
		Par coercivité de $f$, on a $\lim _{n \longrightarrow+\infty} f\left(x_{\psi(n)}\right)=+\infty$, contradiction avec (1). Alors la suite $\left(x_{n}\right)_{n}$ est bornée, et dans ce cas il existe une suite extraite de $\left(x_{n}\right)_{n}$, soit $\left(x_{\varphi(n)}\right)_{n}$ qui converge vers $x^{*} \in \Re^{n}$. En utilisant la continuité de $f$ on a $: f\left(x^{*}\right)= \lim _{n \longrightarrow+\infty} f\left(x_{\varphi(n)}\right)=m$. On en déduit alors que $x^{*}$ solution de $(P)$ et $m>-\infty$.
	\end{proof}
	\begin{theorem}[Unicité] Soit $f: \Re^{n} \longrightarrow \Re$ une fonction strictement convexe, alors il existe au plus $x^{*} \in \Re^{n}$ tel que : $f\left(x^{*}\right)=\displaystyle\min _{x \in \Re^{n}} f(x)$.
	\end{theorem}
	\begin{proof}
		Soit $f$ strictement convexe, supposons qu'il existe $x_{1}^{*}, x_{2}^{*} \in \Re^{n}$ tels que $f\left(x_{1}^{*}\right)= f\left(x_{2}^{*}\right)=\displaystyle\min _{x \in \Re^{n}} f(x)$. Supposons que $x_{1}^{*} \neq x_{2}^{*}$, et comme $f$ est strictement convexe, on $a$ :
		
		$$
		f\left(\frac{1}{2} x_{1}^{*}+\frac{1}{2} x_{2}^{*}\right)<\frac{1}{2} f\left(x_{1}^{*}\right)+\frac{1}{2} f\left(x_{2}^{*}\right)=\displaystyle\min _{x \in \Re^{n}} f(x)
		$$
		
		ce qui est impossible, donc $x_{1}^{*}=x_{2}^{*}$.\\
	\end{proof}

\begin{theorem}[Existence et unicité] 
	Soit $f: \Re^{n} \rightarrow \Re$, vérifiant :
	
	\begin{enumerate}
		\item $f$ continue;
		\item $f$ coercive;
		\item $f$ est strictement convexe.
	\end{enumerate}
	
	Alors il existe un unique $x^{*} \in \Re^{n}$, tel que $f\left(x^{*}\right)=\displaystyle\min _{x \in \Re^{n}} f(x)$.\\
\end{theorem}

\begin{theorem}
	Soit $f: \Re^{n} \rightarrow \Re$, une fonction de classe $C^{1}$. Supposons qu'il existe $\alpha>0$ tel que :
	
	$$
	\forall x, y \in \Re^{n} \quad, \quad\langle\nabla f(x)-\nabla f(y), x-y\rangle \geq \alpha\|x-y\|^{2}
	$$
	
	Alors $f$ est strictement convexe et coercive, en particulier le problème ( $P$ ) admet une solution unique .
\end{theorem}
	
	\subsection{Conditions d'optimalité :}
	\subsubsection{Conditions nécessaires :}
	\begin{theorem} Soit $x^{*}$ un minimum local de $f$ sur $\Re^{n}$, tel que $f: \Re^{n} \rightarrow \Re$ différentiable, alors on $a$ :
		
		\begin{enumerate}
			\item $\nabla f\left(x^{*}\right)=0$. (Condition nécessaire d'optimalité du premier ordre).
			\item $\nabla^{2} f\left(x^{*}\right)$ est une matrice semi-définie positive. (Condition nécessaire d'optimalité du second ordre).
		\end{enumerate}
	\end{theorem}
	
	\begin{proof}
		\begin{enumerate}
			\item $\forall t \in \Re^{*}$ et $\forall d \in \Re^{n}$ on a :
		\end{enumerate}
		
		$$
		\lim _{t \rightarrow 0^{+}} \frac{f\left(x^{*}+t d\right)-f\left(x^{*}\right)}{t}=d^{T} \nabla f\left(x^{*}\right) \geq 0,
		$$
		
		et
		
		$$
		\lim _{t \rightarrow 0^{-}} \frac{f\left(x^{*}+t d\right)-f\left(x^{*}\right)}{t}=d^{T} \nabla f\left(x^{*}\right) \leq 0
		$$
		
		d'où : $d^{T} \nabla f\left(x^{*}\right)=0, \forall d \in \Re^{n}$, alors $\nabla f\left(x^{*}\right)=0$.\\
		2. Soit $d$ une direction arbitraire et $t>0$, le développement de Taylor du second ordre au voisinage de $x^{*}$ est donné par :
		
		$$
		\begin{aligned}
		f\left(x^{*}+t d\right) & =f\left(x^{*}\right)+t d^{T} \nabla f\left(x^{*}\right)+\frac{1}{2} t^{2} d^{T} \nabla^{2} f\left(x^{*}\right) d+o\left(\|t d\|^{2}\right) \\
		& \Longrightarrow f\left(x^{*}+t d\right)-f\left(x^{*}\right)=\frac{1}{2} t^{2} d^{T} \nabla^{2} f\left(x^{*}\right) d+o\left(\|t d\|^{2}\right) \geq 0
		\end{aligned}
		$$
		
		en divisant par $t^{2}$ on obtient :
		
		$$
		\frac{1}{2} d^{T} \nabla^{2} f\left(x^{*}\right) d+\frac{o\left(\|t d\|^{2}\right)}{t^{2}} \geq 0
		$$
		
		d'après la notation de Landeau :
		
		$$
		\lim _{t \rightarrow 0} \frac{o\left(\|t d\|^{2}\right)}{t^{2}}=0
		$$
		
		alors : $d^{T} \nabla^{2} f\left(x^{*}\right) d \geq 0$.\\
	\end{proof}
	
	\begin{remark} 
		\begin{enumerate}
			\item Un minimum local d'une fonction différentiable est un point stationnaire. La réciproque est fausse.
			\item  On dit qu’un point $x$ est un \textbf{point critique} de la fonctionnelle $f$ si 
			\[
			\nabla f(x) = 0.
			\]
			
			On s’intéresse maintenant aux conditions nécessaires et suffisantes faisant intervenir la dérivée seconde.
			
			\medskip
			
			On rappelle qu’une matrice symétrique $A$ est dite :
			\begin{itemize}
				\item \textbf{semi-définie positive} si 
				\[
				\forall u \in \mathbb{R}^n, \quad \langle Au, u \rangle \geq 0,
				\]
				\item \textbf{définie positive} si 
				\[
				\forall u \in \mathbb{R}^n, \; u \neq 0, \quad \langle Au, u \rangle > 0.
				\]
			\end{itemize}
			
			On rappelle également que $A$ est semi-définie positive (respectivement définie positive) si et seulement si
			toutes ses valeurs propres sont positives (respectivement strictement positives).
			
		\end{enumerate}
	\end{remark}
	
	\subsubsection{Conditions suffisantes :}
	Comme le montre le résultat suivant, la dérivée seconde permet souvent de déterminer la nature d’un point critique, c’est-à-dire de déterminer si un point critique est bien un minimum local, un maximum local, ou ni l’un ni l’autre.
	
	\begin{theorem} Soient $f \in C^{2}\left(\Re^{n}\right)$, et $x^{*} \in \Re^{n}$ tels que :
		
		\begin{enumerate}
			\item $\nabla f\left(x^{*}\right)=0$;
			\item $\nabla^{2} f\left(x^{*}\right)$ est définie positive i.e \(
			\forall\, d \in \mathbb{R}^n \setminus \{0\}, \quad 
			\langle \nabla^2 f(x)d,\, d \rangle > 0
			\)
			;\\
			alors $x^{*}$ est un minimum local de $f$.
		\end{enumerate}
	\end{theorem}
	\begin{proof}
		Soit $d$ une direction arbitraire, on $a$ :
		
		$$
		\begin{aligned}
		f\left(x^{*}+t d\right) & =f\left(x^{*}\right)+t d^{T} \nabla f\left(x^{*}\right)+\frac{1}{2} t^{2} d^{T} \nabla^{2} f\left(x^{*}\right) d+o\left(\|t d\|^{2}\right) \\
		& \Longrightarrow f\left(x^{*}+t d\right)-f\left(x^{*}\right)=t^{2}\left(\frac{1}{2} d^{T} \nabla^{2} f\left(x^{*}\right) d+o\left(\frac{\|t d\|^{2}}{t^{2}}\right)\right)
		\end{aligned}
		$$
		
		pour $t$ suffisamment petit on aura :
		
		$$
		f\left(x^{*}+t d\right) \geq f\left(x^{*}\right) \Longrightarrow x^{*} \text { est un minimum local de } f
		$$
	\end{proof}
	
	\begin{remark}[Étude des points critiques]\label{rem:points-critiques}
		Les deux derniers résultats permettent la plupart du temps de déterminer la nature d’un point critique
		$x$ d’une fonction $f$ (i.e. une solution $x$ de l’équation $\nabla f(x) = 0$).
		
		\begin{itemize}
			\item Si la matrice hessienne $\nabla^2 f(x)$ est définie positive, $x$ est un \textbf{minimum local strict} (Théorème 2.2.7).
			\item Si la matrice hessienne $\nabla^2 f(x)$ est définie négative, $x$ est un \textbf{maximum local strict} (Théorème 2.2.7 appliqué à $-f$).
			\item Si la matrice hessienne a des valeurs propres de signes différents, $x$ n’est ni un maximum ni un minimum local (par contraposée du Théorème 2.2.5 2.(Condition nécessaire du second ordre)). On parle alors de \textbf{point selle} ou \textbf{point col} (formellement, un point selle est un point critique où $f$ ne présente pas d’extrémum local).
		\end{itemize}
	\end{remark}
	
	\subsubsection{Condition nécessaire de minimum local sur un ensemble convexe :}
	\begin{theorem}
		Soit $f : \Omega \to \mathbb{R}$ différentiable en un point $x \in U$, où $U \subset \Omega$ est convexe.  
		Si $x$ est un minimum local de $f$ sur $U$, alors pour tout $y \in U$ :
		\[
		\langle \nabla f(x), y - x \rangle \ge 0
		\]
		ou encore
		\[
		df(x)(y-x) \ge 0.
		\]
		
		\textbf{Cas particulier : sous-espace affine}  
		Si $U = x + F$ avec $F$ un sous-espace vectoriel de $\mathbb{R}^n$, alors :
		\[
		\langle \nabla f(x), y - x \rangle = 0, \quad \forall y \in U.
		\]
	\end{theorem}
	\begin{proof} 
		Soit $y = x + h \in U$. Comme $U$ est convexe, tous les points $x + \theta h$ pour $\theta \in [0,1]$ appartiennent à $U$.  
		Par différentiabilité en $x$ :
		\[
		f(x + \theta h) - f(x) = \theta \langle \nabla f(x), h \rangle + \theta \|h\| \varepsilon(\theta h), 
		\quad \text{avec } \lim_{\theta \to 0} \varepsilon(\theta h) = 0.
		\]
		Comme $x$ est un minimum local, le membre de gauche est $\ge 0$, ce qui implique :
		\[
		\langle \nabla f(x), h \rangle \ge 0.
		\]
		
%		Dans le cas d’un sous-espace affine $U = x + F$, si $x+h \in U$, alors $x-h \in U$ également. On a :
%		\[
%		\langle \nabla f(x), h \rangle \ge 0 \quad \text{et} \quad \langle \nabla f(x), -h \rangle \ge 0 \implies \langle \nabla f(x), h \rangle = 0.
%		\]
		
		\textbf{Cas des sous-espaces affines $U = x + F$ :}  
		Si $x+h \in U$, alors $x-h \in U$ également. Ainsi :
		
		\[
		0 \le f(x + h) - f(x) = \langle \nabla f(x), h \rangle + \|h\| \varepsilon(h) \implies \langle \nabla f(x), h \rangle \ge 0,
		\]
		
		\[
		0 \le f(x - h) - f(x) = \langle \nabla f(x), -h \rangle + \|h\| \varepsilon(-h) \implies \langle \nabla f(x), -h \rangle \ge 0.
		\]
		
		On obtient donc la double inégalité :
		
		\[
		\langle \nabla f(x), h \rangle \ge 0 \quad \text{et} \quad \langle \nabla f(x), -h \rangle \ge 0 \implies \langle \nabla f(x), h \rangle = 0,
		\]
		
		ce qui donne finalement :
		
		\[
		\langle \nabla f(x), y - x \rangle = 0.
		\]
	\end{proof}
	
	\begin{remark}[Interprétation géométrique]
		Si $x$ est un minimum local sur un convexe $U$ et $\nabla f(x) \neq 0$, alors le gradient $\nabla f(x)$ pointe vers l’intérieur du convexe, c’est-à-dire qu’il forme un angle aigu avec tout vecteur $y-x$ pour $y \in U$.  
		
		Dans le cas d’un sous-espace affine $U = x + F$, le gradient est orthogonal à $F$ : 
		\[
		\nabla f(x) \in F^\perp.
		\]
	\end{remark}
	
	\subsubsection{Conditions nécessaires et suffisantes :}
	\begin{theorem} $f: \Re^{n} \rightarrow \Re$, une fonction convexe et différentiable, alors une condition nécessaire et suffisante pour que $x^{*} \in \Re^{n}$ soit un minimum global de $f$ sur $\Re^{n}$ est que :
		
		$$
		\nabla f\left(x^{*}\right)=0 .
		$$
	\end{theorem}
	\begin{proof}
		$f$ est convexe alors on $a$ :
		
		$$
		\left\{\begin{array}{c}
		f(x) \geq f\left(x^{*}\right)+\left(x-x^{*}\right)^{T} \nabla f\left(x^{*}\right) \\
		\forall x \in \Re^{n}
		\end{array}\right.
		$$
		
		\begin{itemize}
			\item $x^{*}$ est un minimum global de $f$ sur $\Re^{n}$, d'où :
		\end{itemize}
		
		$$
		\begin{aligned}
		f(x) & \geq f\left(x^{*}\right), \forall x \in \Re^{n}, \\
		& \Longrightarrow\left(x-x^{*}\right)^{T} \nabla f\left(x^{*}\right)=0, \forall x \in \Re^{n} \\
		& \Longrightarrow \nabla f\left(x^{*}\right)=0 .
		\end{aligned}
		$$
		
		$-\nabla f\left(x^{*}\right)=0 \Longrightarrow f(x) \geq f\left(x^{*}\right), \forall x \in \Re^{n}$, alors $x^{*}$ est un minimum global de $f$.\\
	\end{proof}
	
	\begin{exemple} Trouver analytiquement les extrema de la fonction :
		
		$$
		f(x, y)=x^{3}+y^{3}+3 x y
		$$
		
		et déterminer s'il s'agit de minimum ou autres.
	\end{exemple}
	\begin{solution}
		On a :
		
		$$
		\begin{gathered}
		\nabla f(x, y)=\binom{3 x^{2}+3 y}{3 y^{2}+3 x}=\binom{0}{0} \\
		\Longrightarrow\left\{\begin{array}{c}
		\binom{x}{y}=\binom{0}{0} \\
		\text { ou } \\
		\binom{x}{y}=\binom{-1}{-1}
		\end{array}\right.
		\end{gathered}
		$$
		
		$e t$
		
		$$
		\nabla^{2} f(x, y)=\left(\begin{array}{cc}
		6 x & 3 \\
		3 & 6 y
		\end{array}\right)
		$$
		
		alors
		
		$$
		\nabla^{2} f\binom{0}{0}=\left(\begin{array}{ll}
		0 & 3 \\
		3 & 0
		\end{array}\right)
		$$
		
		d'où :
		
		$$
		\left\langle\nabla^{2} f\binom{0}{0}\binom{x}{y},\binom{x}{y}\right\rangle=6 x y
		$$
		
		donc $\nabla^{2} f\binom{0}{0}$ n'est pas définie $\Longrightarrow\binom{0}{0}$ n'est pas extrémal.
		
		$$
		\nabla^{2} f\binom{-1}{-1}=\left(\begin{array}{cc}
		-6 & 3 \\
		3 & -6
		\end{array}\right)
		$$
		
		ses valeurs propres sont les racines de :
		
		$$
		\lambda^{2}+12 \lambda+27=0 \Longrightarrow \lambda_{1}=-9, \text { et }, \lambda_{2}=-3
		$$
		
		$\lambda_{1}, \lambda_{2}<0 \Longrightarrow \nabla^{2} f\binom{-1}{-1}$ est définie négative, dans ce cas le point $\binom{-1}{-1}$ est un maximum local pour $f$ et on a $f\binom{-1}{-1}=1$.
	\end{solution}
	
	\section{Caractérisation des fonctions convexes différentiables et deux fois différentiables}
	
	\begin{theorem}[Convexité et dérivabilité première]\label{thm:conv_deriv1}
		Soit $f : \Omega \to \R$ une fonction différentiable dans l’ouvert $\Omega$ et soit $U \subset \Omega$ un sous-ensemble convexe.  
		\begin{enumerate}%[label=\alph*)]
			\item La fonction $f$ est convexe sur $U$ si et seulement si pour tout $x,y \in U$,
			\[
			f(y) \geq f(x) + \langle \nabla f(x), y - x \rangle
			\quad \text{(ou encore $f(y) \geq f(x) + df(x)(y-x)$).}
			\]
			\item La fonction $f$ est strictement convexe sur $U$ si et seulement si pour tout $x,y \in U$, $x \neq y$,
			\[
			f(y) > f(x) + \langle \nabla f(x), y - x \rangle
			\quad \text{(ou encore $f(y) > f(x) + df(x)(y-x)$).}
			\]
		\end{enumerate}
	\end{theorem}
	
	\begin{proof}%[Démonstration]
		\begin{enumerate}%[label=\alph*)]
			\item $(\Rightarrow)$ : Soient $x,y$ deux points distincts de $U$ et $\theta \in ]0,1[$. Comme $f$ est convexe,
			\[
			f((1-\theta)x + \theta y) \leq (1-\theta)f(x) + \theta f(y).
			\]
			Donc
			\[
			\frac{f(x+\theta(y-x)) - f(x)}{\theta} \leq f(y) - f(x).
			\]
			En passant à la limite $\theta \to 0$, on obtient
			\[
			\langle \nabla f(x), y-x \rangle \leq f(y) - f(x).
			\]
			
			$(\Leftarrow)$ : Réciproquement, supposons que pour tout $x,y \in U$,  
			$f(y) \geq f(x) + \langle \nabla f(x), y-x \rangle$.  
			Soient $x,y \in U$, $\theta \in ]0,1[$. En appliquant l’inégalité aux deux couples 
			$(\theta x + (1-\theta)y,y)$ et $(\theta x + (1-\theta)y,x)$, on obtient
			\[
			\theta f(x) + (1-\theta)f(y) \geq f(\theta x + (1-\theta)y),
			\]
			donc $f$ est convexe.
			
			\item Le raisonnement est analogue en remplaçant les inégalités larges par des strictes, avec un raffinement au passage à la limite (utilisation de deux poids $0 < \theta < \omega < 1$).
		\end{enumerate}
	\end{proof}
	
	\begin{theorem}[Convexité et dérivabilité seconde]\label{thm:conv_deriv2}
		Soit $f : \Omega \to \R$ une fonction deux fois différentiable et soit $U \subset \Omega$ un sous-ensemble convexe.
		\begin{enumerate}%[label=\alph*)]
			\item $f$ est convexe sur $U$ si et seulement si pour tout $x,y \in U$,
			\[
			\langle \nabla^2 f(x)(y-x), \, y-x \rangle \geq 0.
			\]
			\item $f$ est strictement convexe sur $U$ si et seulement si pour tout $x,y \in U$, $x \neq y$,
			\[
			\langle \nabla^2 f(x)(y-x), \, y-x \rangle > 0.
			\]
		\end{enumerate}
		En particulier, si $\Omega = U$ est un ouvert convexe :
		\begin{enumerate}%[label=\alph*)]
			\item $f$ est convexe sur $\Omega$ si et seulement si pour tout $x \in \Omega$, la matrice hessienne $\nabla^2 f(x)$ est semi-définie positive.
			\item $f$ est strictement convexe sur $\Omega$ si et seulement si pour tout $x \in \Omega$, la matrice hessienne $\nabla^2 f(x)$ est définie positive.
		\end{enumerate}
	\end{theorem}
	
	\begin{proof}%[Démonstration]
		Soient $x$ et $y=x+h$ deux points distincts de $U$. Par la formule de Taylor-Maclaurin, il existe $z \in ]x,x+h[$ tel que
		\[
		f(x+h) = f(x) + \langle \nabla f(x), h \rangle + \tfrac12 \langle \nabla^2 f(z)h, h \rangle.
		\]
		Si $\langle \nabla^2 f(z)h, h \rangle \geq 0$, on retrouve la condition de convexité du Théorème \ref{thm:conv_deriv1}.  
		Le raisonnement complet montre que $\nabla^2 f(x)$ doit être semi-définie positive (resp. définie positive).
	\end{proof}
	
	\section{Problèmes d’optimisation convexes}
	
	On rappelle qu’un problème d’optimisation
	\[
	\text{Trouver } x^* \in U \quad \text{tel que} \quad f(x^*) = \displaystyle\min_{x \in U} f(x),
	\]
	est dit \textbf{convexe} si $U$ et $f$ sont convexes.  
	
	\begin{theorem}[Minimum de fonctions convexes]\label{thm:min_conv}
		Soit $U \subset \R^n$ un ensemble convexe.
		\begin{enumerate}%[label=\alph*)]
			\item Si une fonction convexe $f:U\to \R$ admet un minimum local en un point $x$, elle y admet en fait un minimum global sur $U$.
			\item Une fonction $f:U\to\R$ strictement convexe admet au plus un minimum local, qui est en fait un minimum global strict.
			\item Soit $f:\Omega \to \R$ convexe et différentiable, définie sur un ouvert convexe $\Omega \subset \R^n$. Alors $x \in \Omega$ est un minimum global de $f$ si et seulement si $\nabla f(x)=0$.
			\item Soit $f:\Omega \to \R$ définie sur un ouvert $\Omega \supset U$, convexe sur $U$. Alors $x \in U$ est un minimum de $f$ sur $U$ si et seulement si
			\[
			\langle \nabla f(x), y-x \rangle \geq 0, \quad \forall y \in U.
			\]
			En particulier, si $U = x+F$ est un sous-espace affine, alors $x \in U$ est un minimum de $f$ sur $U$ si et seulement si
			\[
			\langle \nabla f(x), y-x \rangle = 0, \quad \forall y \in U.
			\]
		\end{enumerate}
	\end{theorem}
	
	\begin{proof}%[Démonstration]
		\begin{enumerate}%[label=\alph*)]
			\item Par convexité et définition d’un minimum local, on déduit que $x$ est en fait un minimum global.
			\item Si $f$ est strictement convexe et admet un minimum local $x$, alors ce minimum est unique et strictement global.
			\item Si $\nabla f(x)=0$, le Théorème \ref{thm:conv_deriv1} entraîne que $f(y)\geq f(x)$ pour tout $y$, donc $x$ est minimum global.
			\item Même raisonnement, mais restreint à $U$.
		\end{enumerate}
	\end{proof}
	
	\begin{remark}
		\begin{itemize}
			\item Une fonction non strictement convexe peut admettre plusieurs minimums locaux. Cependant, l’ensemble des minimums globaux forme toujours un ensemble convexe.
			\item Dans un problème d’optimisation sans contrainte $\displaystyle\min_{x \in \R^n} f(x)$, si $f$ est différentiable, coercive et strictement convexe, alors il admet une unique solution $x^*$, qui est l’unique solution de l’équation $\nabla f(x)=0$.
		\end{itemize}
	\end{remark}
	
	\subsection{Exercices }
%	\subsubsection{Exercice 1 :}
	\begin{exercise}\label{Exo1chap1}
		Pour chacune des fonctions suivantes:\\
		\begin{enumerate}
			\item $f_{1}\left(x_{1}, x_{2}, x_{3}\right)=3 x_{1}^{2}+3 x_{2}^{2}+3 x_{3}^{2}-2 x_{1}-10 x_{3}-2 x_{1} x_{3}$.\\
			\item $f_{2}\left(x_{1}, x_{2}\right)=x_{1}^{3}+x_{2}^{3}-3 x_{1}-12 x_{2}+20$.\\
			\item $f_{3}\left(x_{1}, x_{2}\right)=x_{1}^{4}+x_{2}^{4}-2\left(x_{1}-x_{2}\right)^{2}$.\\
			\begin{enumerate}
				\item Étudier l'existence de points extrémauts.\\
				\item En utilisant une condition d'optimalité du premier ordre déterminer les points critiques.\\
				\item Préciser à chaque fois leur nature (minimum ou maximum ? local ou global ?).
			\end{enumerate}
		\end{enumerate}
	\end{exercise}
	
%	\subsubsection{Exercice 2 :}
	\begin{exercise}\label{Exo2chap1}
		$f$ est une fonction définie sur $\Re^{2}$ par :
		
		$$
		f(x, y)=x^{4}+y^{4}-2(x-y)^{2}
		$$
		
		\begin{enumerate}
			\item Montrer qu'il existe $(\alpha, \beta) \in \Re^{2}$ et les déterminer tels que :
			
			$$
			f(x, y) \geq \alpha\|(x, y)\|_{2}^{2}+\beta, \quad \forall(x, y) \in \Re^{2}
			$$
			
			En déduire que le problème :
			
			$$
			\displaystyle\min _{(x, y) \in \Re^{2}} f(x, y)
			$$
			
			admet au moins une solution.\\
			\item  La fonction $f$ est -elle convexe sur $\Re^{2}$ ?.\\
			\item Déterminer les points critiques de $f$, et préciser leur nature.\\
		\end{enumerate}
	\end{exercise}
	\begin{exercise}\label{Exo3chap1}
%		Exercice 3 :\\
		Soit $Q$, une matrice de $M_{n}(\Re)$ symétrique et définie positive. Soit $b$, un vecteur de $\Re^{n}$.
		
		\begin{enumerate}
			\item Démontrer qu'il existe $\delta>0$ tel que pour tout vecteur $x$ de $\Re^{n},\langle Q x, x\rangle \geq \delta\|x\|^{2}$.
			\item Soit $f$, la fonction définie sur $\Re^{n}$ par :
	
		
		$$
		f(x)=\frac{1}{2}\langle Q x, x\rangle-\langle b, x\rangle ;
		$$
		
		pour $x \in \Re^{n}$.\\
		\begin{enumerate}
			\item Démontrer que $f$ est coercive.\\
			\item  Démontrer que le problème :
			
			$$
			(P)\left\{\begin{array}{c}
			\displaystyle\min f(x) \\
			x \in \Re^{n}
			\end{array}\right.
			$$
			
			possède une solution unique que l'on déterminera complètement.
		\end{enumerate}
		\end{enumerate}
	\end{exercise}
	
	
	
	\chapter{Optimisation sous contraintes d'égalité : Multiplicateurs de Lagrange}
	\label{chap:lagrange}
	
	Après avoir étudié l'optimisation sans contrainte, où l'on cherche un extremum sur $\R^n$ tout entier, nous abordons les problèmes où les solutions doivent appartenir à un sous-ensemble strict de l'espace. Le cas le plus simple et le plus fondamental est celui où ce sous-ensemble est défini par des contraintes d'égalité. La méthode des multiplicateurs de Lagrange, développée au 18ème siècle, est un outil élégant et puissant pour résoudre ce type de problèmes.
	
	\section{Position du problème et intuition géométrique}
	
	Un problème d'optimisation sous contraintes d'égalité se formule comme suit :
	$$
	(P_E) \quad
	\begin{array}{ll}
	\text{minimiser (ou maximiser)} & f(x) \\
	\text{pour} & x \in \R^n \\
	\text{sous contrainte} & g_i(x) = 0, \quad i=1, \dots, m
	\end{array}
	$$
	où $f: \R^n \to \R$ est la \textbf{fonction objectif} (ce que l'on veut optimiser) et $g: \R^n \to \R^m$ (avec $m < n$) est la \textbf{fonction de contrainte}. L'ensemble des points admissibles (ceux qui respectent les contraintes) est noté $\mathcal{C} = \{x \in \R^n \mid g(x)=0\}$.
	
	\paragraph{Intuition géométrique .}
	Imaginons que $f(x,y)$ représente l'altitude sur une carte, et que la contrainte $g(x,y)=0$ représente un sentier que vous êtes obligé de suivre. Chercher le minimum de $f$ sous la contrainte $g$ revient à trouver le point le plus bas sur ce sentier.
	
	En un point optimal $x^*$, vous ne pouvez ni monter ni descendre en suivant le sentier. Cela signifie que la direction de la plus grande pente, donnée par le gradient $\grad f(x^*)$, doit être \textbf{perpendiculaire au sentier}. Si ce n'était pas le cas, le gradient aurait une composante le long du sentier, et en vous déplaçant dans la direction opposée, vous pourriez descendre plus bas, ce qui contredirait l'optimalité de $x^*$.
	
	Or, nous savons que le gradient d'une fonction est toujours orthogonal à ses lignes de niveau. Le sentier $\mathcal{C}$ est précisément la ligne de niveau zéro de la fonction $g$. Le gradient $\grad g(x^*)$ est donc lui aussi orthogonal au sentier.
	
	Puisque $\grad f(x^*)$ et $\grad g(x^*)$ sont tous deux orthogonaux au même sentier (dans le cas d'une seule contrainte), ils doivent être \textbf{colinéaires}. Il existe donc un scalaire, noté $\lambda^*$, tel que :
	$$\grad f(x^*) + \lambda^* \grad g(x^*) = 0$$
	Cette colinéarité est l'idée fondamentale de la méthode de Lagrange.
	
	\section{Le théorème de Lagrange}
	
	\begin{definition}[Point régulier]
		Un point admissible $x^* \in \mathcal{C}$ est un \textbf{point régulier} si les vecteurs gradients des contraintes en ce point, $\{\grad g_1(x^*), \dots, \grad g_m(x^*)\}$, forment une famille linéairement indépendante.
	\end{definition}
	
	Autrement dit, les contraintes sont qualifiées au point \(x^*\) si et seulement si la famille
	\(\{\nabla g_i(x^*) ,\;1\le i\le p\}\) est une famille libre.
	
	\begin{remark}
		\textit{En clair : cette condition technique évite les points "pathologiques" sur la contrainte (comme un point de rebroussement). Elle garantit que les gradients des contraintes décrivent correctement la géométrie locale de l'ensemble admissible.}
	\end{remark}
	
	\begin{definition}[Lagrangien]
		Le \textbf{Lagrangien} associé au problème $(P_E)$ est la fonction $\mathcal{L}: \R^n \times \R^m \to \R$ définie par :
		$$\mathcal{L}(x, \lambda) = f(x) + \sum_{i=1}^m \lambda_i g_i(x) = f(x) + \lambda\T g(x)$$
		Le vecteur $\lambda \in \R^m$ est appelé le vecteur des \textbf{multiplicateurs de Lagrange}.
	\end{definition}
	\begin{remark}
		\textit{Astuce : on construit le Lagrangien en prenant la fonction objectif et en y "ajoutant" chaque contrainte, multipliée par son propre multiplicateur $\lambda_i$.}
	\end{remark}
	
	On commence par des conditions nécessaires :
	
	\begin{theorem}[Condition nécessaire, ordre 1]
		On suppose que \(f\) et \(g_i,\;1\le i\le p\) sont continûment différentiables sur \(D\).
		
		Si \(f\) restreinte à \(C\) présente un extrémum en \(x^*\in C\) et si les contraintes sont qualifiées,
		alors le lagrangien admet un point critique \((x^*,\lambda^*)\) avec \(\lambda^*=(\lambda^*_1,\lambda^*_2,\dots,\lambda^*_p)\) tel que
		\[
		\nabla L(x^*,\lambda^*)=0_{\mathbb{R}^{n+p}}.
		\]
		Les \(\lambda^*_i\) sont appelés multiplicateurs de Lagrange.
	\end{theorem}
	
	\begin{remark}
		\begin{enumerate}
			\item Si \((x^*,\lambda^*)\) est un point critique du lagrangien, alors \(x^*\in C\) car \(\dfrac{\partial L}{\partial \lambda_j}(x^*,\lambda^*)=g_j(x^*)=0\).
			\item Si la contrainte n'est pas qualifiée en \(x^*\), il se peut que \(f\) restreinte à \(C\) présente un extrémum en \(x^*\in C\) mais que les multiplicateurs de Lagrange \(\lambda^*\) n'existent pas.
		\end{enumerate}
	\end{remark}
	
	La version « faible » pour les conditions suffisantes :
	
	\begin{theorem}[Condition suffisante d'ordre 2]
		On suppose que \(f\) et \(g_j,\;1\le j\le p\) sont deux fois continûment différentiables sur \(D\) et que
		\begin{itemize}
			\item \((x^*,\lambda^*)\) est un point critique du lagrangien,
			\item les contraintes en \(x^*\) sont qualifiées,
			\item la matrice \(\nabla^2_x L(x^*,\lambda^*)\) est définie positive,
		\end{itemize}
		alors \(x^*\) est un minimum du problème contraint.
	\end{theorem}
	
	\begin{example}[Optimisation sous contrainte]\label{eq3.5}
		On veut optimiser
		\[
		f(x,y)=x^{3}y - x^{2} + 3x + y^{2} + 2y
		\]
		sous la contrainte
		\[
		g(x,y)=2x+y-1=0.
		\]
		
		Le Lagrangien associé est
		\[
		L(x,y,\lambda)=x^{3}y - x^{2} + 3x + y^{2} + 2y + \lambda(2x+y-1).
		\]
		
		\begin{enumerate}
			\item \textbf{Gradient et hessienne bordée.} \\
			\[
			\nabla L(x,y,\lambda)=
			\begin{pmatrix}
			3x^{2}y - 2x + 3 + 2\lambda\\[4pt]
			x^{3} + 2y + 2 + \lambda\\[4pt]
			2x + y - 1
			\end{pmatrix}.
			\]
			La hessienne bordée (matrice symétrique ajoutant la contrainte) s'écrit
			\[
			\nabla^{2}L(x,y,\lambda)=
			\begin{pmatrix}
			6xy - 2 & 3x^{2} & 2 \\[4pt]
			3x^{2}  & 2      & 1 \\[4pt]
			2      & 1      & 0
			\end{pmatrix}.
			\]
			
			\item \textbf{Point critique.} \\
			On trouve un seul point critique de $L$ :
			\[
			(x,y,\lambda)=(-1,\,3,\,7).
			\]
			En ce point,
			\[
			\nabla^{2}L(-1,3,7)=
			\begin{pmatrix}
			-20 & 3 & 2 \\[4pt]
			3  & 2 & 1 \\[4pt]
			2  & 1 & 0
			\end{pmatrix}.
			\]
			
			\item \textbf{Qualification de la contrainte.} \\
			\[
			\nabla g(x,y)=\begin{pmatrix}2\\[2pt]1\end{pmatrix},
			\qquad \nabla g(-1,3)=\begin{pmatrix}2\\[2pt]1\end{pmatrix}\neq 0,
			\]
			donc la contrainte est qualifiée en $(-1,3)$.
			
			\item \textbf{Déterminant et conclusion.} \\
			\[
			\det\!\big(\nabla^{2}L(-1,3,7)\big)=24>0.
			\]
			Par le critère des multiplicateurs (et la signature de la hessienne bordée), on en déduit que $f$ présente un \emph{maximum contraint} en $( -1,3)$.  
			De plus,
			\[
			f(-1,3)=(-1)^{3}\cdot 3 - (-1)^{2} + 3(-1) + 3^{2} + 2\cdot 3 = 8.
			\]
		\end{enumerate}
	\end{example}
	
	\begin{example}\label{ex:3.6}
		Dans l'exemple \eqref{eq3.5}, \(\nabla^2_x L(-1,3,7)=
		\begin{pmatrix}
		-20 & 3\\
		3 & 2
		\end{pmatrix}\)
		a une valeur propre négative, l'autre positive. Ce théorème ne permet pas de conclure.
	\end{example}
	
	Dans l'exemple suivant, le théorème permet de conclure :
	
	\begin{example}\label{ex:3.7}
		Minimiser \(f(x,y,z)=x+y+z\) sous la contrainte \(g(x,y,z)=x^2+y^2+z^2-r^2\).
		
		\begin{enumerate}
			\item Le lagrangien est
			\[
			L(x,y,z,\lambda)=x+y+z+\lambda\big(x^2+y^2+z^2-r^2\big).
			\]
			Le gradient est
			\[
			\nabla L(x,y,z,\lambda)=
			\begin{pmatrix}
			1+2\lambda x\\[4pt]
			1+2\lambda y\\[4pt]
			1+2\lambda z\\[4pt]
			x^2+y^2+z^2-r^2
			\end{pmatrix},
			\]
			et la Hessienne (par rapport à \((x,y,z,\lambda)\) structurée ci-dessous) :
			\[
			\nabla^2 L(x,y,z,\lambda)=
			\begin{pmatrix}
			2\lambda & 0 & 0 & 2x\\
			0 & 2\lambda & 0 & 2y\\
			0 & 0 & 2\lambda & 2z\\
			2x & 2y & 2z & 0
			\end{pmatrix},
			\]
			dont le déterminant est
			\[
			\det\big(\nabla^2 L(x,y,z,\lambda)\big) = -16\lambda^2\big(x^2+y^2+z^2\big) < 0.
			\]
			
			\item Points critiques : en annulant les composantes du gradient on trouve
			\(x=y=z=-\dfrac{1}{2\lambda}\), et la contrainte donne \(3x^2=r^2\). Les points critiques sont donc
			\[
			A=(x^*,\lambda^*)=\Big(\tfrac{r}{\sqrt{3}},\tfrac{r}{\sqrt{3}},\tfrac{r}{\sqrt{3}}, -\tfrac{\sqrt{3}}{2r}\Big)
			\quad\text{et}\quad B=-A.
			\]
			
			Vérifions que la contrainte est qualifiée pour chacun de ces points : \(\nabla g(x,y,z) = (2x,2y,2z)^T\),
			donc \(\nabla g(x^*)\neq 0\) et \(\nabla g(-x^*)\neq 0\) ; ces deux points critiques sont donc candidats à être des extrémas contraints.
			
			\item De l'égalité \(\nabla^2_x L(x,y,z,\lambda) = 2\lambda I_3\) on déduit :
			\begin{itemize}
				\item au point \(A\) les valeurs propres sont \(\lambda_1=\lambda_2=\lambda_3=-\dfrac{\sqrt{3}}{r}<0\) donc un maximum ;
				\item au point \(B=-A\) les valeurs propres sont \(\lambda_1=\lambda_2=\lambda_3=\dfrac{\sqrt{3}}{r}>0\) donc un minimum.
			\end{itemize}
		\end{enumerate}
	\end{example}
	
	Voici une version plus raffinée de ce théorème :
	
	\begin{theorem}[Condition suffisante d'ordre 2 — version raffinée]
		On suppose que \(f\) et \(g_1,g_2,\dots,g_p\) sont deux fois continûment différentiables sur \(D\) et que
		\begin{enumerate}
			\item \((x^*,\lambda^*)\) est un point critique du lagrangien,
			\item les contraintes en \(x^*\) sont qualifiées,
			\item pour tout \(h\in\mathbb{R}^n\) tel que \(\forall i\in\{1,\dots,p\},\; \nabla g_i(x^*)\cdot h = 0\), on a
			\[
			h^T\big(\nabla^2_x L(x^*,\lambda^*)\big) h > 0.
			\]
		\end{enumerate}
		Alors \(x^*\) est un minimum du problème contraint.
		
		Si on remplace la troisième condition par \(h^T\big(\nabla^2_x L(x^*,\lambda^*)\big) h < 0\), alors \(x^*\) est un maximum.
	\end{theorem}
	
	\begin{remark}
		La différence avec le théorème précédent tient à la condition 3 : on demande ici que la matrice \(\nabla^2_x L(x^*,\lambda^*)\) soit définie positive seulement pour les vecteurs \(h\) qui sont orthogonaux à tous les gradients des contraintes.
	\end{remark}
	
	\begin{proposition}[règle pratique]
		La troisième condition du théorème précédent est équivalente à : toutes les racines du polynôme
		\[
		p(\lambda)=\det\begin{pmatrix} A-\lambda I_n & D \\ D^T & 0 \end{pmatrix}
		\]
		sont positives, où \(I_n\) est la matrice identité de \(\mathbb{R}^{n\times n}\),
		\[
		A=\nabla^2_x L(x^*,\lambda^*)\in\mathbb{R}^{n\times n},\qquad
		D=(\nabla g_j(x^*))_{1\le j\le p}\in\mathbb{R}^{n\times p},
		\]
		(c'est-à-dire \(D\) est la matrice dont les \(p\) colonnes sont les gradients des contraintes).
	\end{proposition}
	
	\begin{example}\label{ex:3.8}
		Dans l'exemple (3.5),
		\[
		\nabla^2 L(-1,3,7)=
		\begin{pmatrix}
		-20 & 3 & -2\\
		3 & 2 & -1\\
		-2 & -1 & 0
		\end{pmatrix},
		\]
		et
		\[
		p(\lambda)=
		\begin{vmatrix}
		-20-\lambda & 3 & -2\\[2pt]
		3 & 2-\lambda & -1\\[2pt]
		-2 & -1 & -\lambda
		\end{vmatrix}=24+5\lambda,
		\]
		qui a pour racine \(\lambda=-\tfrac{24}{5}<0\). On a donc un maximum en \(f(-1,3)=8\).
	\end{example}
	
	\section{En résumé}
	
	Pour une ou plusieurs contraintes d'égalité, on utilisera la méthodologie suivante :
	
	\paragraph{Contraintes égalités}
	\begin{enumerate}
		\item Déterminer les points critiques du lagrangien.
		\item Pour chaque point critique \((x^*,\lambda^*)\) :
		\begin{enumerate}
			\item vérifier que les contraintes en \(x^*\) sont qualifiées,
			\item trouver les racines du polynôme
			\[
			p(t)=\det\begin{pmatrix} A-tI_n & D \\ D^T & 0\end{pmatrix}.
			\]
			\begin{itemize}
				\item[i.] si les racines sont toutes strictement positives, on a un minimum en \(f(x^*)\);
				\item[ii.] si les racines sont toutes strictement négatives, on a un maximum en \(f(x^*)\);
				\item[iii.] sinon on ne peut pas conclure.
			\end{itemize}
		\end{enumerate}
	\end{enumerate}
	
	
	\begin{exercise}
		\textbf{Énoncé.} Minimiser 
		\[
		f(x,y) = x^2 + y^2
		\]
		sous la contrainte 
		\[
		x + y = 1.
		\]
	\end{exercise}
	
	\begin{exercise}
		\textbf{Énoncé.} Trouver les extrema (max et min) de 
		\[
		f(x,y) = xy
		\]
		sur la contrainte 
		\[
		x^2 + y^2 = 1.
		\]
	\end{exercise}
	
	\begin{exercise}
		\textbf{Énoncé.} Minimiser 
		\[
		f(x,y,z) = x^2 + y^2 + z^2
		\]
		sous les contraintes
		\[
		\begin{cases}
		x + y + z = 1, \\
		x - y = 0.
		\end{cases}
		\]
	\end{exercise}
	
	\begin{exercise}[contre-exemple : contraintes non qualifiées]
		\textbf{Énoncé.} Soit $f(x,y) = x$ et la contrainte $g(x,y) = x^2 + y^2 = 0$. Trouver les extrema sur l’ensemble 
		\[
		C = \{ (x,y) : x^2 + y^2 = 0 \}.
		\]
	\end{exercise}
	
	\begin{exercise}
		\textbf{Énoncé.} Minimiser 
		\[
		f(x,y,z) = x + y + z
		\]
		sous les contraintes
		\[
		\begin{cases}
		g_1(x,y,z) = x + y + z - 1 = 0, \\
		g_2(x,y,z) = 2x + 2y + 2z - 2 = 0.
		\end{cases}
		\]
	\end{exercise}
	
	%===============================================================
	\chapter{Théorème de Karush-Kuhn-Tucker (KKT)}
	\label{chap:kkt}
	
	Ce chapitre généralise l'approche de Lagrange aux problèmes qui comportent à la fois des contraintes d'égalité et d'inégalité. Le résultat central, connu sous le nom de conditions de Karush-Kuhn-Tucker (KKT), est la pierre angulaire de l'optimisation non linéaire.
	
	\section{Conditions d’optimalité non qualifiées}
	
	Dans cette section étendue, nous caractérisons les solutions d'un problème d'optimisation sous contraintes mixtes (égalités et inégalités). Nous considérons le problème général de minimisation suivant, qui étend le cas des contraintes d'égalité vu précédemment. Ce cadre est essentiel pour modéliser des problèmes réels où les variables doivent satisfaire à la fois des égalités strictes (comme des bilans massiques en ingénierie) et des inégalités (comme des bornes sur les ressources). Nous rappelons les bases avant d'approfondir.
	
	Nous considérons maintenant des problèmes de la forme générale :
	$$
	(P) \quad
	\begin{array}{ll}
	\text{minimiser} & f(x) \\
	\text{s.c.} & h_i(x) = 0, \quad i=1, \dots, p \\
	& g_j(x) \le 0, \quad j=1, \dots, q
	\end{array}
	$$
	l'ensemble admissible est $U=\{x\in\mathbb{R}^{n} \mid h(x)=0 \text{ et } g(x)\le0\}$.
	
	
	Nous supposons que les fonctions \( f \), \( h \) et \( g \) sont de classe \( C^1 \) (continûment différentiables) sur un ouvert \( \Omega \) contenant \( U \), pour garantir l'existence et la continuité des gradients. Pour des analyses plus avancées, nous pourrons supposer \( C^2 \) si nécessaire pour les conditions du second ordre.
	
	\subsection{Théorème de Fritz John : Conditions nécessaires non qualifiées}
	
	Le théorème de Fritz John fournit des conditions nécessaires d'optimalité sans imposer de conditions supplémentaires sur les contraintes (d'où le terme "non qualifiées"). Ces conditions généralisent les multiplicateurs de Lagrange aux problèmes avec inégalités. Elles sont particulièrement utiles lorsque les contraintes peuvent être dégénérées, mais elles nécessitent des qualifications pour être pleinement exploitables.
	
	\begin{theorem}[Conditions nécessaires d'optimalité non qualifiées ou théorème de Fritz John]\label{theo4.1.1}
		Supposons que \( x^* \in U \) est une solution locale du problème de minimisation. Alors il existe des vecteurs \( \lambda^* = (\lambda^*_1, \dots, \lambda^*_p)^\top \in \mathbb{R}^p \), \( \mu^* = (\mu^*_1, \dots, \mu^*_q)^\top \in \mathbb{R}^q_+ \) (composantes non négatives), et \( \mu^*_0 \in \mathbb{R}_+ \) (non négatif), non tous nuls (c'est-à-dire \( (\lambda^*_1, \dots, \lambda^*_p, \mu^*_0, \mu^*_1, \dots, \mu^*_q)^\top \in \mathbb{R}^{p+q+1} \setminus \{0\} \)), tels que :
		\[
		\mu^*_0 \nabla f(x^*) + \sum_{i=1}^p \lambda^*_i \nabla h_i(x^*) + \sum_{j=1}^q \mu^*_j \nabla g_j(x^*) = 0,
		\]
		et pour tout \( j \in \{1, 2, \dots, q\} \),
		\[
		\mu^*_j g_j(x^*) = 0.
		\]
		Cela signifie que \( \mu^*_j = 0 \) dès que \( g_j(x^*) < 0 \) (contrainte inactive).
	\end{theorem}
	
	\textbf{Interprétation géométrique :} À un point d'optimum local \( x^* \), le gradient \( \nabla f(x^*) \) doit être orienté de manière à ne pas permettre de descente dans l'ensemble admissible \( U \). Les gradients des contraintes actives définissent les directions normales à la frontière de \( U \). La combinaison linéaire exprime que \( \nabla f(x^*) \) (pondéré par \( \mu_0^* \)) est opposé à une combinaison des normales aux contraintes, assurant l'équilibre. Si \( \mu_0^* = 0 \), cela indique une dégénérescence des contraintes, où l'optimalité est imposée uniquement par les contraintes sans lien direct avec \( f \).
	
	\textbf{Dérivation intuitive :} Considérons le cône des directions admissibles en \( x^* \), qui sont les directions \( d \) telles que \( \langle \nabla h_i(x^*), d \rangle = 0 \) pour toutes contraintes d'égalité, et \( \langle \nabla g_j(x^*), d \rangle \leq 0 \) pour les contraintes actives (\( g_j(x^*)=0 \)). Pour que \( x^* \) soit un minimum local, on doit avoir \( \langle \nabla f(x^*), d \rangle \geq 0 \) pour tout tel \( d \). Par des arguments de séparation (théorème de Farkas ou dualité en optimisation convexe), cela implique l'existence des multiplicateurs comme ci-dessus, avec \( \mu_0^* \) potentiellement nul. Une preuve rigoureuse utilise le théorème des alternatives de Motzkin ou des lemmes de Farkas généralisés.
	
	\begin{example}[Application du théorème de Fritz John]
		Considérons le problème : minimiser \( f(x,y) = x \) sous \( g_1(x,y) = y \leq 0 \), \( g_2(x,y) = -y \leq 0 \). Ici, \( U = \{ (x,0) \mid x \in \mathbb{R} \} \), et tout point est optimal car \( f \) est constante sur \( y=0 \). À \( (0,0) \), les deux contraintes sont actives, mais dépendantes (\( \nabla g_2 = -\nabla g_1 \)). Le théorème donne \( \mu_0^* \partial_x f + \mu_1^* \partial_y g_1 + \mu_2^* \partial_y g_2 = (\mu_0^*, 0) + (0, \mu_1^*) + (0, -\mu_2^*) = 0 \), avec \( \mu_0^* = 0 \), \( \mu_1^* = \mu_2^* > 0 \) possible, illustrant la non-qualification.
	\end{example}
	
	\subsubsection{Vocabulaire et notations}
	
	Les réels \( \lambda^*_1, \dots, \lambda^*_p \) (associés aux égalités) et les réels positifs \( \mu^*_1, \dots, \mu^*_q \) (associés aux inégalités) sont appelés \textbf{multiplicateurs de Lagrange}. Notez que les \( \lambda_i^* \) peuvent être de signe quelconque, contrairement aux \( \mu_j^* \geq 0 \), qui reflètent le fait que les inégalités sont unidirectionnelles (poussant dans une direction).
	
	La relation \( \forall j \in \{1, \dots, q\}, \quad \mu^*_j g_j(x^*) = 0 \) est appelée \textbf{relation de complémentarité} (ou complementary slackness). Elle assure que seules les contraintes saturées (actives) contribuent à l'équation via leurs multiplicateurs non nuls. Cela reflète l'idée que les contraintes inactives n'influencent pas localement l'optimalité.
	
	\begin{definition}[Contraintes actives/inactives]
		Pour un point \( x \in U \), une contrainte d'inégalité \( g_j \) est dite \textbf{active} (ou saturée) si \( g_j(x) = 0 \) (le point est sur la frontière correspondante de \( U \)). Sinon, si \( g_j(x) < 0 \), elle est \textbf{inactive} (le point est à l'intérieur pour cette contrainte).
		
		On note \( I(x) = \{ j \in \{1, \dots, q\} \mid g_j(x) = 0 \} \) l'ensemble des indices des contraintes actives en \( x \).
	\end{definition}
	
	La relation de complémentarité implique que dans la combinaison linéaire du théorème, seuls les gradients \( \nabla g_j(x^*) \) pour \( j \in I(x^*) \) peuvent intervenir (car \( \mu_j^* = 0 \) sinon). Cela réduit le système à ne considérer que les contraintes pertinentes localement.
	
	
	\begin{example}[Illustration de contraintes actives]
		Considérons \( f(x,y) = x + y \), à minimiser sous \( g_1(x,y) = x \leq 0 \), \( g_2(x,y) = y \leq 0 \). À \( (0,0) \), les deux contraintes sont actives, et les multiplicateurs satisfont \( \mu_0 (1,1) + \mu_1 (1,0) + \mu_2 (0,1) = 0 \) avec \( \mu_1, \mu_2 \geq 0 \). Une solution est \( \mu_0 = 1 \), \( \mu_1 = -1 \), mais comme \( \mu_1 < 0 \) invalide, cela montre que (0,0) n'est pas un minimum (en fait, c'est un maximum). Correction : pour minimisation, le point optimal est à l'infini, mais si borné, ajuster.
	\end{example}
	
	\begin{remark}
		\begin{itemize}
			\item Les multiplicateurs de Lagrange ne sont pas uniques : on peut les multiplier par une constante positive sans changer les équations (propriété d'homogénéité due à la forme linéaire).
			\item Si \( \mu_0^* = 0 \), l'équation ne dépend plus de \( \nabla f(x^*) \), et ne fournit aucune information spécifique sur la fonction objectif. C'est pourquoi ces conditions sont dites "non qualifiées" : elles peuvent être satisfaites même pour des points non optimaux si les contraintes sont mal posées ou redondantes. Nous introduirons des qualifications pour garantir \( \mu_0^* > 0 \), menant aux conditions de Karush-Kuhn-Tucker (KKT).
			\item En pratique, si \( \mu_0^* > 0 \), on peut normaliser en divisant par \( \mu_0^* \), obtenant les conditions KKT classiques : \( \nabla f(x^*) + \sum \lambda_i^* \nabla h_i + \sum \mu_j^* \nabla g_j = 0 \), avec complémentarité et \( \mu_j^* \geq 0 \).
			\item Le théorème de Fritz John est valable pour des optima locaux, mais pour des optima globaux, des conditions supplémentaires comme la convexité sont nécessaires.
		\end{itemize}
	\end{remark}
	
	\section{Qualification des contraintes}
	
	Les conditions non qualifiées de Fritz John peuvent échouer à caractériser correctement l'optimalité si les contraintes sont "mal décrites", c'est-à-dire si leurs gradients sont linéairement dépendants, menant à des solutions triviales non informatives (avec \( \mu_0^* = 0 \)). Cette section introduit plusieurs conditions de qualification pour éviter cela, assurant que les multiplicateurs sont pertinents. Nous ne nous limitons pas à deux qualifications, mais en présentons plusieurs courantes, en commençant par les plus générales et en allant vers des cas spécifiques comme les problèmes convexes.
	
	\subsection{Mauvaises descriptions des contraintes}
	
	L'ensemble \( U \) reste inchangé si l'on modifie les contraintes de manière redondante, mais cela peut introduire des dépendances linéaires entre gradients, rendant les conditions non qualifiées inutiles.
	
	Par exemple, remplaçons \( h(x) = 0 \) par \( \phi(x) = \begin{pmatrix} h(x) \\ -h_1(x) \end{pmatrix} = 0 \in \mathbb{R}^{p+1} \). Alors \( \nabla \phi_{p+1}(x) = -\nabla h_1(x) = -\nabla \phi_1(x) \), créant une dépendance linéaire. En tout point \( x \in U \), on peut trouver une combinaison non triviale des gradients avec \( \lambda = (\lambda_1, 0, \dots, 0, \lambda_1)^\top \neq 0 \), et \( \mu_0 = 0 \), sans impliquer \( \nabla f \).
	
	Plus généralement :
	\begin{itemize}
		\item Ajouter une combinaison linéaire des \( h_i \) (égalités) crée des dépendances.
		\item Pour les inégalités, ajouter des combinaisons linéaires positives des \( g_j \) actives ne change pas \( U \), mais peut causer des dépendances.
		\item Des contraintes non lisses ou avec cusps (ex. \( g(x) = |x| \leq 0 \)) peuvent violer les qualifications.
	\end{itemize}
	
	Une "bonne" représentation des contraintes requiert donc des propriétés comme l'indépendance linéaire ou des conditions plus faibles sur les cônes tangents. Cependant, il n'est pas toujours possible d'obtenir une telle indépendance, d'où la nécessité de qualifications variées adaptées au contexte (convexe, non convexe, etc.).
	
	\begin{example}[Dépendance linéaire]
		Considérons deux contraintes d'égalité redondantes : \( h_1(x) = x_1 + x_2 = 0 \), \( h_2(x) = 2x_1 + 2x_2 = 0 \). Les gradients sont dépendants (\( \nabla h_2 = 2 \nabla h_1 \)), menant à des multiplicateurs triviaux avec \( \mu_0^* = 0 \), même si le point est optimal pour certaines \( f \).
	\end{example}
	
	\subsection{Qualifications des contraintes}
	
	Nous présentons plusieurs qualifications courantes, en ordre de généralité décroissante. Chacune garantit, sous hypothèses, que \( \mu_0^* > 0 \) dans le théorème de Fritz John, menant aux conditions KKT.
	
	\begin{definition}[Qualification des contraintes de Mangasarian-Fromovitz (QCMF)]\label{Def4.2.2}
		Un point réalisable \( x^* \in U \) vérifie la QCMF si :
		\begin{enumerate}
			\item[i)] Les vecteurs \( \nabla h_i(x^*) \), pour \( i = 1, \dots, p \), sont linéairement indépendants.
			\item[ii)] Il existe un vecteur non nul \( d \in \mathbb{R}^n \setminus \{0\} \) tel que :
			\[
			\forall i \in \{1, \dots, p\}, \quad \langle \nabla h_i(x^*), d \rangle = 0,
			\]
			et
			\[
			\forall j \in I(x^*), \quad \langle \nabla g_j(x^*), d \rangle < 0.
			\]
		\end{enumerate}
	\end{definition}
	
	\textbf{Interprétation détaillée :} La condition i) assure une bonne paramétrisation locale des égalités (par théorème des fonctions implicites). La condition ii) garantit l'existence d'une direction \( d \) tangent aux égalités et \textbf{strictement descendante} pour toutes les contraintes actives, permettant d'"échapper" à la frontière dans l'intérieur relatif de \( U \). Cela évite les points dégénérés comme les cusps ou les arêtes pointues. La QCMF est relativement faible et valable pour des problèmes non convexes.
	
	En pratique, la QCMF est souvent vérifiée dans les problèmes bien posés, et elle implique que dans le théorème de Fritz John, on peut avoir \( \mu_0^* > 0 \). Elle est équivalente à la non-vacuité du cône tangent strict.
	
	\begin{definition}[Qualification des contraintes de Kuhn-Tucker (QCKT) ou Qualification d'Indépendance Linéaire (LICQ)]
		Un point réalisable \( x^* \in U \) vérifie la QCKT (ou LICQ) si les vecteurs \( \nabla h_i(x^*) \) pour \( i = 1, \dots, p \), et \( \nabla g_j(x^*) \) pour \( j \in I(x^*) \), sont linéairement indépendants.
	\end{definition}
	
	La QCKT est une condition d'indépendance linéaire totale sur les contraintes actives (égalités + inégalités actives). Elle est plus facile à vérifier algébriquement (via le rang de la matrice jacobienne active, qui doit être égal au nombre de contraintes actives). C'est une condition plus forte que la QCMF, impliquant cette dernière.
	
	\begin{proposition}[QCKT implique QCMF]
		Tout point réalisable \( x^* \in U \) vérifiant la QCKT vérifie également la QCMF.
	\end{proposition}
	
	\begin{proof}
		Supposons que \( x^* \) vérifie la QCKT. Soit \( r = |I(x^*)| \leq q \) le nombre de contraintes actives. Considérons l'application linéaire :
		\[
		\psi : \mathbb{R}^n \to \mathbb{R}^{p+r}, \quad x \mapsto \left( (\langle \nabla h_i(x^*), x \rangle)_{1 \leq i \leq p}, (\langle \nabla g_j(x^*), x \rangle)_{j \in I(x^*)} \right)^\top.
		\]
		La matrice de \( \psi \) est formée des \( p + r \) vecteurs lignes \( \nabla h_i(x^*) \) et \( \nabla g_j(x^*) \) pour \( j \in I(x^*) \), qui sont linéairement indépendants par QCKT. Ainsi, \( \psi \) est surjective (rang maximal \( p+r \)).
		
		En particulier, il existe un vecteur \( d \in \mathbb{R}^n \setminus \{0\} \) tel que \( \psi(d) = (0, \dots, 0, -1, \dots, -1)^\top \) (zéros pour les \( p \) premières composantes, et -1 pour les \( r \) suivantes). Puisque \( \psi(0) = 0 \), \( d \neq 0 \). Ce \( d \) satisfait exactement :
		- \( \langle \nabla h_i(x^*), d \rangle = 0 \) pour tout \( i \),
		- \( \langle \nabla g_j(x^*), d \rangle = -1 < 0 \) pour tout \( j \in I(x^*) \).
		
		Ainsi, \( x^* \) vérifie la condition ii) de la QCMF, et i) est déjà incluse dans QCKT.
	\end{proof}
	
	\begin{remark}
		La QCKT (LICQ) est suffisante pour la plupart des applications numériques, car elle permet de résoudre le système des conditions KKT via des méthodes comme Newton ou SQP (Sequential Quadratic Programming). Elle assure aussi l'unicité locale des multiplicateurs.
	\end{remark}
	
	\begin{definition}[Condition de Slater]
		Pour un problème convexe (où \( f \) et les \( g_j \) sont convexes, et les \( h_i \) affines), un point réalisable \( x^* \in U \) vérifie la condition de Slater si les vecteurs \( \nabla h_i(x^*) \) sont linéairement indépendants, et s'il existe un point strictement réalisable \( \bar{x} \in \mathbb{R}^n \) tel que \( h(\bar{x}) = 0 \) et \( g_j(\bar{x}) < 0 \) pour tout \( j = 1, \dots, q \) (intérieur relatif non vide pour les inégalités).
	\end{definition}
	
	\textbf{Interprétation :} La condition de Slater est une qualification "globale" pour les problèmes convexes, garantissant la forte dualité (égalité des valeurs primale et duale) et l'existence de multiplicateurs avec \( \mu_0^* = 1 \). Elle est plus forte que la QCMF dans le contexte convexe, mais ne requiert pas d'indépendance des gradients des inégalités actives. Elle est particulièrement utile en programmation linéaire ou convexe quadratique.
	
	\begin{proposition}[Slater implique QCMF dans le cas convexe]
		Si les fonctions sont convexes et la condition de Slater est satisfaite, alors la QCMF l'est aussi localement autour des points optimaux.
	\end{proposition}
	
	\begin{proof}[Esquisse]
		Le point strictement réalisable \( \bar{x} \) permet de construire une direction \( d = \bar{x} - x^* \) (ajustée pour les égalités), telle que \( \langle \nabla g_j(x^*), d \rangle \approx g_j(\bar{x}) - g_j(x^*) < 0 \) par convexité et linéarisation. L'indépendance des \( \nabla h_i \) assure la tangence.
	\end{proof}
	
	\begin{example}[Vérification de la condition de Slater]
		Pour minimiser \( f(x) = \|x\|^2 \) sous \( g(x) = A x - b \leq 0 \) (convexe), si l'intérieur de $ \{x | A x < b\} $ est non vide, Slater est vérifiée. Si les inégalités sont affines et l'ensemble strict est vide (ex. égalités déguisées), elle échoue.
	\end{example}
	
	\begin{definition}[Qualification de Rang Constant (CRCQ)]
		Un point réalisable \( x^* \in U \) vérifie la CRCQ si, dans un voisinage de \( x^* \), le rang de la matrice formée par les gradients des contraintes actives (égalités + inégalités actives en points voisins) est constant.
	\end{definition}
	
	La CRCQ est une condition plus faible que la LICQ, permettant des dépendances locales tant que le rang ne varie pas. Elle est utile pour des algorithmes où l'ensemble actif change, et implique la QCMF sous certaines hypothèses.
	
	\begin{remark}
		Il existe d'autres qualifications comme la Pseudo-Normalité (pour des problèmes non différentiables), la CPLD (Constant Positive Linear Dependence), ou l'ACQ (Abadie Constraint Qualification), qui équivaut à l'égalité entre cône tangent et cône linéarisé. Pour un cours de L3, les QCMF, LICQ/QCKT et Slater suffisent souvent, mais la CRCQ est introduite pour illustrer la variété.
	\end{remark}
	
	\begin{example}[Vérification de QCKT]
		Pour minimiser \( f(x,y) = x^2 + y^2 \) sous \( g(x,y) = x + y - 1 \leq 0 \). À \( x^* = (0.5, 0.5) \), la contrainte est active, \( \nabla g = (1,1) \), indépendant (pas d'égalité). QCKT vérifiée. Si on ajoute une contrainte redondante \( 2x + 2y - 2 \leq 0 \), QCKT échoue mais CRCQ peut tenir si rang constant =1.
	\end{example}
	
	\begin{proposition}[Implications entre qualifications]
		En général : LICQ (QCKT) \(\implies\) MFCQ (QCMF) \(\implies\) CPLD \(\implies\) ACQ. Pour convexe : Slater \(\implies\) MFCQ.
	\end{proposition}
	
	\section{Conditions de Karush-Kuhn-Tucker (KKT)}
	
	\begin{definition}[Lagrangien du problème mixte]
		On appelle Lagrangien du problème $(P)$ la fonction $\mathcal{L}: \R^n \times \R^p \times \R^q \to \R$ définie par:
		$$\mathcal{L}(x, \lambda, \mu) = f(x) + \sum_{i=1}^{p}\lambda_{i}h_{i}(x) + \sum_{j=1}^{q}\mu_{j}g_{j}(x)$$
		où $\lambda$ est le vecteur des multiplicateurs pour les contraintes d'égalité et $\mu$ pour les inégalités.
	\end{definition}
	
	\begin{theorem}[Conditions de Karush-Kuhn-Tucker (KKT)]\label{Theo4.3.2}
		Soit $x^*$ un minimum local du problème $(P)$ qui vérifie une condition de qualification QCMF. Alors il existe des multiplicateurs $\lambda^* \in \R^p$ et $\mu^* \in \R^q$ tels que le triplet $(x^*, \lambda^*, \mu^*)$ vérifie les 4 conditions suivantes :
		\begin{enumerate}
			\item \textbf{Stationnarité :} Le gradient du Lagrangien par rapport à $x$ est nul.
			$$\grad_x \mathcal{L}(x^*, \lambda^*, \mu^*) = \nabla f(x^*) + \sum_{i=1}^p \lambda_i^* \nabla h_i(x^*) + \sum_{j=1}^q \mu_j^* \nabla g_j(x^*) = 0$$
			\item \textbf{Admissibilité primale :} Le point $x^*$ respecte toutes les contraintes.
			$$h(x^*) = 0 \quad \text{et} \quad g(x^*) \le 0$$
			\item \textbf{Admissibilité duale :} Les multiplicateurs des contraintes d'inégalité sont positifs ou nuls.
			$$\mu^* \ge 0$$
			\item \textbf{Complémentarité :} Pour chaque contrainte d'inégalité, soit la contrainte est active, soit son multiplicateur est nul (ou les deux).
			$$\mu_j^* g_j(x^*) = 0 \quad \text{pour tout } j=1,\dots,q$$
		\end{enumerate}
	\end{theorem}
	
	\begin{proof}
		Il suffit de montrer que l’on a nécessairement \(\mu^*_0\) non nul dans les coefficients donnés par le Théorème \ref{theo4.1.1} (et alors on divise par \(\mu^*_0 > 0\) pour obtenir le résultat). Supposons par l’absurde que \(\mu^*_0 = 0\). On a alors \(\lambda^* = (\lambda^*_1, \ldots, \lambda^*_p)^\top \in \mathbb{R}^p\) et \(\mu^* = (\mu^*_1, \ldots, \mu^*_q)^\top \in \mathbb{R}^q_+\) non tous deux nuls tels que
		
		\[
		\sum_{i=1}^p \lambda^*_i \nabla h_i(x^*) + \sum_{j=1}^q \mu^*_j \nabla g_j(x^*) = 0
		\]
		
		et pour tout \(j \in \{1, \ldots, q\}\), \(\mu^*_j g_j(x^*) = 0\).
		
		Par définition, pour tout \(j \notin I(x^*)\), \(\mu^*_j = 0\). On a ensuite deux cas. Si pour tous \(j \in I(x^*)\), \(\mu^*_j = 0\), alors on obtient \(\lambda^* = (\lambda^*_1, \ldots, \lambda^*_p)^\top \neq 0\) tel que
		
		\[
		\sum_{i=1}^p \lambda^*_i \nabla h_i(x^*) = 0,
		\]
		
		ce qui contredit le point i) de QCMF. Si il existe \(j_0 \in I(x^*)\) tel que \(\mu^*_{j_0} > 0\), alors, en prenant un vecteur \(d\) vérifiant le point ii) de QCMF (Définition \ref{Def4.2.2}),
		
		\[
		0 = \sum_{i=1}^p \lambda^*_i \langle \nabla h_i(x^*), d \rangle + \sum_{j=1}^q \mu^*_j \langle \nabla g_j(x^*), d \rangle \leq \mu^*_{j_0} \langle \nabla g_{j_0}(x^*), d \rangle < 0.
		\]
		
		On obtient encore une contradiction, ce qui conclut la preuve par l’absurde.
	\end{proof}
	
	\begin{remark}
		La condition de complémentarité est la plus subtile. Elle signifie que si une contrainte est inactive ($g_j(x^*) < 0$), son multiplicateur $\mu_j^*$ doit être nul. En d'autres termes, seules les contraintes actives peuvent avoir un multiplicateur non nul et influencer la condition de stationnarité.
	\end{remark}
	
	\begin{remark}
		Le Théorème \ref{Theo4.3.2} peut être énoncé pour toute condition de qualification des contraintes (QC).
	\end{remark}
	\begin{example}[Minimisation avec contrainte]
	% Énoncé du problème
	On cherche à minimiser la fonction :
	\[
	f(x, y) = (x - 2)^2 + (y - 1)^2
	\]
	sous la contrainte :
	\[
	g_1(x, y) = x + y - 2 \leq 0.
	\]
	
	% Introduction aux conditions KKT
	Nous utilisons les conditions de Karush-Kuhn-Tucker (KKT) pour résoudre ce problème d'optimisation avec une contrainte d'inégalité. Les conditions KKT pour un problème de la forme \(\displaystyle\min f(x, y)\) sous \(g_1(x, y) \leq 0\) sont :
	\begin{enumerate}
		\item Stationnarité : \(\nabla f(x, y) + \lambda \nabla g_1(x, y) = 0\), où \(\lambda\) est le multiplicateur de Lagrange.
		\item Faisabilité : \(g_1(x, y) \leq 0\).
		\item Complémentarité : \(\lambda g_1(x, y) = 0\).
		\item Non-négativité du multiplicateur : \(\lambda \geq 0\).
	\end{enumerate}

	% Vérification des contraintes de qualification
	\subsection*{Vérification des contraintes de qualification}
	Nous vérifions ici deux conditions courantes : la qualification de Slater et la qualification d'indépendance linéaire des contraintes (LICQ).
	
	\begin{itemize}
		\item \textbf{Qualification de Slater (pour les contraintes convexes)} : L'ensemble réalisable est défini par \(g_1(x, y) \leq 0\), où \(g_1\) est affine (donc convexe). L'ensemble intérieur est non vide, car il existe un point strictement réalisable, par exemple \((0, 0)\) où \(g_1(0, 0) = -2 < 0\). Ainsi, la condition de Slater est satisfaite, ce qui implique que les conditions KKT sont nécessaires et suffisantes pour l'optimalité (dans le cas convexe).
		
		\item \textbf{LICQ (Linear Independence Constraint Qualification)} : Au point optimal potentiel, les gradients des contraintes actives doivent être linéairement indépendants. Ici, il n'y a qu'une seule contrainte \(g_1\), et son gradient est \(\nabla g_1(x, y) = (1, 1)^\top\), qui est non nul. Ainsi, le seul gradient est linéairement indépendant, et LICQ est satisfaite.
	\end{itemize}
	
	Puisque les contraintes sont linéaires, ces conditions sont toujours vérifiées pour ce type de problème. Les conditions KKT sont donc applicables, et tout point satisfaisant KKT est un candidat optimal.
	
	% Étape 1 : Calcul des gradients
	Calculons les gradients de \(f\) et \(g_1\) :
	\[
	\nabla f(x, y) = \begin{pmatrix} \frac{\partial f}{\partial x} \\ \frac{\partial f}{\partial y} \end{pmatrix} = \begin{pmatrix} 2(x - 2) \\ 2(y - 1) \end{pmatrix},
	\]
	\[
	\nabla g_1(x, y) = \begin{pmatrix} \frac{\partial g_1}{\partial x} \\ \frac{\partial g_1}{\partial y} \end{pmatrix} = \begin{pmatrix} 1 \\ 1 \end{pmatrix}.
	\]
	
	% Étape 2 : Application des conditions KKT
	La condition de stationnarité donne :
	\[
	\nabla f(x, y) + \lambda \nabla g_1(x, y) = \begin{pmatrix} 2(x - 2) \\ 2(y - 1) \end{pmatrix} + \lambda \begin{pmatrix} 1 \\ 1 \end{pmatrix} = \begin{pmatrix} 0 \\ 0 \end{pmatrix}.
	\]
	Cela se traduit par le système :
	\[
	\begin{cases}
	2(x - 2) + \lambda = 0, \\
	2(y - 1) + \lambda = 0.
	\end{cases}
	\]
	De plus, nous avons :
	\begin{itemize}
		\item Faisabilité : \(x + y - 2 \leq 0\),
		\item Complémentarité : \(\lambda (x + y - 2) = 0\),
		\item Non-négativité : \(\lambda \geq 0\).
	\end{itemize}
	
	% Étape 3 : Résolution des cas
	Analysons les deux cas possibles pour la condition de complémentarité.
	
	\subsection*{Cas 1 : \(\lambda = 0\)}
	Si \(\lambda = 0\), alors :
	\[
	2(x - 2) = 0 \implies x = 2, \quad 2(y - 1) = 0 \implies y = 1.
	\]
	Le point candidat est \((x, y) = (2, 1)\). Vérifions la faisabilité :
	\[
	g_1(2, 1) = 2 + 1 - 2 = 1 > 0.
	\]
	Ce point ne satisfait pas la contrainte \(x + y - 2 \leq 0\), donc il n'est pas admissible.
	
	\subsection*{Cas 2 : \(\lambda \neq 0\), donc \(x + y - 2 = 0\)}
	Si \(x + y = 2\), alors \(y = 2 - x\). De la condition de stationnarité :
	\[
	2(x - 2) + \lambda = 0 \implies \lambda = -2(x - 2) = 2(2 - x),
	\]
	\[
	2(y - 1) + \lambda = 0.
	\]
	Substituons \(\lambda = 2(2 - x)\) dans la deuxième équation :
	\[
	2(y - 1) + 2(2 - x) = 0 \implies y - 1 + 2 - x = 0 \implies y - x + 1 = 0.
	\]
	Puisque \(y = 2 - x\), substituons :
	\[
	(2 - x) - x + 1 = 2 - 2x + 1 = 3 - 2x = 0 \implies x = \frac{3}{2}.
	\]
	Alors, \(y = 2 - \frac{3}{2} = \frac{1}{2}\). Calculons \(\lambda\) :
	\[
	\lambda = 2\left(2 - \frac{3}{2}\right) = 2 \cdot \frac{1}{2} = 1.
	\]
	Vérifions les conditions KKT :
	\begin{itemize}
		\item Faisabilité : \(g_1\left(\frac{3}{2}, \frac{1}{2}\right) = \frac{3}{2} + \frac{1}{2} - 2 = 0 \leq 0\).
		\item Complémentarité : \(\lambda (x + y - 2) = 1 \cdot 0 = 0\).
		\item Non-négativité : \(\lambda = 1 \geq 0\).
	\end{itemize}
	Le point \(\left(\frac{3}{2}, \frac{1}{2}\right)\) satisfait toutes les conditions KKT.
	
	% Étape 4 : Vérification de la convexité
	Pour confirmer que \(\left(\frac{3}{2}, \frac{1}{2}\right)\) est un minimum global, vérifions la convexité de \(f\) et la nature de la contrainte. Le hessien de \(f\) est :
	\[
	H_f = \begin{pmatrix}
	\frac{\partial^2 f}{\partial x^2} & \frac{\partial^2 f}{\partial x \partial y} \\
	\frac{\partial^2 f}{\partial y \partial x} & \frac{\partial^2 f}{\partial y^2}
	\end{pmatrix} = \begin{pmatrix}
	2 & 0 \\
	0 & 2
	\end{pmatrix}.
	\]
	Les valeurs propres sont \(2, 2 > 0\), donc \(H_f\) est défini positif, et \(f\) est strictement convexe. La contrainte \(g_1(x, y) = x + y - 2 \leq 0\) est linéaire, donc convexe. Puisque \(f\) est strictement convexe et l'ensemble réalisable est convexe, le problème admet au plus un minimum global. Comme \(\left(\frac{3}{2}, \frac{1}{2}\right)\) satisfait les conditions KKT, c'est le minimum global.
	
	% Étape 5 : Évaluation de la fonction
	Évaluons \(f\) au point critique :
	\[
	f\left(\frac{3}{2}, \frac{1}{2}\right) = \left(\frac{3}{2} - 2\right)^2 + \left(\frac{1}{2} - 1\right)^2 = \left(-\frac{1}{2}\right)^2 + \left(-\frac{1}{2}\right)^2 = \frac{1}{4} + \frac{1}{4} = \frac{1}{2}.
	\]
	
	% Étape 6 : Illustration graphique
	\textbf{Illustration géométrique.}
	
	La figure ci-dessous montre la région réalisable définie par \(x + y \leq 2\), la frontière de la contrainte \(x + y = 2\), quelques courbes de niveau de \(f(x, y) = (x - 2)^2 + (y - 1)^2\), et le point optimal \(\left(\frac{3}{2}, \frac{1}{2}\right)\). Les courbes de niveau de \(f\) sont des cercles centrés en \((2, 1)\), et le point optimal est le point de tangence entre la courbe de niveau de \(f\) et la droite \(x + y = 2\).
	
		\begin{figure}[H]
			\centering
			\begin{tikzpicture}[scale=2]
			% Axes
			\draw[->] (-0.25, 0) -- (3, 0) node[right] {$x$};
			\draw[->] (0, -0.25) -- (0, 3) node[above] {$y$};
			
			% Étiquettes des axes
			\foreach \x in {1, 2} \draw (\x, 0.1) -- (\x, -0.1) node[below] {\x};
			\foreach \y in {1, 2} \draw (0.1, \y) -- (-0.1, \y) node[left] {\y};
			
			% Région réalisable (x + y <= 2)
			\fill[blue!20, opacity=0.3] (0, 0) -- (2, 0) -- (0, 2) -- cycle;
			\draw[thick] (0, 2) -- (2, 0) node[midway, above, sloped] {$x + y = 2$};
			
			% Courbes de niveau de f(x, y)
			\draw[red!50, dashed] (2, 1) circle (0.5); % f = 0.25
			\draw[red!50, dashed] (2, 1) circle (0.707); % f = 0.5
			\draw[red!50, dashed] (2, 1) circle (1); % f = 1
			
			% Point optimal (3/2, 1/2)
			\fill[red] (1.5, 0.5) circle (1.5pt) node[above right] {$ \left( \frac{3}{2}, \frac{1}{2} \right) $};
			
			% Centre des courbes de niveau (2, 1)
			\fill[cyan] (2, 1) circle (1pt) node[above right] {$(2, 1)$};
			
			% Légende
			\node at (3, 2.5) [align=left, scale=0.8] {
				Région réalisable : bleu hachuré\\
				Courbes de niveau de $f$ : cercles rouges\\
				Point optimal : $\left( \frac{3}{2}, \frac{1}{2} \right)$
			};
			\end{tikzpicture}
			\caption{Représentation graphique du problème d'optimisation.}
			\label{fig:optimisation}
		\end{figure}
		
	\end{example}
	
	\section{Cas des problèmes convexes}
	La puissance des conditions KKT est décuplée dans le cadre de l'optimisation convexe.
	
	\begin{definition}[Problème d'optimisation convexe]
		Le problème $(P)$ est dit convexe si la fonction objectif $f$ est convexe, les fonctions de contrainte d'inégalité $g_j$ sont convexes, et les fonctions de contrainte d'égalité $h_i$ sont affines.
	\end{definition}
	
	\begin{theorem}[KKT pour les problèmes convexes]	
		Si $ (P) $ est un problème de programmation convexe, et si $ (QC)_{x^{*}}$ est vérifiée, alors les deux conditions KKT en $ x\in D $ sont suffisantes pour que $ x^{*} $ soit un minimum global de $ f $.
	\end{theorem}
	\begin{remark}
		\textit{C'est un résultat majeur : dans le cas convexe, les conditions KKT ne sont plus seulement nécessaires, elles deviennent suffisantes ! Trouver un point KKT, c'est trouver LA solution globale.}
	\end{remark}
	
	\begin{example}
		Considérons le problème
		\[\begin{aligned}
		&\displaystyle\min_{x=(x_1,x_2)\in\mathbb{R}^2}\; f(x)=(x_1-1)^2+(x_2-2)^2,\\[4pt]
		&\text{s.c.}\quad g_1(x)=-x_1\le0,\quad g_2(x)=-x_2\le0,\quad g_3(x)=x_1+x_2-2\le0.
		\end{aligned}\]
		Le domaine réalisable est le triangle $\{x_1\ge0,\ x_2\ge0,\ x_1+x_2\le2\}$.
		
		\textbf{Calcul des conditions KKT}\\
		Les gradients des contraintes sont
		\[\nabla g_1=(-1,0),\quad \nabla g_2=(0,-1),\quad \nabla g_3=(1,1).
		\]
		qui sont linéairement indépendants pour tout $ x $ donc tout point réalisable vérifie les QCKT (Ce qui prouve l'existence des multiplicateurs de KKT).
		
		Soit $\lambda=(\lambda_1,\lambda_2,\lambda_3)\ge0$ les multiplicateurs KKT. La stationnarit\'e s'\'ecrit
		\[
		\nabla f(x)+\sum_{i=1}^3\lambda_i\nabla g_i(x)=0,
		\]
		c'est-\`a-dire
		\[\begin{pmatrix}2(x_1-1)\\[4pt]2(x_2-2)\end{pmatrix}
		+\lambda_1\begin{pmatrix}-1\\0\end{pmatrix}+\lambda_2\begin{pmatrix}0\\-1\end{pmatrix}+\lambda_3\begin{pmatrix}1\\1\end{pmatrix}=\begin{pmatrix}0\\0\end{pmatrix}.
		\]
		De plus $\lambda_i g_i(x)=0$ (compl\'ementarit\'e) et $g_i(x)\le0$, $\lambda_i\ge0$.
		
		\textbf{Étude des cas}\\
		
		\textbf{Cas A : toutes les contraintes inactives}\\
		Si $x_1>0$ et $x_2>0$ et $x_1+x_2<2$, alors $\lambda_1=\lambda_2=\lambda_3=0$ et
		\[2(x_1-1)=0,\quad 2(x_2-2)=0\]
		d'o\`u $(x_1,x_2)=(1,2)$, point non admissible car $1+2>2$. Ce cas est rejet\'e.
		
		\textbf{Cas B : $g_3$ active, $g_1,g_2$ inactives}\\
		Supposons $x_1>0$, $x_2>0$ et $x_1+x_2=2$. Alors $\lambda_1=\lambda_2=0$, $\lambda_3\ge0$ et
		\[\begin{cases}
		2(x_1-1)+\lambda_3=0,\\[4pt]
		2(x_2-2)+\lambda_3=0.
		\end{cases}\]
		En substituant $x_2=2-x_1$ et en éliminant $\lambda_3$ on obtient l'\'equation pour $x_1$ :
		\[\frac{d}{dx_1}\bigl((x_1-1)^2+(2-x_1-2)^2\bigr)=4x_1-2=0\Longrightarrow x_1=\tfrac12,
		\]
		d'o\`u $x^\star=(\tfrac12,\tfrac32)$. Le multiplicateur $\lambda_3$ s'en d\'eduit :
		\[2(\tfrac12-1)+\lambda_3=0\Longrightarrow \lambda_3=1.
		\]
		On v\'erifie la compl\'ementarit\'e et la faisabilit\'e : $x^\star$ est admissible.
		
		\textbf{Cas C : une contrainte bordée}\\
		On peut aussi examiner les ar\^etes $x_1=0$ ou $x_2=0$. En v\'erifiant les valeurs de $f$ sur ces ar\^etes on constate qu'elles sont sup\'erieures \`a $f(x^\star)$. (D\'etail laiss\'e en exercice.)
		
		%		\section*{Conclusion}
		La solution du problème est
		\[x^\star=\Bigl(\tfrac12,\tfrac32\Bigr),\qquad f(x^\star)=\tfrac12,
		\]
		avec multiplicateurs KKT $\lambda=(0,0,1)$. Par la convexit\'e stricte de $f$ et la convexit\'e du domaine, il s'agit de l'unique minimum global.
		
		%	\end{document}
	\end{example}
	
	
	\begin{example}[Inégalité avec conditions KKT]
		Montrons que pour tout \(x_1\ge0,\;x_2\ge0\) on a
		\[
		x_1^2 + x_2^2 \le 4\,e^{\,x_1+x_2-2},
		\]
		équivalente à
		\[
		(x_1^2 + x_2^2)e^{-x_1-x_2} \le 4e^{-2}.
		\]
		
		Posons
		\[
		f(x)= - (x_1^2 + x_2^2)e^{-x_1-x_2},\qquad x=(x_1,x_2),
		\]
		et considérons le problème contraint
		\[
		\displaystyle\min_{x\in\mathbb{R}^2} f(x)\quad\text{s.c.}\quad x_1\ge0,\ x_2\ge0,
		\]
		dont la valeur minimale attendue est \(-4e^{-2}\).
		
		Les contraintes s’écrivent 
		\[
		g_1(x)=-x_1\le0,\qquad g_2(x)=-x_2\le0.
		\]
		Les fonctions \(f,g_1,g_2\) sont \(C^1\) sur \(\mathbb{R}^2\) et
		\[
		\nabla g_1(x)=\begin{pmatrix}-1\\0\end{pmatrix},\qquad
		\nabla g_2(x)=\begin{pmatrix}0\\-1\end{pmatrix},
		\]
		qui sont linéairement indépendants: la qualification de contraintes (LICQ / QCKT) est satisfaite en tout point réalisable. Les conditions KKT s’appliquent donc.
		
		Calculons le gradient de \(f\). En notant \(S=x_1^2+x_2^2\) et \(t=x_1+x_2\), on a
		\[
		f(x) = -S e^{-t},\quad
		\nabla f(x)=
		\begin{pmatrix}
		(S-2x_1)e^{-t}\\[2mm]
		(S-2x_2)e^{-t}
		\end{pmatrix}.
		\]
		
		D’après KKT, pour tout point optimal \(x\) il existe \(\mu=(\mu_1,\mu_2)\ge0\) tels que
		\[
		\nabla f(x)+\mu_1\nabla g_1(x)+\mu_2\nabla g_2(x)=0,
		\qquad
		\mu_1 g_1(x)=\mu_2 g_2(x)=0.
		\]
		
		On obtient le système
		\[
		\begin{cases}
		(S-2x_1)e^{-t}-\mu_1=0,\\[1mm]
		(S-2x_2)e^{-t}-\mu_2=0,\\[1mm]
		\mu_1\ge0,\ \mu_2\ge0,\\[1mm]
		\mu_1 x_1=\mu_2 x_2=0.
		\end{cases}
		\]
		
		\textbf{Cas A — contraintes inactives:} \(x_1>0,\ x_2>0\) donc \(\mu_1=\mu_2=0\). On cherche les points critiques de \(f\):
		\[
		\begin{cases}
		S-2x_1=0,\\
		S-2x_2=0
		\end{cases} \implies x_1=x_2.
		\]
		En posant \(x_1=x_2=x\), on obtient \(2x^2-2x=0 \implies x=0 \text{ ou } x=1\). Seul \((1,1)\) donne \(f(1,1)=-2e^{-2}\).
		
		\textbf{Cas B — une contrainte active:} par symétrie, supposons \(x_1=0<x_2\). Alors \(\mu_1\ge0,\ \mu_2=0\) et le système devient
		\[
		\begin{cases}
		x_2^2 e^{-x_2}-\mu_1=0,\\
		(x_2^2-2x_2)e^{-x_2}=0
		\end{cases} \implies x_2=2, \ \mu_1=4e^{-2}.
		\]
		Par symétrie, on obtient également le point \((2,0)\) avec \(f(2,0)=-4e^{-2}\).
		
		\textbf{Cas C — deux contraintes actives:} \(x_1=x_2=0\) donne \(\mu_1=\mu_2=0\), soit \(f(0,0)=0\), non minimal.
		
		Il reste à justifier qu’il s’agit bien d’un minimum global de $f$. Ici $f$ n’est pas coercive, mais un
		argument similaire suffit : on montre qu’à l’infini $f$ tend vers 0. 
		
		Pour tout $x_1, x_2 \geqslant 0$, par inégalité
		triangulaire
		$$ \|x\| = \|(x_1, 0) + (0, x_2)\| \leqslant \|(x_1, 0)\| + \|(0, x_2)\| = |x_1| + |x_2| = x_1 + x_2 ,$$
		donc par croissance de $t \mapsto -e^{-t}$,
		$$ \|x\| \leqslant (x_1 + x_2) \implies -\|x\| \geqslant -(x_1 + x_2) \implies e^{-\|x\|} \geqslant e^{-(x_1+x_2)} \implies \|x\|^2e^{-\|x\|} \geqslant \|x\|^2e^{-(x_1+x_2)} $$
		$$ -\|x\|^2e^{-\|x\|} \leqslant -\|x\|^2e^{-(x_1+x_2)} = -(x_1^2+x_2^2)e^{-(x_1+x_2)} = f(x) $$$$\implies 0 \xleftarrow[\|x\| \to +\infty]{} -\|x\|^2e^{-\|x\|} \leqslant f(x) \leqslant 0 \xrightarrow[\|x\| \to +\infty]{} 0 $$
		Donc $f(x) \xrightarrow[\|x\| \to +\infty]{} 0$. Ainsi, sur $[0, +\infty[^2$, $f$ tend vers 0 à l’infini. 
		
		Il existe donc un voisinage $B(0, R) \cap
		[0, +\infty[^2$ tel que $f$ est supérieure à $\frac{1}{2}f(0, 2) < 0$ en dehors de ce compact. Donc la valeur $f(0, 2)$ est bien
		un minimum global de $f$ sur $[0, +\infty[^2$.
		
		Ainsi les minima globaux sont atteints en \((0,2)\) et \((2,0)\), et la valeur minimale du problème est \(-4e^{-2}\).
		
		On conclut que pour tout \(x_1,x_2\ge0\),
		\[
		(x_1^2+x_2^2)e^{-x_1-x_2}\le 4e^{-2},
		\]
		l’égalité étant réalisée uniquement en \((0,2)\) et \((2,0)\).
	\end{example}
	
	
	
	% Exercice 2 : Problème KKT simple (x>=1, y>=2)
	\begin{exercise}[Un problème KKT simple]
		\textbf{Énoncé :} Résoudre
		\[\displaystyle\min_{(x,y)\in\mathbb{R}^2} f(x,y)=x^2+y^2\quad\text{s.c.}\quad x\ge1,\ \ y\ge2.\]
		
	\end{exercise}
	
	% Exercice 3 : Problème quadratique sur disque et demi-plan
	\begin{exercise}[Problème KKT quadratique]
		\textbf{Énoncé :} 
		
		Minimiser:
		\[
		f(x, y) = (x - 3)^2 + (y - 5)^2
		\]
		sous les contraintes :
		\[
		g_1(x, y) = x^2 + y^2 - 20 \leq 0, \quad g_2(x, y) = x + y - 4 \leq 0.
		\]
	\end{exercise}
	
	% Exercice 4 : égalité et inégalité
	\begin{exercise}[Égalité et inégalité]
		\textbf{Énoncé :} Maximiser $f(x,y)=y$ sous
		\[x^2+y^2=1,\qquad x+y\ge0.
		\]		
	\end{exercise}
	
	%    \end{document}
	%===============================================================
	
	
	\chapter{Généralités sur la programmation linéaire}
	\label{chap:pl_generalites}
	
	\section{Introduction}
	
	En mathématiques appliquées et en recherche opérationnelle, les problèmes de \textbf{programmation linéaire (PL)} sont des problèmes d'optimisation dans lesquels la \textbf{fonction objectif} et les \textbf{contraintes} sont toutes deux linéaires.  
	Autrement dit, la programmation linéaire vise à \textit{trouver la meilleure décision possible} (au sens de maximiser un profit ou minimiser un coût) parmi un ensemble de solutions réalisables, défini par des contraintes linéaires.  
	\section{Les conditions de formulation}
	
	La programmation linéaire, comme étant un modèle, admet certaines hypothèses (ou conditions) que le décideur doit vérifier avant de pouvoir l'utiliser pour représenter correctement son problème réel. Ces hypothèses garantissent que la modélisation est cohérente et que les outils mathématiques de la PL s'appliquent.  
	
	Ces hypothèses sont :
	\begin{itemize}
		\item[$\diamond $] les variables de décision du problème sont positives ;
		\item[$\diamond $] le critère de sélection de la meilleure décision est décrit par une fonction linéaire de ces variables;
		\item[$\diamond $] les restrictions relatives aux variables de décision (exemple : limitation des ressources) peuvent être exprimées par un ensemble d'équations linéaires. Ces équations forment l'ensemble des contraintes;
		\item[$\diamond $]  les paramètres du problème en dehors des variables de décisions ont une valeur connue avec certitude.
	\end{itemize}
	\section{Les étapes de formulation de PL}
	Généralement, il y a trois étapes pour pouvoir construire le modèle d'un programme linéaire :
	\begin{enumerate}
		\item identifier les variables du problème à valeur non connue (variables de décision) et leur nature ou type (discret, continu, binaire)
		\item identifier les restrictions (les contraintes) du problème et les exprimer par un système d'équation linéaire ou un système d'inéquation linéaire ;
		\item  identifier l'objectif ou le critère de sélection qui se traduit par la fonction économique ou la fonction "objectif". Spécifier si la fonction objectif est à maximiser ou à minimiser.
	\end{enumerate}
	
	\section{Formes générales d'un programme linéaire}
	
	La forme la plus générale d'un problème de PL est la suivante :
	\begin{equation*}\label{key}
	\begin{array}{lc}
	opt\:\: z =\displaystyle\sum_{j=1}^{n}c_{j}x_{j}\\
	\\
	sc \left\lbrace \begin{array}{llc}
	\displaystyle\sum_{j=1}^{n}a_{ij}x_{j}\leq (\geq) b_{i}&i=\overline{1,p}& (2a)\\
	\displaystyle\sum_{j=1}^{n}a_{ij}x_{j}= b_{i}&i=\overline{1,q}& (2b)\\
	x_{j}\geq 0&j=\overline{1,q}& (3)\\
	x_{j}\in \mathbb{R}& j\in q+1,\ldots,n& (4).
	\end{array}\right.
	\end{array}
	\end{equation*}
	Dans cette formulation générale :
	\begin{itemize}
		\item[$\diamond $] $ opt $ pour optimisation peut-être une maximisation ou une minimisation ;
		\item[$\diamond $] les $ x_{j} $ représentent les variables de décision au nombre de $ n $ ;
		\item[$\diamond $] $ z $ est la fonction objectif ou économique ;
		\item[$\diamond $] $ a_{ij} $, $ c_{j} $ et $ b_{i} $ sont des constantes, c'est-à-dire des données du problème.
		\item[$\diamond $] les équations $ (2a) $ et $ (2b) $ sont les contraintes réelles du problème. Avec $ (2a) $ les contraintes d'inégalités et $ (2b) $ les contraintes d'égalités,
		\item[$\diamond $] $ (3) $ sont des contraintes de non-négativité ou de signes,
		\item[$\diamond $] $ (4) $ implique que les variables sont sans restriction de signes.
	\end{itemize}
	
	\subsection{Forme canonique mixte}
	
	Il s'agit d'un problème de programmation linéaire (encore appelé programme linéaire) écrit sous la forme suivante.
	
	\[
	\displaystyle\max_{(x_1, \ldots, x_n)} f(x_1, \ldots, x_n) = c_1 x_1 + c_2 x_2 + \cdots + c_n x_n = \sum_{j=1}^n c_j x_j
	\]
	sous contraintes :
	\[
	\begin{cases}
	\text{contraintes d'inégalités : } & \forall i \in I_1, \quad \sum_{j=1}^n a_{ij} x_j = a_{i1} x_1 + a_{i2} x_2 + \cdots + a_{in} x_n \leq b_i \\
	\text{contraintes d'égalités : } & \forall i \in I_2, \quad \sum_{j=1}^n a_{ij} x_j = b_i \\
	\text{contraintes de signes : } & \forall j \in J_1, \quad x_j \geq 0 \\
	& \forall j \in J_2, \quad x_j \text{ de signe quelconque}
	\end{cases}
	\]
	
	L'ensemble $I=I_{1} \cup I_{2}$ est l'ensemble des indices de contraintes avec $\operatorname{card}(I)=m$. Autrement dit, il y a $m$ contraintes. L'ensemble $J=J_{1} \cup J_{2}$ est l'ensemble des indices des variables avec $\operatorname{card}(J)=n$. Il y a $n$ variables.
	
	\subsection{Forme canonique pure}
	
	Sous cette forme, il n'y a pas de contraintes d'égalité c'est-à-dire $I_{2}=\varnothing$ et $J_{2}=\varnothing$.\\
	On note
	
	$$
	\begin{aligned}
	x & =(x_{1}, \ldots, x_{n})^{t} \in \mathbb{R}^{n} \\
	c & =(c_{1}, \ldots, c_{n})^{t} \in \mathbb{R}^{n} \\
	b & =(b_{1}, \ldots, b_{m})^{t} \in \mathbb{R}^{m}
	\end{aligned}
	$$
	
	et la matrice $A$ de taille $m \times n$ :
	
	$$
	A=\left(\begin{array}{cccc}
	a_{11} & a_{12} & \cdots & a_{1 n} \\
	\vdots & & & \vdots \\
	a_{m 1} & a_{m 2} & \cdots & a_{m n}
	\end{array}\right)
	$$
	
	Un programme linéaire (PL) est dit sous forme canonique pure s'il s'écrit :
	
	\begin{align*}
	& \displaystyle\max_{x}\left[f\left(x_{1}, \ldots, x_{n}\right)=c^{t} x=c_{1} x_{1}+c_{2} x_{2}+c_{3} x_{3}+\ldots +c_{n} x_{n}\right] \\
	& \left\{\begin{aligned}
	\text { s.t. } & A x \leq b \\
	& x \geq 0
	\end{aligned}\right. \tag{2.2}
	\end{align*}
	
	\begin{remarque}
		Il est toujours possible de ramener :
		
		\begin{itemize}
			\item[$\diamond $] L'optimisation par maximisation équivaut à une minimisation, et vice versa. En effet, maximiser une fonction \( z \) revient à minimiser \(-z\), et minimiser \( z \) revient à maximiser \(-z\) :
			\begin{equation}
			\displaystyle\max(z) = -\displaystyle\min(-z) \quad \text{et} \quad \displaystyle\min(z) = -\displaystyle\max(-z).
			\end{equation}
			\item[$\diamond $] toutes les inégalités à des inégalités de même type en multipliant les deux membres par \(-1\);
			\item[$\diamond $] toutes les variables à être non-négatives.
			si $ x_{j}\in\mathbb{R} $, on remplace $ x_{j} $ par $ x_{j}^{+}-x_{j}^{-} $ avec $ x_{j}^{+}=\displaystyle\max(x_{j},0)\geq0 $ et $ x_{j}^{-}=\displaystyle\max(-x_{j},0)\leq0 $,
			
			\begin{exemple}
				Considérons le problème :
				\[
				\text{Minimiser } z = x_1 - 2x_2
				\]
				sous contraintes :
				\[
				\begin{aligned}
				-x_1 + 4x_2 &\geq 5 \\
				2x_1 + 3x_2 &= 4 \\
				x_1 \geq 0, \, x_2 \text{ sans restriction de signe}
				\end{aligned}
				\]
				Transformons en maximisation :
				\[
				\text{Maximiser } (-z) = -x_1 + 2x_2
				\]
				
				Décomposons $x_2 = x_2^+ - x_2^-$, avec $x_2^+, x_2^- \geq 0$. 
				
				Le programme devient :
				\[
				\text{Maximiser } (-z) = -x_1 + 2x_2^+ - 2x_2^-
				\]
				sous contraintes :
				\[
				\begin{aligned}
				-x_1 + 4x_2^+ - 4x_2^- \geq 5 \\
				2x_1 + 3x_2^+ - 3x_2^- \geq 4 \\
				x_1, x_2^+, x_2^-\geq 0
				\end{aligned}
				\]
			\end{exemple}
			\item[$\diamond $] les contraintes d'égalités en contraintes d'inégalités sous la forme
			\begin{equation}\label{key}
			\sum_{j=1}^{n}a_{ij}x_{j}=b_{i}\Leftrightarrow\left\lbrace \begin{array}{cc}
			\displaystyle\sum_{j=1}^{n}a_{ij}x_{j}\geq b_{i}&\\
			&i=\overline{1,p}\\
			\displaystyle\sum_{j=1}^{n}a_{ij}x_{j}\leq b_{i}&
			\end{array}\right. 
			\end{equation}
		\end{itemize}
	\end{remarque}
	
	\subsection{Forme standard}
	
	Sous cette forme, $I_{1}=\varnothing$ et $J_{2}=\varnothing$. Un programme linéaire (PL) est dit sous forme standard s'il s'écrit :
	
	\begin{align*}
	& \displaystyle\max_{x}\Big[f\left(x_{1}, \ldots, x_{n}\right)=c^{t} x=c_{1} x_{1}+c_{2} x_{2}+c_{3} x_{3}+\ldots +c_{n} x_{n}\Big] \\
	& \left\{\begin{aligned}
	\text { s.t. } & A x=b \\
	& x \geq 0
	\end{aligned}\right. \tag{2.3}
	\end{align*}
	
	
	On dit de plus que le PL est sous forme standard simpliciale si $A$ de taille $m \times n$ avec $m \leq n$, se décompose en :
	
	$$
	A=\left(I_{m} \mid H\right)
	$$
	
	où $I_{m}$ désigne la matrice identité de taille $m \times m$ et $H$ est une matrice de taille $m \times(n-m)$.
	\section{Exemple de formulation}
	%	\subsection{Exemple 1}
	\begin{exemple}
		Un agriculteur veut allouer $ 150 $ hectares de surface irrigable entre culture de tomates et celles de piments. Il dispose de $ 480 $ heures de main d’œuvre et de $ 440 \:m^{3}$ d’eau. Un hectare de tomates demande 1 heure de main d’œuvre, $ 4 \:m^{3}$ d’eau et donne un bénéfice net de $ 1000000 $ $ FCFA $. Un hectare de piments demande $ 4 $ heures de main d’œuvre, $ 2 \:m^{3}$ d’eau et donne un bénéfice net de $ 2000000 $ $ FCFA $.
		
		Le bureau du périmètre irrigué veut protéger le prix des tomates et ne lui permet pas de cultiver plus de $ 90 $ hectares de tomates. Quelle est la meilleure allocation de ses ressources ?
		\textbf{Formulation}\\
		
		\textbf{Étape 1} Identification des variables de décision.\\
		\begin{itemize}
			\item[$\diamond $] $ x_{1}:$ la surface allouée à la culture des tomates
			\item[$\diamond $] $ x_{2}: $ la surface allouée à la culture des piments.
		\end{itemize}
		On vérifie bien que $ x_{1}\geq0 $ et $ x_{2}\geq0 $.\\
		
		\textbf{Étape 2} Identification des contraintes\\
		\begin{itemize}
			\item[$ \diamond $] Terrain : $ x_{1}+x_{2}\leq150 $
			\item[$\diamond $] Eau : $ 4x_{1}+2x_{2}\leq440 $
			\item[$\diamond $] Main d'œuvre : $ x_{1}+4x_{2}\leq480 $
			\item[$\diamond $] La limitation du bureau du périmètre irrigué $ x_{1}\leq 90 $.
		\end{itemize}
		\textbf{Étape 3} Identification de la fonction objectif\\
		$ z=10^{6}x_{1}+2.10^{6}x_{2} $.\\
		
		Le programme linéaire qui modélise le problème d'agriculture est :
		\begin{equation}\label{key}
		\begin{array}{lc}
		\displaystyle\max\:\:10^{6}x_{1}+2.10^{6}x_{2}  \\
		sc \left\lbrace \begin{array}{llc}
		x_{1}+x_{2}\leq150 \\
		4x_{1}+2x_{2}\leq440\\
		x_{1}+4x_{2}\leq480 \\
		x_{1}\leq 90 \\
		x_{1}\geq0,\: x_{2}\geq0.
		\end{array}\right.
		\end{array}
		\end{equation}
	\end{exemple}
	
	
	\begin{exemple}[problème de production]
		Une usine fabrique 2 produits $P_1$ et $P_2$ en utilisant un certain nombre de ressources : équipement, main d'œuvre, matières premières. Ces besoins sont indiqués dans le tableau ci-dessous. Par ailleurs, chaque ressource est disponible en quantité limitée.
		
		\begin{center}
			\begin{tabular}{|l|c|c|c|}
				\hline\hline
				& $P_1$ & $P_2$ & Disponibilité \\
				\hline\hline
				équipement & 3 & 9 & 81 \\
				\hline
				main d'œuvre & 4 & 5 & 55 \\
				\hline
				matière première & 2 & 1 & 20 \\
				\hline\hline
			\end{tabular}
		\end{center}
		
		Les deux produits $P_1$ et $P_2$ rapportent à la vente respectivement des bénéfices de 6 euros et 4 euros par unité. Quelles quantités de produits $P_1$ et $P_2$ (valeurs non-nécessairement entières) doit produire l'usine afin de maximiser le bénéfice total venant de la vente des 2 produits ?
		
		\begin{itemize}
			\item \textbf{Choix des variables (les inconnues)} : $x_1$ et $x_2$ sont respectivement les quantités des produits $P_1$ et $P_2$ fabriqués ($x_1, x_2 \in \mathbb{R}$).
			
			\item \textbf{Choix de la fonction objectif à maximiser} : La fonction objectif $f$ correspond au bénéfice total. Elle vaut $f(x_1, x_2) = 6x_1 + 4x_2$. Le problème se traduit donc par
			\[ \displaystyle\max_{(x_1, x_2)} \Big[f(x_1, x_2) = 6x_1 + 4x_2\Big]. \]
			
			\item \textbf{Détermination des contraintes.}
			\begin{itemize}
				\item La disponibilité de chacune des ressources s'écrit :
				\[
				\begin{cases}
				3x_1 + 9x_2 \leq 81 \\
				4x_1 + 5x_2 \leq 55 \\
				2x_1 + x_2 \leq 20
				\end{cases}
				\]
				\item Positivité des variables : $x_1, x_2 \geq 0$.
			\end{itemize}
		\end{itemize}
		
		En résumé, le problème de production se modélise sous la forme
		\begin{equation}
		\left\{
		\begin{aligned}
		& \displaystyle\max_{(x_1,x_2)} \Big[f(x_1,x_2) = 6x_1 + 4x_2\Big] \\
		& \text{s.t. } && 3x_1 + 9x_2 \le 81 \\
		&      && 4x_1 + 5x_2 \le 55 \\
		&      && 2x_1 + x_2 \le 20 \\
		&      && x_1, x_2 \ge 0
		\end{aligned}
		\right.
		\end{equation}
	\end{exemple}
	
	%	\textbf{Exemple 2}\\
	\begin{exemple}
		Un spécialiste en médecine a fabriqué un médicament (des pilules) pour guérir les sujets atteints d’un rhume. Ces pilules sont fabriquées selon deux formats :
		\begin{itemize}
			\item[$\diamond $] Petite taille : elle contient 2 grains d’aspirine, 5 grains de bicarbonate et 1 grain de codéine.
			\item[$\diamond $] Grande taille : elle contient 1 grain d’aspirine, 8 grains de bicarbonate et 6 grains de codéine.
		\end{itemize}
		Pour guérir la maladie, le sujet a besoin de 12 grains d’aspirine, 74 grains de bicarbonate et 24 grains de codéine. Déterminer le nombre de pilules minimales à prescrire au sujet pour qu’il soit guéri.\\
		
		\textbf{Formulation}
		\begin{equation}\label{key}
		\begin{array}{lc}
		\displaystyle\min\:\:x_{1}+x_{2}  \\
		sc \left\lbrace \begin{array}{llc}
		2x_{1}+x_{2}\geq12 \\
		5x_{1}+8x_{2}\geq74\\
		x_{1}+6x_{2}\geq24 \\
		x_{1}\geq0,\: x_{2}\geq0.
		\end{array}\right.
		\end{array}
		\end{equation}
	\end{exemple}
	
	\chapter{Résolution graphique}
	
	Dans le cas d'un (PL) à deux variables, on peut envisager une résolution graphique. Les contraintes où apparaissent des inégalités correspondent géométriquement à des demi-plans. L'intersection de ces demi-plans forme l'ensemble des variables satisfaisant à toutes les contraintes.
	
	\section{Technique de résolution graphique dans $\mathbb{R}^{2}$}
	
	La résolution graphique d'un programme linéaire ne peut se faire que pour les dimensions 2 et 3. Elle est donc peu efficace, mais permet de bien comprendre le rôle de PL.
	Dans $\mathbb{R}^{2}$, la  représentation graphique des éléments du problème s'effectue dans le cadre d'un repère où les valeurs de $x_{1}$ figurent sur l'axe des abscisses et celle de $x_{2}$ sur l'axe des ordonnées.\\
	Les étapes de résolution sont les suivantes :
	\begin{itemize}
		\item[$ \diamond $]\textbf{ Étape 1} \\
		Tracer le domaine des solutions admissibles, cette étape consiste à représenter graphiquement les différentes contraintes du système dans le repère. Et cette représentation permet de localiser l'ensemble $ D$.
		\item[$ \diamond $] \textbf{Étape 2} \\
		Tracer les droites isomarge  (resp isocoût si on est dans un problème de minimisation ) ie les droites $z=c_{1}x_{1}+c_{2}x_{2}$.
		Dans le cas d'une minimisation pour diminuer la valeur de $z=c_{1}x_{1}+c_{2}x_{2}$, il suffit de déplacer la droite d'isocoût en suivant le sens inverse d'un vecteur $\overrightarrow{u}\begin{pmatrix}
		c_{1}\\
		c_{2}
		\end{pmatrix}$ et pour augmenter la valeur de $z$, il faut suivre le sens du vecteur $\overrightarrow{u}$. Il est donc nécessaire de représenter le sens du vecteur $\overrightarrow{u}$.
		\item[$ \diamond $] \textbf{Étape 3} \\
		L'optimum est alors la solution réalisable située sur la droite isomarge la plus haute pour le cas de maximisation et respectivement sur la droite isocoût la plus basse le cas de minimisation.
	\end{itemize}
	
	Dans cette section, nous allons étudier deux types de problèmes :
	
	\subsection{Résolution graphique d'un problème de maximisation}
	
	Reprenons l'exemple de production.
	
	$$
	\left\{\begin{aligned}
	\displaystyle\max_{\left(x_{1}, x_{2}\right)}\Big[f\left(x_{1}, x_{2}\right)=6 x_{1}+4 x_{2}\Big] &  \\
	\text { s.t. } \quad 3 x_{1}+9 x_{2} & \leq 81 \\
	4 x_{1}+5 x_{2} & \leq 55 \\
	2 x_{1}+x_{2} & \leq 20 \\
	x_{1}, x_{2} & \geq 0
	\end{aligned}\right.
	$$
	
	A la fonction objectif $f$ correspond une droite $f\left(x_{1}, x_{2}\right)=6 x_{1}+4 x_{2}=$ constante, de coefficient directeur $\left(-\dfrac{6}{4}, 1\right)$. La constante précédente qui définie la droite doit être la plus grande possible (maximisation) et rencontrer l'ensemble des variables qui satisfont les contraintes. Pour déterminer cette valeur maximale, on fait donc "glisser" la droite (translation parallèle à la direction de la droite) du haut vers le bas jusqu'à rencontrer l'ensemble des variables satisfaisant les contraintes. Le maximum de $f$ sur cet ensemble des contraintes est alors atteint. On obtient ainsi la solution optimale $\left(x_{1}, x_{2}\right)=\left(\dfrac{15}{2},5\right)$ et ce qui donne une valeur maximale $\displaystyle\max (f)=65$.
	
	On remarque que l'ensemble des contraintes (la partie hachurée de la figure \ref{fig2.1}) est un polygone convexe et que le maximum de $f$ est atteint en un sommet de ce polygone. Cette observation est, en fait, un résultat général que l'on établira dans les sections suivantes.
	
	\begin{figure}[h]
		\begin{center}
			\includegraphics[scale=0.35]{2025_10_05_27e91bcb17e07a381837g-03}
			%			\captionsetup{labelformat=empty}
			\caption{Résolution graphique du problème de production}
			\label{fig2.1}
		\end{center}
	\end{figure}
	
	\subsection{Résolution graphique d'un problème de minimisation}
	
	Pour résoudre le problème linéaire dans le cas d'une minimisation, nous allons appliquer une démarche équivalente, mais dans ce cas, nous cherchons la droite (représentant la fonction économique) la plus rapprochée de l'origine.
	
	Résoudre le problème linéaire suivant :
	
	$$
	\left\{\begin{aligned}
	\displaystyle\min \quad & 6 x_{1}+9 x_{2}  \\
	\text { s.t } \quad & 3 x_{1}+2 x_{2}  \geq 18 \\
	& x_{1}+3 x_{2}  \geq 12 \\
	& x_{1}+x_{2}  \geq 8 \\
	& x  \geq 0
	\end{aligned}\right.
	$$
	
	Représentons graphiquement l'ensemble-solution puis représentons la famille de droites correspondante à :
	
	$$
	\begin{aligned}
	6 x_{1}+9 x_{2}=p & \Longleftrightarrow x_{2}=-\dfrac{2}{3} x_{1}+\dfrac{p}{9} \\
	& \Longleftrightarrow \text { droite de pente } -\dfrac{2}{3}
	\end{aligned}
	$$
	
	\begin{figure}[h]
		\begin{center}
			\includegraphics[scale=0.35]{2025_10_05_27e91bcb17e07a381837g-04}
			%			\captionsetup{labelformat=empty}
			\caption{Résolution graphique d'un problème de minimisation}
		\end{center}
	\end{figure}
	
	\textbf{Exercice 2.1} Résoudre graphiquement le problème suivant
	
	$$
	\left\{\begin{aligned}
	\displaystyle\min \quad & -50 x_{1}-30 x_{2}  \\
	\text { s.t } \quad & 10 x_{1}+6 x_{2}  \leq 45 \\
	& 4 x_{1}+6 x_{2}  \leq 36 \\
	& 2 x_{1}+6 x_{2}  \leq 27 \\
	& x  \geq 0
	\end{aligned}\right.
	$$
	
	\chapter{Résolution algébrique}
	
	Pour résoudre des problèmes de programmation linéaire ayant un nombre de variables (principales) supérieur à trois, nous devons obligatoirement faire appel à une autre méthode de résolution que celle déjà étudiée. Toutefois, avant d'aborder la méthode du simplexe comme technique de résolution d'un problème de programmation linéaire, nous allons découvrir une version modifiée "la méthode algébrique" qui permettra d'étudier les différentes étapes à franchir pour arriver à une solution optimale.
	
	\section{Variable d'écarts}
	
	La méthode de résolution que nous étudions dans ce chapitre nécessite que les contraintes du modèle soient exprimées sous forme d'équations linéaires au lieu d'inéquations.
	
	On peut facilement transformer une inéquations linéaire ayant un signe $\leq$ en une équation linéaire en additionnant une variable non négative dite variable d'écart.
	
	\begin{proposition}
		Tout PL sous forme standard s'écrit de façon équivalente en un PL sous forme canonique pure et inversement.
	\end{proposition}
	\label{prop:2.1}
	
	%	\section{Démonstration}
	\begin{proof}
		\begin{enumerate}
			\item Soit un PL sous forme canonique pure. On a
			$$
			\begin{aligned}
			A x \leq b \Longleftrightarrow & \displaystyle\sum_{j=1}^{n} a_{i j} x_{j}+e_{i}=b_{i}, \quad \forall i=1, \ldots, m \\
			& \text { où } e_{i}=b_{i}-\displaystyle\sum_{j=1}^{n} a_{i j} x_{j} \geq 0
			\end{aligned}
			$$
			
			Ainsi,
			
			$$
			\left\{\begin{array}{r} 
			A x \leq b , \\
			x \geq 0 .
			\end{array} \Leftrightarrow \left\{\begin{array}{r} 
			(A \mid I_{m}) \begin{pmatrix} x \\ e \end{pmatrix} = b , \\
			\begin{pmatrix} x \\ e \end{pmatrix} \geq 0 .
			\end{array} \Leftrightarrow \left\{\begin{array}{r}
			\tilde{A} \tilde{x}=b \\
			\tilde{x} \geq 0 .
			\end{array}\right.\right.\right.
			$$
			
			où $e=(e_{1}, \ldots, e_{m})^{t}$ sont appelées variables d'écart et $\tilde{A}$ est une matrice de taille $m \times n+m$.\\
			\item  Soit un PL sous forme standard. On a
			
			$$
			A x=b \Longleftrightarrow\left\{\begin{array}{c} 
			A x \leq b , \\
			A x \geq b .
			\end{array} \Longleftrightarrow \left\{\begin{array}{c}
			A x \leq b, \\
			-A x \leq-b .
			\end{array} \Longleftrightarrow\begin{pmatrix} A \\ -A \end{pmatrix} x \leq\begin{pmatrix} b \\ -b \end{pmatrix} \Longleftrightarrow \tilde{\tilde{A}} x \leq \tilde{\tilde{b}}\right.\right.
			$$
			
			où $\tilde{\tilde{A}}$ est une matrice de taille $2 m \times n$ et $\tilde{\tilde{b}} \in \mathbb{R}^{2 m}$.
		\end{enumerate}
	\end{proof}
	
	\begin{remarque}
		Pour transformer un programme linéaire sous forme canonique en sa forme standard (où toutes les contraintes sont des égalités avec un second membre positif), on introduit des variables supplémentaires non-négatives. La méthode varie selon le type d'inégalité :
		\begin{itemize}
			\item \textbf{Pour une contrainte de type "inférieur ou égal" ($\leq$)} : On ajoute une \textbf{variable d'écart} ($e_i \geq 0$) pour transformer l'inégalité en une égalité.
			\[ a_{i,1}x_1 + \dots + a_{i,n}x_n \leq b_i \quad \Rightarrow \quad a_{i,1}x_1 + \dots + a_{i,n}x_n + e_i = b_i \]
			
			\item \textbf{Pour une contrainte de type "supérieur ou égal" ($\geq$)} : On soustrait une \textbf{variable d'excès} ($e_i \geq 0$) et on ajoute une \textbf{variable artificielle} ($r_i \geq 0$).
			\[ a_{i,1}x_1 + \dots + a_{i,n}x_n \geq b_i \quad \Rightarrow \quad a_{i,1}x_1 + \dots + a_{i,n}x_n - e_i + r_i = b_i \]
			
			\item \textbf{Pour une contrainte de type "égal" (=)} : On ajoute simplement une \textbf{variable artificielle} ($r_i \geq 0$).
			\[ a_{i,1}x_1 + \dots + a_{i,n}x_n = b_i \quad \Rightarrow \quad a_{i,1}x_1 + \dots + a_{i,n}x_n + r_i = b_i \]
		\end{itemize}
		Les \textbf{variables d'écart et d'excès} représentent la différence entre les deux membres de l'inéquation. Les \textbf{variables artificielles} sont des outils de calcul temporaires, indispensables pour démarrer l'algorithme du simplexe. Elles doivent impérativement être nulles dans la solution finale, ce qui est généralement réalisé à l'aide de la méthode des deux phases.
	\end{remarque}
	
	\section{Solutions de base réalisables}
	
	On considère désormais (sauf mention contraire) un PL toujours sous forme standard.
	
	\begin{definition}
		On appelle solution réalisable tout vecteur $x$ qui satisfait les contraintes du PL i.e. tel que $A x=b$ et $x \geq 0$.
	\end{definition}
	\label{def:2.1}
	
	\begin{hypothese}
		On suppose désormais que la matrice $A$ est de taille $m \times n$ avec $\operatorname{rang}(A)=m \leq n$.\\
		On rappelle que le rang de $A$ est le nombre maximal de lignes de $A$ linéairement indépendantes. C'est aussi le nombre de colonnes de $A$ linéairement indépendantes. L'hypothèse de rang plein n'est pas restrictive car :
	\end{hypothese}
	
	\begin{itemize}
		\item si $\operatorname{rang}(A)<m$ le système $A x=b$ n'a pas de solution en général.
		\item Si $\operatorname{rang}(A)<m$ et $b \in \operatorname{Im}(A)$, il y a des équations redondantes dans le système $A x=b$, qu'on peut donc supprimer pour obtenir un nouveau système de rang plein. Sous l'hypothèse de rang plein, le système $A x=b$ admet toujours des solutions.
		\item Si $m<n$, il y en a une infinité.
		\item Si $m=n$, la solution est unique et vaut $x=A^{-1} b$, dans ce cas, il n'y a rien à maximiser...
	\end{itemize}
	
	\begin{remarque}
		Pour la forme canonique pure avec $A x \leq b$, l'hypothèse de rang plein n'est pas non plus restrictive car si $\operatorname{rang}(A)<m$, il y a des contraintes d'inégalités redondantes qu'on peut supprimer pour obtenir un système de rang plein.
	\end{remarque}
	\label{rem:2.1}
	
	\begin{definition}
		Soit $B \subseteq\{1, \ldots, n\}$ un ensemble d'indices avec $\operatorname{card}(B)=m$ tel que les colonnes $A_{j}, j \in B$, de $A$ sont linéairement indépendantes. Autrement dit, la matrice carrée $A_{B}$ formée des colonnes $A_{j}$, $j \in B$, est inversible. On dit que l'ensemble $B$ des indices est une base.
	\end{definition}
	\label{def:2.2}
	
	\begin{itemize}
		\item Les variables $x_{B}=\left(x_{j}, j \in B\right)$ sont appelées variables de base.
		\item Les variables $x_{H}=\left(x_{j}, j \notin B\right)$ sont appelées variables hors-base.
	\end{itemize}
	
	\begin{remarque}
		\begin{itemize}
			\item Sous l'hypothèse de rang plein, il existe toujours une base non vide.
			\item Quitte à re-numéroter les indices, on peut toujours écrire les décompositions par blocs :
		\end{itemize}
	\end{remarque}
	\label{rem:2.2}
	
	$$
	\begin{aligned}
	A & =\left(A_{B} \mid A_{H}\right) \text { où } A_{H} \text { est la matrice formée des colonnes } A_{j}, j \notin B \\
	x & =\left(x_{B} \mid x_{H}\right)^{t}
	\end{aligned}
	$$
	
	Le système $A x=b$ est équivalent à $A_{B} x_{B}+A_{H} x_{H}=b$.
	
	Par la relation précédente et du fait que la matrice $A_{B}$ est inversible, on peut fixer les variables hors-base et les variables de base sont alors complètement déterminées.
	
	\begin{definition}
		On dit que $x=\left(x_{B} \mid x_{H}\right)^{t}$ est solution de base associée à la base $B$ si $$ x_{H}=0 $$.
	\end{definition}
	\label{def:2.3}
	
	\textbf{Propriétés des solutions de base réalisables :} Si $x=\left(x_{B} \mid x_{H}\right)^{t}$ est une solution de base réalisable alors $x_{H}=0$ et $x_{B}=A_{B}^{-1} b$.
	
	%	\textbf{Exemple 2.1} 
	\begin{exemple}
		Reprenons l'exemple du problème de production de l'introduction. Sous forme standard (en introduisant des variables d'écart), le PL s'écrit
		
		$$
		\left\{\begin{array}{rl} 
		\displaystyle\max_{( x_{1} , x_{2} )} & f ( x_{1} , x_{2} ) = 6 x_{1} + 4 x_{2}  \\
		\text { s.t. }  & 3 x_{1} + 9 x_{2} \leq 81  \\
		& 4 x_{1} + 5 x_{2} \leq 55  \\
		& 2 x_{1} + x_{2} \leq 20  \\
		& x_{1} , x_{2} \geq 0 
		\end{array} \right. \Longrightarrow \left\{\begin{array}{rl} 
		\displaystyle\max_{( x_{1} , x_{2} , e_{1} , e_{2} , e_{3} )} & f ( x_{1} , x_{2} , e_{1} , e_{2} , e_{3} ) = 6 x_{1} + 4 x_{2} + 0 e_{1} + 0 e_{2} + 0 e_{3}  \\
		\text { s.t. } & 3 x_{1} + 9 x_{2} + e_{1} = 81  \\
		& 4 x_{1} + 5 x_{2} + e_{2} = 55  \\
		& 2 x_{1} + x_{2} + e_{3} = 20  \\
		& x_{1} , x_{2} , e_{1} , e_{2} , e_{3} \geq 0 
		\end{array} \right.
		$$
		
		On a $m=3, n=5$ et $\operatorname{rang}(A)=m=3$.
		
		Nous sommes alors en présence d'un système de plus d'inconnues que d'équations.
		
		Rappelons qu'il est possible de résoudre un système d'équations ayant plus d'inconnues que d'équations ( $m \leq n$ ) en fixant arbitrairement (par une valeur ou un paramètre) les valeurs de $n-m$ variables.
		
		\begin{description}
			\item[Programme initial] \quad\\	
			Pour déterminer le programme initial, on pose habituellement à zéro les variables principales du modèle, ceci correspond à ne fabriquer aucun produit :
			
			$$
			x_{1}=0 \text { et } x_{2}=0 .
			$$
			
			Notre système de 3 équations à 5 inconnues devient alors un système de 3 équations à 3 inconnues que l'on va pouvoir résoudre. Une base est donnée par $B=\{3,4,5\}$ avec $A_{B}=\left(\begin{array}{lll}1 & 0 & 0 \\ 0 & 1 & 0 \\ 0 & 0 & 1\end{array}\right)$. La première solution de base réalisable correspondante est $x=\left(x_{1}, x_{2}, e_{1}, e_{2}, e_{3}\right)^{t}=(\underbrace{0,0}_{x_{H}}, \underbrace{81,55,20}_{x_{B}=A_{B}^{-1} b})^{t}$, la valeur de la fonction économique est : 0. La première solution de base réalisable correspond
			
			\begin{figure}[H]
				\begin{center}
					\includegraphics[scale=0.35]{2025_10_05_27e91bcb17e07a381837g-07}
					%			\captionsetup{labelformat=empty}
					\caption{Résolution graphique du problème de production}
				\end{center}
			\end{figure}
			
			au point extrême $(0,0)^{t}$ de la résolution graphique. On peut évidemment améliorer la valeur de la fonction économique en fabriquant quelques unités du produit $P_{1}$ ou du produit $P_{2}$.
			
			\item[Révision du programme initial] \quad\\
			On doit maintenant examiner la possibilité d'améliorer la valeur de la fonction économique en introduisant l'une ou l'autre des variables principales (soit $x_{1}$, soit $x_{2}$ mais non les deux simultanément) dans le programme de base.\\
			Le choix de la variable (actuellement hors base) à introduire dans le programme s'effectue à l'aide d'un critère qui permet d'améliorer le plus rapidement possible la valeur de la fonction\\
			économique, c-à-d on choisira la variable qui donne la meilleure augmentation de la valeur de la fonction économique. Dans notre exemple la fonction économique est
			
			$$
			f\left(x_{1}, x_{2}\right)=6 x_{1}+4 x_{2}+0 e_{1}+0 e_{2}+0 e_{3}
			$$
			
			l'augmentation sera la plus élevée si nous introduisons la variable $x_{1}$ dans le programme, Donc : $x_{1}$ est une variable entrante de base\\
			Il nous reste maintenant deux choses à déterminer. D'une part, calculer la quantité du $P_{1}$ que l'on doit fabriquer et d'autre part, quelle variable actuellement dans le programme doit devenir une variable hors programme c.-à-d. déterminer la variable sortante.
			
			\item[Détermination de la valeur entrante] \quad\\			
			Déterminons la plus grande valeur pouvant être prise par $x_{1}$ dans chaque équation.
			
			\begin{align}
			x_{1} & =27-3 x_{2}-\dfrac{1}{3} e_{1}  \label{2.44}\\
			x_{1} & =\dfrac{55}{4}-\dfrac{5}{4} x_{2}-\dfrac{1}{4} e_{2}  \label{2.5}\\
			x_{1} & =10-\dfrac{1}{2} x_{2}-\dfrac{1}{2} e_{3} \label{2.6}
			\end{align}
			
			
			Pour s'assurer que toutes les contraintes soient respectées en même temps, on choisit parmi ces valeurs la plus petite quantité positive :
			
			$$
			x_{1}=10 .
			$$
			
			Ce minimum s'obtient dans l'équation \eqref{2.6}, ce qui se traduit par $e_{3}=0$ et donc c'est la variable $e_{3}$ qui sort du programme :
			
			\begin{equation}
			e_{3}\text{ est une variable hors base.}
			\end{equation}
			
			\item[Détermination de la nouvelle solution : solution de base № 2]\quad\\
			Une deuxième base est donnée par $B=\{1,3,4\}$ avec $A_{B}=\left(\begin{array}{lll}3 & 1 & 0 \\ 4 & 0 & 1 \\ 2 & 0 & 0\end{array}\right)$. La première solution de base réalisable correspondante est $x=\left(x_{1}, x_{2}, e_{1}, e_{2}, e_{3}\right)^{t}=(\underbrace{10}_{x_{B}}, \underbrace{0}_{x_{H}}, \underbrace{51,15}_{x_{B}}, \underbrace{0}_{x_{H}})^{t}$, la valeur de la fonction économique est : 60.\\
			Il correspond au point-sommet $(10,0)^{t}$ de la résolution graphique. S'agit-il d'une solution optimal? Toutefois, avant de tester si le programme est optimal, il faut d'abord transformer les contraintes de mon problème pour exprimer les variables de base en fonction des variables hors-base:
			
			\begin{itemize}
				\item les variables de base sont : $x_{1}, e_{1}, e_{2}$.
				\item les variables hors-base sont : $x_{2}, e_{3}$.
			\end{itemize}
			
			\begin{align}
			\left\{\begin{array}{l}e_{1}=81-3 x_{1}-9 x_{2} \\ e_{2}=55-4 x_{1}-5 x_{2} \\ e_{3}=20-2 x_{1}-x_{2}\end{array}\right.& \Rightarrow\left\{\begin{array}{l}e_{1}=81-3\left(10-\dfrac{1}{2} x_{2}-\dfrac{1}{2} e_{3}\right)-9 x_{2} \\ e_{2}=55-4\left(10-\dfrac{1}{2} x_{2}-\dfrac{1}{2} e_{3}\right)-5 x_{2} \\ x_{1}=10-\dfrac{1}{2} x_{2}-\dfrac{1}{2} e_{3}\end{array}\right.\nonumber\\
			&\Rightarrow \begin{cases}e_{1} & =51-\dfrac{15}{2} x_{2}+\dfrac{3}{2} e_{3} \\ e_{2} & =15-3 x_{2}-2 e_{3} \\ x_{1} & =10-\dfrac{1}{2} x_{2}-\dfrac{1}{2} e_{3}\end{cases}\label{equ2.7}
			\end{align}
%			\[
%			
%			\]
%			
%			\[ \]
			
			Il nous reste à exprimer la valeur de la fonction économique, en fonction des variables hors-base $x_{2}$ et $e_{3}$ à l'aide de la substitution :
			
			$$
			f_{o b j}=6 x_{1}+4 x_{2}=60+x_{2}-3 e_{3} .
			$$
			
			Cette dernière expression nous indique qu'on peut améliorer encore la valeur de la fonction économique puisque le coefficient de la variable $x_{2}$ est positif.
			
			\item[Détermination de la troisième base : variable entrante est $x_{2}$]\quad\\
			Déterminons la valeur de $x_{2}$, par le même principe, on obtient que la variable $x_{2}$ peut prendre la valeur $20, \dfrac{102}{15}$ ou 5 , Ainsi $x_{2}$ devant respecter les différentes contraintes, nous choisirons $x_{2}=5$ qui se réalise dans l'équation $e_{2}=15-3 x_{2}+2 e_{3}$ et qui impose que $e_{2}$ devienne nulle, ainsi : $e_{2}$ devient variable hors base.
			
			\begin{itemize}
				\item les variables de base sont : $x_{1}, e_{1}, x_{2}$.
				\item les variables hors-base sont : $e_{2}, e_{3}$.
			\end{itemize}
			
			Exprimons à nouveau les variables dans le programme en fonction des variables hors-base.
			
			$$
			\left\{\begin{aligned}
			x_{1} & =\dfrac{15}{2}+\dfrac{1}{6} e_{2}-\dfrac{1}{3} e_{3} \\
			e_{1} & =\dfrac{27}{2}+\dfrac{5}{2} e_{2}-\dfrac{7}{2} e_{3} \\
			x_{2} & =5-\dfrac{1}{3} e_{2}+\dfrac{2}{3} e_{3}
			\end{aligned}\right.
			$$
			
			La valeur de la fonction économique est alors :
			
			$$
			f_{o b j}=6 \times \dfrac{15}{2}+4 \times 5=65
			$$
			
			Puis on exprime la valeur de la fonction économique, en fonction des variables hors-base $e_{2}$ et $e_{3}$ à l'aide de la substitution :
			
			$$
			\begin{aligned}
			f_{o b j} & =6 \times\left(\dfrac{15}{2}+\dfrac{1}{6} e_{2}-\dfrac{1}{3} e_{3}\right)+4 \times\left(5-\dfrac{1}{3} e_{2}+\dfrac{2}{3} e_{3}\right), \\
			& =65-\dfrac{1}{3} e_{2}-\dfrac{7}{3} e_{3} .
			\end{aligned}
			$$
			
			Pouvons-nous améliorer encore la valeur de la fonction économique?\\
			La réponse est non puisque les coefficients des variables hors-base $e_{2}$ et $e_{3}$ sont négatifs.
		\end{description}
	\end{exemple}
	\begin{remarque}
		Un programme de base est optimal si tous les coefficients des variables hors-base, dans l'expression de la fonction économique, sont $\leq 0$.
	\end{remarque}
	\label{rem:2.3}
	
	On constate qu'elle correspond bien à la solution optimale obtenue lors de la résolution graphique.
	
	\begin{remarque}
		Il y a au plus $C_{n}^{m}$ solutions de base (toutes ne sont pas réalisables).
	\end{remarque}
	\label{rem:2.4}
	
	\section{Propriétés géométriques des solutions de base réalisables}
	
	On note
	
	\begin{equation*}
	\mathbf{D}_{R}=\left\{x \in \mathbb{R}^{n} \mid A x=b, x \geq 0\right\}, \tag{2.8}
	\end{equation*}
	
	l'ensemble des solutions réalisables d'un PL sous forme standard.\\
	Commençons par rappeler les notions d'ensemble convexe et de polyèdre :\\
	\begin{minipage}[t]{0.45\textwidth}
		\centering
		\includegraphics[max width=\textwidth, center]{2025_10_05_27e91bcb17e07a381837g-10(1)}
	\end{minipage}
	\begin{minipage}[t]{0.45\textwidth}
		\centering
		\includegraphics[max width=\textwidth, center]{2025_10_05_27e91bcb17e07a381837g-10}
	\end{minipage}
	
	\begin{definition}
		Un ensemble $E$ est dit convexe si $\forall x, y \in E, \lambda x+(1-\lambda) y \in E$ pour tout $0 \leq \lambda \leq 1$.
	\end{definition}
	\label{def:2.4}
	
	\begin{definition}
		Un polyèdre $P \subseteq \mathbb{R}^{n}$ est l'ensemble de solutions d'un système fini d'inégalités linéaires, i.e.,
		
		$$
		P=\left\{x \in \mathbb{R}^{n} \mid A x \leq b\right\}
		$$
		
		où $A$ est une matrice $m \times n$ et $b \in \mathbb{R}^{m}$, et $m$ et $n$ sont deux entiers positifs.
		
		Un polyèdre borné est appelé polytope. En d'autres termes, un polyèdre $P \subseteq \mathbb{R}^{n}$ est un polytope s'il existe $x, y \in \mathbb{R}^{n}$ tel que $x \leq z \leq y$ pour tout $z \in P$.
		
		Noter que tout polyèdre est convexe, c'est-à-dire si $x$ et $y$ sont deux points quelconques de $P$, alors $\lambda x+(1-\lambda) y \in P$ pour tout $0 \leq \lambda \leq 1$.
	\end{definition}
	\label{def:2.5}
	
	\begin{remarque}
		L'ensemble $\mathbf{D}_{R}$ des solutions réalisables est un polyèdre convexe, fermé.
	\end{remarque}
	\label{rem:2.5}
	
	%	\textbf{Exemple 2.2} 
	\begin{exemple}
		L'ensemble $\mathbf{D}_{R}=\left\{x \in \mathbb{R}^{3} \mid 2 x_{1}+3 x_{2}+x_{3}=3, x_{1}, x_{2}, x_{3} \geq 0\right\}$ est représenté sur la figure \ref{2.4}.
	\end{exemple}
	
	\begin{definition}
		Un modèle mathématique est dit non borné si son domaine réalisable est non borné, ç-à-d si au moins l'une de ses variables peut tendre vers l'infini sans quitter le domaine réalisable.
	\end{definition}
	\label{def:2.6}
	
	\begin{definition}
		Un point $x \in \mathbf{D}_{R}$ est un sommet (ou point extrême) si et seulement s'il n'existe pas $y, z \in \mathbf{D}_{R}, y \neq z$ tels que $x=\lambda y+(1-\lambda) z$ avec $0<\lambda<1$.
	\end{definition}
	
	\begin{theorem}
		$x$ est une solution de base réalisable si et seulement si $x$ est un sommet de $\mathbf{D}_{R}$.
	\end{theorem}
	\label{thm:2.1}
	
	On a vu sur l'exemple de production que la solution était atteinte sur un sommet de $\mathbf{D}_{R}$ et correspond donc à une solution de base.
	
	\begin{figure}[h]
		\centering
		\begin{tikzpicture}[scale=1.2,
		x={(-0.707cm,-0.707cm)}, y={(1cm,0cm)}, z={(0cm,1cm)}]
		
		% Définir les sommets
		\coordinate (O) at (0,0,0);
		\coordinate (A) at (2,0,0);
		\coordinate (B) at (0,2,0);
		\coordinate (C) at (0,0,3);
		
		% Axes
		\draw[->, thick] (O) -- (3,0,0) node[anchor=north east]{$x_1$};
		\draw[->, thick] (O) -- (0,3,0) node[anchor=north west]{$x_2$};
		\draw[->, thick] (O) -- (0,0,4) node[anchor=south]{$x_3$};
		
		% Triangle avec hachures diagonales
		\fill[pattern=north east lines, pattern color=gray!40] (A) -- (B) -- (C) -- cycle;
		\draw[thick] (A) -- (B) -- (C) -- cycle;
		
		% Graduations sur les axes
		\foreach \x in {1, 2} {
			\draw (\x,0,0.1) -- (\x,0,-0.1) node[below] {\small \x};
		}
		\foreach \y in {1, 2} {
			\draw (0.1,\y,0) -- (-0.1,\y,0) node[below left] {\small \y};
		}
		\foreach \z in {1, 2, 3} {
			\draw (0,0.1,\z) -- (0,-0.1,\z) node[left] {\small \z};
		}
		
		\end{tikzpicture}
		\caption{Polyèdre des solutions réalisables}
		\label{2.4}
	\end{figure}
	
	\begin{proposition}
		Si le domaine réalisable $\mathbf{D}_{R}$ n'est pas vide, alors ce polytope comprend au moins un point extrême.
	\end{proposition}
	\label{prop:2.2}
	
	Le résultat général s'énonce ainsi :
	
	\begin{theorem}
		L'optimum de la fonction objectif $f$ sur $\mathbf{D}_{R}$, s'il existe, est atteint en au moins un sommet de $\mathbf{D}_{R}$.
	\end{theorem}
	\label{thm:2.2}
	
	%	\textbf{Démonstration} Voir les propriétés géométriques des polyèdres.
	
	Pour résoudre un PL sous forme standard, il suffit de se restreindre aux solutions de base réalisables (les sommets de $\mathbf{D}_{R}$). Tout se passe donc avec les solutions de base.
	
	L'ensemble $\mathbf{D}_{R}$ n'est pas nécessairement borné. En fait pour un PL, 3 situations (et seulement 3) peuvent se produire :
	
	\begin{enumerate}
		\item $\mathbf{D}_{R} = \varnothing$ : le PL n'a pas de solution.
		\item $\mathbf{D}_{R} \neq \varnothing$ mais la fonction objectif $f$ n'est pas majorée sur $\mathbf{D}_{R}$ : le maximum de $f$ vaut $+\infty$. Si $\mathbf{D}_{R}$ est borné, ce cas est exclu.
		\item $\mathbf{D}_{R} \neq \varnothing$ et la fonction objectif $f$ est majorée sur $\mathbf{D}_{R}$ : le PL admet une solution optimale (non nécessairement unique).
	\end{enumerate}
	
	\begin{remarque}
		Il y a au plus $C_{n}^{m}$ solutions de base réalisables. Pour déterminer une solution de base, on doit résoudre un système linéaire $(x_{B} = A_{B}^{-1} b)$. La résolution d'un système linéaire par une méthode directe de type Gauss/LU requiert de l'ordre de $m^{3}$ opérations. Si on explore toutes les solutions de base et que l'on compare les coûts correspondants, on effectue de l'ordre de $m^{3} C_{n}^{m}$ opérations. Ce nombre est vite très grand avec $n$ et $m$. Par exemple, avec $n=20$ et $m=10$, on a $3 \times 10^{8}$ opérations. Dans la méthode du simplexe, on va explorer seulement les sommets qui permettent d'augmenter la fonction objectif. On va réduire ainsi le nombre de solutions de base à explorer et donc le nombre de systèmes linéaires à résoudre.
	\end{remarque}
	\label{rem:2.6}
	
	\section{Marche à suivre de la méthode algébrique}
	
	La méthode algébrique est une recherche systématique de programmes de base (points sommets) jusqu'à l'obtention d'un programme optimal.
	
	Pour ce faire, il s'agit :
	
	\begin{colorframed}
		\begin{enumerate}
			\item De structurer le problème sous forme standard.
			\item De déterminer les variables de base qui serviront de départ au cheminement vers la solution optimale (programme optimal).
			\item D'expliciter la fonction économique et de déterminer si elle peut être améliorée : recherche de l'éventuelle variable (hors-base) admettant le plus grand coefficient positif.
			\item En introduisant cette variable dans le programme, on choisira la plus petite valeur positive obtenue à l'aide du système d'équations calculé lors de l'étape précédente. Cela induira la variable sortante.
			\item Pour déterminer une nouvelle variable de base, on doit exprimer la fonction économique ainsi que les nouvelles variables de base en fonction des variables hors-base.
			\item Retourner à 3) jusqu'à l'obtention du programme de base optimal.
			\item Donner le programme optimal en précisant la valeur de toutes les variables ainsi que la valeur optimisée de la fonction économique.
		\end{enumerate}
	\end{colorframed}
	\chapter{Méthode du simplexe}
	
	La présentation algébrique de la méthode du simplexe nous a permis de donner une interprétation pratique et économique aux différentes opérations à effectuer dans la résolution d'un problème de programmation linéaire pouvant comporter deux variables ou plus. Toutefois, la méthode algébrique devient rapidement laborieuse à mesure que le nombre de contraintes et le nombre de variables augmentent.
	
	Nous allons donc systématiser toutes ces opérations algébriques sous forme d'un simple tableau dit « tableau du simplexe ». Ceci permettra de résoudre plus facilement des problèmes ayant beaucoup plus de variables et de contraintes. Malgré l'efficacité de cette méthode, il n'en demeure pas moins que la résolution à l'aide d'un ordinateur devient indispensable à mesure que les problèmes prennent de fortes dimensions.
	
	\section{Principe de l'algorithme du simplexe}
	
	La méthode du simplexe est due à G. Dantzig (1947). Elle comporte 2 phases.
	\begin{itemize}
		\item \textbf{Phase I - Initialisation} : Trouver une solution de base réalisable (ou bien détecter l'impossibilité).
		\item \textbf{Phase II - Progression} : On passe d'un sommet à un sommet voisin pour augmenter la fonction objectif (ou bien on détecte une fonction objectif $f$ non majorée).
	\end{itemize}
	
	\subsection{L'algorithme du simplexe : la phase 2}
	
	Dans cette section, nous allons présenter la Phase II de la méthode du simplexe. La Phase I qui sert plus à initialiser la Phase II, sera abordée plus tard. Cette phase s'applique à des problèmes du type
	\[
	\left\{
	\begin{array}{rl}
	\displaystyle\max & c^{t} x \\
	\text{s.t.} & A x \leq b \\
	& x \geq 0
	\end{array}
	\quad \text{où} \quad
	\begin{array}{rl}
	\displaystyle\min_{x} & c^{t} x \\
	\text{s.t.} & A x \leq b \\
	& x \geq 0
	\end{array}
	\right.
	\]
	où $A$ est une matrice de taille $m \times n$.
	
	On fera l'hypothèse que $b \geq 0$. Cette supposition est cruciale pour la Phase II. Ceci garantit que $0 \in K = (x \geq 0, A x \leq b)$. De plus, nous savons que le point $0$ est un sommet. Ce point servira de point de départ de l'algorithme du simplexe. En gros, l'algorithme va pivoter autour de ce point pour trouver un meilleur sommet. On poursuit l'algorithme jusqu'à l'obtention de la solution optimale.
	
	Supposons qu'on dispose d'une solution de base réalisable $\bar{x}$ d'un PL sous forme standard. À une permutation près des colonnes, la matrice $A$ peut s'écrire
	\[
	A = \left( A_{B} \mid A_{H} \right)
	\]
	avec $A_{B}$ une matrice carrée de taille $m \times m$, inversible, correspondant aux variables de base et $A_{H}$ une matrice de taille $m \times (n-m)$, correspondant aux variables hors-base. On décompose également (à la même permutation près des composantes)
	\[
	\bar{x} = \left( \bar{x}_{B} \mid \bar{x}_{H} \right)^{t}
	\]
	Le but est de trouver une autre base $B^{*}$ et une solution de base $x^{*}$ associée telles que
	\[
	f(x^{*}) > f(\bar{x}) \quad \text{( $x^{*}$ est meilleur que $\bar{x}$ )}
	\]
	La méthode du simplexe consiste à faire entrer une variable hors-base dans la nouvelle base (variable entrante) et faire sortir à la place une variable de base (variable sortante).
	
	\subsubsection{Calcul des coûts réduits et variable entrante}
	
	On écrit la fonction objectif $f$ avec les variables de base/hors-base. On a $b = A x = A_{B} x_{B} + A_{H} x_{H}$ avec $A_{B}$ inversible donc $x_{B} = A_{B}^{-1} (b - A_{H} x_{H})$. On obtient donc
	\[
	\begin{aligned}
	f(x) &= c^{t} x = c_{B}^{t} x_{B} + c_{H}^{t} x_{H} \quad \text{avec} \quad c = \left( c_{B} \mid c_{H} \right)^{t} \\
	&= c_{B}^{t} A_{B}^{-1} (b - A_{H} x_{H}) + c_{H}^{t} x_{H} \\
	&= c_{B}^{t} A_{B}^{-1} b + \left( c_{H}^{t} - c_{B}^{t} A_{B}^{-1} A_{H} \right) x_{H}
	\end{aligned}
	\]
	Or $x_{B} = A_{B}^{-1} b$ (car $x_{H} = 0$) et $c_{B}^{t} A_{B}^{-1} b = c^{t} x = f(x)$ donc
	\[
	f(x) = f(x) + \left( c_{H}^{t} - c_{B}^{t} A_{B}^{-1} A_{H} \right) x_{H}
	\]
	Le vecteur
	\[
	L_{H}^{t} = c_{H}^{t} - c_{B}^{t} A_{B}^{-1} A_{H}
	\]
	s'appelle vecteur des coûts réduits.
	
	\begin{itemize}
		\item Si les coûts réduits sont tous négatifs, i.e. $L_{H}^{t} \leq 0$, il n'est alors pas possible d'augmenter la fonction objectif $f$ : l'algorithme se termine normalement, c'est-à-dire qu'on a trouvé une solution de base réalisable optimale.
		\item Dans le cas contraire (i.e. $\exists (L_{H})_{i} > 0$), on a intérêt à faire entrer dans la base la variable hors-base qui a le coût réduit positif le plus grand possible.
	\end{itemize}
	
	On note $e \in B$ l'indice de la variable entrante. On choisit $e$ tel que
	\[
	(L_{H})_{e} = \displaystyle\max_{j} \left\{ (L_{H})_{j}, (L_{H})_{j} > 0 \right\}
	\]
	ce qu'on note par
	\[
	e = \operatorname{argmax}_{j} \left\{ (L_{H})_{j}, (L_{H})_{j} > 0 \right\}
	\]
	\begin{remarque}
		Si on traite d'un problème de minimisation, c'est-à-dire avec $\displaystyle\min f(x)$, alors la variable entrante $x_{e}$ est déterminée par l'indice
		\[
		e = \operatorname{argmin}_{j} \left\{ (L_{H})_{j}, (L_{H})_{j} < 0 \right\}
		\]
	\end{remarque}
	\label{rem:3.1}
	
	\subsubsection{Variable sortante}
	
	Une fois l'indice $e$ choisi, il faut déterminer quelle variable doit quitter la base. En maintenant la relation $A x = b$, on augmente $x_{e}$ jusqu'à annuler une des variables de base. Cette variable sera alors la variable sortante.
	\[
	\begin{aligned}
	A x = b &\Longleftrightarrow A_{B} x_{B} + A_{e} x_{e} = b \quad \text{où $A_{e}$ désigne la $e$-ième colonne de $A$} \\
	&\Longleftrightarrow x_{B} = A_{B}^{-1} (b - A_{e} x_{e}) \\
	&\Longleftrightarrow x_{B} = \bar{x}_{B} - A_{B}^{-1} A_{e} x_{e} \\
	&\Longleftrightarrow x_{B} = \bar{x}_{B} - z x_{e}
	\end{aligned}
	\]
	où
	\[
	z = A_{B}^{-1} A_{e} \in \mathbb{R}^{m}
	\]
	\begin{itemize}
		\item Si $z \leq 0$, on peut augmenter $x_{e}$ autant qu'on veut, on aura toujours la positivité de la variable de base $x_{B}$. La fonction objectif n'est pas majorée sur $\mathbf{D}_{R}$, i.e. le maximum de $f$ vaut $+\infty$. Dans ce cas, l'algorithme s'arrête.
		\item Sinon (i.e. il existe $z_{i} > 0$), pour avoir la positivité $(\bar{x}_{B})_{i} - z_{i} x_{e} \geq 0$ pour tout $i$, on choisit la variable sortante pour laquelle le rapport $\dfrac{(\bar{x}_{B})_{i}}{z_{i}}$ pour $i = 1, \ldots, m$ avec $z_{i} > 0$, est le plus petit possible. On choisit
		\[
		s = \operatorname{argmin}_{i} \left\{ \dfrac{(\bar{x}_{B})_{i}}{z_{i}}, z_{i} > 0 \right\}
		\]
		On a, dans ce cas, $x_{s} = 0$ et $x_{B} \geq 0$. La valeur de la variable entrante est donnée par
		\[
		x_{e} = \displaystyle\min_{i} \left\{ \dfrac{(\bar{x}_{B})_{i}}{z_{i}}, z_{i} > 0 \right\}
		\]
	\end{itemize}
	
	%	\section{}
	\begin{colorframed}
		{\color{titleblue}\textbf{Algorithme (Résumé de la méthode du simplexe en phase 2 - progression)}}
		
		\medskip
		\begin{enumerate}
			\item \begin{itemize}
				\item Calcul des variables de base réalisables : étant donné $A = (A_{B} \mid A_{H})$, on calcule les variables de base réalisables $x_{B} = A_{B}^{-1} b \geq 0$.
				\item Calcul des coûts réduits : $L_{H}^{t} = c_{H}^{t} - c_{B}^{t} A_{B}^{-1} A_{H}$.
				\item Si $L_{H} \leq 0$ alors $\bar{x}_{B}$ est une solution optimale (arrêt de l'algorithme).
			\end{itemize}
			\item Variable entrante : $e = \operatorname{argmax}_{j} \left\{ (L_{H})_{j}, (L_{H})_{j} > 0 \right\}$.
			\item Variable sortante : Calcul de $z = A_{B}^{-1} A_{e}$ puis $s = \operatorname{argmin}_{i} \left\{ \dfrac{(\bar{x}_{B})_{i}}{z_{i}}, z_{i} > 0 \right\}$.
			\item On obtient une nouvelle base $B^{*}$ et une nouvelle matrice $A_{B^{*}}$ dans laquelle la colonne $A_{e}$ remplace la colonne $A_{s}$. Calcul de $A_{B^{*}}^{-1}$ et retour en 1.
		\end{enumerate}
	\end{colorframed}
	
	Nous allons maintenant systématiser cet algorithme en présentant la méthode du simplexe sous la forme d'un tableau.
	
	\section{Représentation tabulaire du simplexe}
	
	Afin de simplifier la présentation des différentes itérations de la méthode du simplexe, nous systématisons cette procédure sous forme de tableaux sur deux types de problèmes : de maximisation et de minimisation.
	
	\subsection{Cas général de la méthode tabulaire du simplexe}
	
	Considérons le programme linéaire suivant, exprimé sous sa forme standard, qui vise à maximiser une fonction objectif sous des contraintes d'égalité :
	
	\[
	\left\{
	\begin{array}{rcl}
	\displaystyle \displaystyle\max_{\left(x_{1}, \ldots, x_{n}\right)} z & = & c_{1} x_{1} + c_{2} x_{2} + c_{3} x_{3} + \cdots + c_{n} x_{n} \\
	\text{s.t.} & & a_{11} x_{1} + a_{12} x_{2} + \cdots + a_{1n} x_{n} = b_{1} \\
	& & a_{21} x_{1} + a_{22} x_{2} + \cdots + a_{2n} x_{n} = b_{2} \\
	& & \vdots \\
	& & a_{m1} x_{1} + a_{m2} x_{2} + \cdots + a_{mn} x_{n} = b_{m} \\
	& & x_{1}, \ldots, x_{n} \geq 0
	\end{array}
	\right.
	\]
	
	Ce programme peut être reformulé de manière équivalente en introduisant la variable $z$ dans les contraintes, comme suit :
	
	\begin{equation}\label{pro3.2}
	\left\{
	\begin{array}{rcl}
	\displaystyle \displaystyle\max z & & \\
	\text{s.t.} & & a_{11} x_{1} + a_{12} x_{2} + \cdots + a_{1n} x_{n} = b_{1} \\
	& & a_{21} x_{1} + a_{22} x_{2} + \cdots + a_{2n} x_{n} = b_{2} \\
	& & \vdots \\
	& & a_{m1} x_{1} + a_{m2} x_{2} + \cdots + a_{mn} x_{n} = b_{m} \\
	& & c_{1} x_{1} + c_{2} x_{2} + c_{3} x_{3} + \cdots + c_{n} x_{n} - z = 0 \\
	& & x_{1}, \ldots, x_{n} \geq 0
	\end{array}
	\right.
	\end{equation}
	
	Sans perte de généralité, nous pouvons supposer que les variables $x_{1}, \ldots, x_{m}$ sont les variables de base. Le problème \eqref{pro3.2} peut donc être représenté par le tableau suivant, dit tableau du simplexe :
	
	
	\begin{center}
		\begin{tabular}{c|ccccccccc|c|c}
			\toprule
			Variables de base & \multicolumn{4}{c}{$\overbrace{\phantom{xxxxxxxxxxxxxxxxx}}^{\text{variable de base}}$} & \multicolumn{4}{c}{$\overbrace{\phantom{xxxxxxxxxxxxxxxxx}}^{\text{variable hors base}}$}& & & \\
			& $x_1$ & $x_2$ & $\cdots$ & $x_m$ & $x_{m+1}$ & $x_{m+2}$ & $\cdots$ & $x_n$ & $z$ & Termes de droite & $R$ \\
			\midrule
			$x_1$ & $a_{11}$ & $a_{12}$ & $\cdots$ & $a_{1m}$ & $a_{1m+1}$ & $a_{1m+2}$ & $\cdots$ & $a_{1n}$ & 0 & $b_1$ &\\
			$x_2$ & $a_{21}$ & $a_{22}$ & $\cdots$ & $a_{2m}$ & $a_{2m+1}$ & $a_{2m+2}$ & $\cdots$ & $a_{2n}$ & 0 & $b_2$ &\\
			$\vdots$ & $\vdots$ & $\vdots$ & $\ddots$ & $\vdots$ & $\vdots$ & $\vdots$ & $\ddots$ & $\vdots$ & $\vdots$ & $\vdots$& \\
			$x_m$ & $a_{m1}$ & $a_{m2}$ & $\cdots$ & $a_{mm}$ & $a_{mm+1}$ & $a_{mm+2}$ & $\cdots$ & $a_{mn}$ & 0 & $b_m$ &\\
			\hline
			$z$ & $c_1$ & $c_2$ & $\cdots$ & $c_m$ & $c_{m+1}$ & $c_{m+2}$ & $\cdots$ & $c_n$ & $-1$ & 0& \\
			\hline
		\end{tabular}
	\end{center}
	
	\subsection{Problème de maximisation}
	Pour bien visualiser la similitude qui existe entre l'approche algébrique et celle des tableaux du simplexe, nous allons résoudre à nouveau le problème de production.\\
	
	Le problème, après l'introduction des variables d'écart, était le suivant :
	
	\begin{itemize}
		\item Maximiser la fonction économique :
		\[
		f(x_{1}, x_{2}, e_{1}, e_{2}, e_{3}) = 6x_{1} + 4x_{2} + 0e_{1} + 0e_{2} + 0e_{3}
		\]
		\item Sous les contraintes :
		\[
		\left\{
		\begin{aligned}
		3x_{1} + 9x_{2} + e_{1} &= 81 \\
		4x_{1} + 5x_{2} + e_{2} &= 55 \\
		2x_{1} + x_{2} + e_{3} &= 20 \\
		x_{1}, x_{2}, e_{1}, e_{2}, e_{3} &\geq 0
		\end{aligned}
		\right.
		\]
	\end{itemize}
	
	Que l'on peut écrire sous la forme :
	
	\[
	\left\{
	\begin{array}{rcl}
	3x_1 + 9x_2 + e_1 + 0e_2 + 0e_3 & = & 81 \\
	4x_1 + 5x_2 + 0e_1 + e_2 + 0e_3 & = & 55 \\
	2x_1 + x_2 + 0e_1 + 0e_2 + e_3 & = & 20 \\
	x_1, x_2, e_1, e_2, e_3 & \geq & 0
	\end{array}
	\right.
	\]
	
	%	\subsection*{Remarque 3.2}
	\begin{remarque}
		Cette présentation sous forme d'un tableau de nombres est appelée matrice du problème de programmation linéaire. Toutes les propriétés que vous connaissez sur les calculs à propos des matrices peuvent être alors utilisées.
	\end{remarque}
	
	
	\textbf{ÉTAPE no 1 :} Matrice du programme de base no 1 :
	
	\begin{center}
		\begin{tabular}{c|cccccc|c|c}
			\toprule
			Variables de base & $x_{1}$ & $x_{2}$ & $e_{1}$ & $e_{2}$ & $e_{3}$ & $z$ & $b$&R \\
			\midrule
			$e_{1}$ & 3 & 9 & 1 & 0 & 0 & 0 & 81& \\
			$e_{2}$ & 4 & 5 & 0 & 1 & 0 & 0 & 55 &\\
			$e_{3}$ & 2 & 1 & 0 & 0 & 1 & 0 & 20 &\\
			\hline
			$z$ & 6 & 4 & 0 & 0 & 0 & -1 & 0 &\\
			\bottomrule
		\end{tabular}
	\end{center}
	
	\textbf{ÉTAPE no 2 :} Matrice du programme de base no 2 :
	
	
	Pour augmenter le plus rapidement possible la valeur de la fonction économique, il faut donner une valeur positive à $ x_{1} $ puisque c'est qui a le plus grand coefficient positif sur la dernière ligne du tableau. Quelle est la plus grande valeur que l'on peut attribuer à $ x_{1} $ ?	Pour la déterminer, on prend les rapports des éléments de la colonne de b sur les éléments de la colonne des $ x_{1} $. On trouve alors :
	\begin{center}
		\begin{tabular}{c|>{\columncolor{gray!25}}cccccc|c|c}
			\toprule
			Variables de base & \textcolor{blue}{$x_{1}$} & $x_{2}$ & $e_{1}$ & $e_{2}$ & $e_{3}$ & $z$ & $b$&R \\
			\midrule
			$e_{1}$ & 3 & 9 & 1 & 0 & 0 & 0 & 81 & $\dfrac{81}{3}=27$ \\
			$e_{2}$ & 4 & 5 & 0 & 1 & 0 & 0 & 55 & $\dfrac{55}{4}$ \\
			\rowcolor{gray!25}\textcolor{red}{$e_{3}$} & 2 & 1 & 0 & 0 & 1 & 0 & 20 & $\dfrac{20}{2}=10$ \\
			\hline
			$z$ & 6 & 4 & 0 & 0 & 0 & -1 & 0 & \\
			\bottomrule
		\end{tabular}
	\end{center}
	
	Le plus petit de ces rapports est $ 10 $ et c'est la plus grande valeur qu'on peut attribuer à $ x_{1} $. Nous utiliserons donc l'élément de la troisième ligne première colonne; encadrons-le pour le mettre en évidence. Nous l'appellerons pivot de la 1ère étape.
	\begin{center}
		\begin{tabular}{c|>{\columncolor{gray!25}}cccccc|c|c}
			\toprule
			Variables de base & \textcolor{blue}{$x_{1}$} & $x_{2}$ & $e_{1}$ & $e_{2}$ & $e_{3}$ & $z$ & $b$ &R\\
			\midrule
			$e_{1}$ & 3 & 9 & 1 & 0 & 0 & 0 & 81 & $\dfrac{81}{3}=27$ \\
			$e_{2}$ & 4 & 5 & 0 & 1 & 0 & 0 & 55 & $\dfrac{55}{4}$ \\
			\rowcolor{gray!25}\textcolor{red}{$e_{3}$} & \tikz[baseline=(char.base)]{\node[draw=red,circle,very thick,inner sep=2pt] (char) {2};} & 1 & 0 & 0 & 1 & 0 & 20 & $\dfrac{20}{2}=10$ \\
			\hline
			$z$ & 6 & 4 & 0 & 0 & 0 & -1 & 0 & \\
			\bottomrule
		\end{tabular}
	\end{center}
	
	Nous devons maintenant annuler les autres éléments de la première colonne. Pour ce faire, nous effectuons des opérations sur les lignes. En premier lieu, nous allons rendre le pivot unitaire en divisant les termes de la troisième ligne par 2 :
	
	\[
	L_{3} \leftarrow \dfrac{1}{2} L_{3}
	\]
	\begin{center}
		\begin{tabular}{c|cccccc|c|c}
			\toprule
			Variables de base & $x_{1}$ & $x_{2}$ & $e_{1}$ & $e_{2}$ & $e_{3}$ & $z$ & $b$&R \\
			\midrule
			$e_{1}$ & 3 & 9 & 1 & 0 & 0 & 0 & 81& \\
			$e_{2}$ & 4 & 5 & 0 & 1 & 0 & 0 & 55 &\\
			$x_{1}$ & 1 & $\dfrac{1}{2}$ & 0 & 0 & $\dfrac{1}{2}$ & 0 & 10 &\\
			\hline
			$z$ & 6 & 4 & 0 & 0 & 0 & -1 & 0& \\
			\bottomrule
		\end{tabular}
	\end{center}
	
	En deuxième lieu, on soustrait 4 fois la troisième ligne à la deuxième, 3 fois la troisième à la
	première et finalement 6 fois la troisième à la dernière :
	\[
	\begin{aligned}
	L_{1} &\leftarrow L_{1} - 3L_{3} \\
	L_{2} &\leftarrow L_{2} - 4L_{3} \\
	L_{4} &\leftarrow L_{4} - 6L_{3}
	\end{aligned}
	\]
	
	\begin{center}
		\begin{tabular}{c|cccccc|c|c}
			\toprule
			Variables de base & $x_{1}$ & $x_{2}$ & $e_{1}$ & $e_{2}$ & $e_{3}$ & $z$ & $b$&R \\
			\midrule
			$e_{1}$ & 0 & $\dfrac{15}{2}$ & 1 & 0 & $-\dfrac{3}{2}$ & 0 & 51& \\
			$e_{2}$ & 0 & 3 & 0 & 1 & -2 & 0 & 15& \\
			$x_{1}$ & 1 & $\dfrac{1}{2}$ & 0 & 0 & $\dfrac{1}{2}$ & 0 & 10& \\
			\hline
			$z$ & 0 & 1 & 0 & 0 & -3 & -1 & -60& \\
			\bottomrule
		\end{tabular}
	\end{center}
	Ce tableau représente alors le système d'équations que nous avions obtenu après la deuxième étape de la résolution algébrique \eqref{equ2.7}. On constate également que la valeur obtenue en bas à droite du tableau correspond à l'opposé de la valeur numérique obtenue dans la fonction économique :
	\[
	z = 60 + x_{2} - 3e_{3}
	\]
	On constate que l'on peut encore augmenter la valeur de la fonction économique en introduisant une valeur pour $ x_{2} $.
	
	
	\textbf{ÉTAPE no 3 :} Matrice du programme de base no 3 :
	
	Nous avons
	\begin{center}
		\begin{tabular}{c|c>{\columncolor{gray!25}}ccccc|c|c}
			\toprule
			Variables de base & $x_{1}$ & \textcolor{blue}{$x_{2}$} & $e_{1}$ & $e_{2}$ & $e_{3}$ & $z$ & $b$&R \\
			\midrule
			$e_{1}$ & 0 & $\dfrac{15}{2}$ & 1 & 0 & $-\dfrac{3}{2}$ & 0 & 51 & $\dfrac{51}{\frac{15}{2}}= \dfrac{102}{15}$ \\
			\rowcolor{gray!25}\textcolor{red}{$e_{2}$} & 0 & \tikz[baseline=(char.base)]{\node[draw=red,circle,very thick,inner sep=2pt] (char) {3};} & 0 & 1 & -2 & 0 & 15 & $\dfrac{15}{3}=5$ \\
			$x_{1}$ & 1 & $\dfrac{1}{2}$ & 0 & 0 & $\dfrac{1}{2}$ & 0 & 10 & $\dfrac{10}{\dfrac{1}{2}}=20$ \\
			\hline
			$z$ & 0 & 1 & 0 & 0 & -3 & -1 & -60 & \\
			\bottomrule
		\end{tabular}
	\end{center}
	
	Le plus petit rapport étant sur la deuxième ligne, notre pivot sera l'élément de la deuxième ligne, deuxième colonne. Rendons-le unitaire en le divisant par 3 :
	\[
	L_{2} \leftarrow \dfrac{1}{3} L_{2}
	\]
	
	\begin{center}
		\begin{tabular}{c|ccccc|c|c}
			\toprule
			Variables de base & $x_{1}$ & $x_{2}$ & $e_{1}$ & $e_{2}$ & $e_{3}$ & $z$ & $b$ \\
			\midrule
			$e_{1}$ & 0 & $\dfrac{15}{2}$ & 1 & 0 & $-\dfrac{3}{2}$ & 0 & 51 \\
			$x_{2}$ & 0 & 1 & 0 & $\dfrac{1}{3}$ & $-\dfrac{2}{3}$ & 0 & 5 \\
			$x_{1}$ & 1 & $\dfrac{1}{2}$ & 0 & 0 & $\dfrac{1}{2}$ & 0 & 10 \\
			\hline
			$z$ & 0 & 1 & 0 & 0 & -3 & -1 & -60 \\
			\bottomrule
		\end{tabular}
	\end{center}
	
	Effectuons à nouveau les opérations sur les lignes afin d'annuler les autres éléments de la 2ème colonne :
	\[
	\begin{aligned}
	L_{1} &\leftarrow L_{1} - \dfrac{15}{2} L_{2} \\
	L_{3} &\leftarrow L_{3} - \dfrac{1}{2} L_{2} \\
	L_{4} &\leftarrow L_{4} - 1 L_{2}
	\end{aligned}
	\]
	
	\begin{center}
		\begin{tabular}{c|ccccc|c|c}
			\toprule
			Variables de base & $x_{1}$ & $x_{2}$ & $e_{1}$ & $e_{2}$ & $e_{3}$ & $z$ & $b$ \\
			\midrule
			$e_{1}$ & 0 & 0 & 1 & $-\dfrac{5}{2}$ & $\dfrac{7}{2}$ & 0 & $\dfrac{27}{2}$ \\
			$x_{2}$ & 0 & 1 & 0 & $\dfrac{1}{3}$ & $-\dfrac{2}{3}$ & 0 & 5 \\
			$x_{1}$ & 1 & 0 & 0 & $-\dfrac{1}{6}$ & $\dfrac{1}{3}$ & 0 & $\dfrac{15}{2}$ \\
			\hline
			$z$ & 0 & 0 & 0 & $-\dfrac{1}{3}$ & $-\dfrac{7}{3}$ & -1 & -65 \\
			\bottomrule
		\end{tabular}
	\end{center}
	
	Il n'y a plus de coefficients positifs sur la dernière ligne du tableau et il n'est donc plus possible d'augmenter la valeur de la fonction économique. La solution optimale est donc
	\[
	(x_{1}, x_{2}, e_{1}, e_{2}, e_{3})^{t} = \left( \dfrac{15}{2}, 5, \dfrac{27}{2}, 0, 0 \right)^{t}
	\]
	qui correspond à un bénéfice maximum de $ z = 65 $ produits en produisant $ \dfrac{15}{2} $ unités du produit $ P_{1} $ et 5 unités du produit $ P_{2} $.
	
	\subsection{Problème de minimisation}
	
	Nous avons vu dans que tout programme de minimisation peut se transformer en un problème de maximisation.
	Pour la résolution d'un problème de minimisation par l'algorithme du simplexe, nous devons adapter le critère d'optimalité. Ainsi pour que la solution de base soit optimale il faut que les coûts marginaux de toutes les variables hors base soient positifs ou nuls.
	Pour mieux comprendre ce cas, je vous propose l'exemple suivant :
	
	Considérons le programme de minimisation sous forme canonique suivant :
	\[
	\left\{
	\begin{array}{r}
	\displaystyle\min -x_{1} - 2x_{2} \\
	\text{s.t.} \quad -3x_{1} + 2x_{2} \leq 2 \\
	-x_{1} + 2x_{2} \leq 4 \\
	x_{1} + x_{2} \leq 5 \\
	x_{1} \geq 0, \quad x_{2} \geq 0
	\end{array}
	\right.
	\]
	
	Après introduction des variables d'écart :
	\[
	z = -x_{1} - 2x_{2} + 0e_{1} + 0e_{2} + 0e_{3}
	\]
	Sous les contraintes :
	\[
	\left\{
	\begin{aligned}
	-3x_{1} + 2x_{2} + e_{1} &= 2 \\
	-x_{1} + 2x_{2} + e_{2} &= 4 \\
	x_{1} + x_{2} + e_{3} &= 5 \\
	x_{1} + 2x_{2} - z &= 0 \\
	x_{1}, x_{2}, e_{1}, e_{2}, e_{3} &\geq 0
	\end{aligned}
	\right.
	\]
	
	\textbf{ÉTAPE no 1 :} Matrice du programme de base no 1 : On a donc le tableau du simplexe suivant
	où les variables $ e_{1} $, $ e_{2} $ et $ e_{3} $ jouent le rôle des variables de base :
	
	\begin{center}
		\begin{tabular}{c|ccccc|c|c}
			\toprule
			Variables de base & $x_{1}$ & $x_{2}$ & $e_{1}$ & $e_{2}$ & $e_{3}$ & $-z$ & $b$ \\
			\midrule
			$e_{1}$ & -3 & 2 & 1 & 0 & 0 & 0 & 2 \\
			$e_{2}$ & -1 & 2 & 0 & 1 & 0 & 0 & 4 \\
			$e_{3}$ & 1 & 1 & 0 & 0 & 1 & 0 & 5 \\
			\hline
			$-z$ & 1 & 2 & 0 & 0 & 0 & 1 & 0 \\
			\bottomrule
		\end{tabular}
	\end{center}
	Cette base n'est pas optimale puisque les coûts relatifs des variables hors base dans la dernière
	ligne ne sont pas positifs ou nuls.
	\[\displaystyle\min\left\lbrace\dfrac{b_{i}}{a_{i2}}, \, a_{i2}>0\right\rbrace   \]
	\textbf{ÉTAPE no 2 :}  Matrice du programme de base no 2 :
	
	Choisissons comme variable entrante dans la nouvelle base la variable hors base ayant le plus petit coût négatif négatif, c'est à dire $ x_{2} $. cherchons la variable sortante en appliquant la règle
	du plus petit rapport :
	
	\begin{center}
		\begin{tabular}{c|c|>{\columncolor{gray!25}}c|c|c|c|c|c|c}
			\hline
			Variables de base & $x_{1}$ & \textcolor{blue}{$x_{2}$} & $e_{1}$ & $e_{2}$ & $e_{3}$ & $-z$ & $b$ & R \\
			\hline
			\rowcolor{gray!25}
			\textcolor{red}{$e_{1}$} & $-3$ & \tikz[baseline=(char.base)]{\node[draw=red,circle,very thick,inner sep=2pt] (char) {2};} & 1 & 0 & 0 & 0 & 2 & $\dfrac{2}{2}=1$ \\
			\hline
			$e_{2}$ & $-1$ & 2 & 0 & 1 & 0 & 0 & 4 & $\dfrac{4}{2}=2$ \\
			\hline
			$e_{3}$ & 1 & 1 & 0 & 0 & 1 & 0 & 5 & $\dfrac{5}{1}=5$ \\
			\hline
			$-z$ & 1 & 2 & 0 & 0 & 0 & 1 & 0 & \\
			\hline
		\end{tabular}
	\end{center}
	
	
	Le plus petit de ces rapports est 1 et c'est la plus petite valeur qu'on peut attribuer à $ x_{2} $. Nous utiliserons donc l'élément de la première ligne deuxième colonne, autrement, la variable $ e_{1} $ quitte la base et sera remplacée par la variable $ x_{2} $, ainsi on obtient la tableau suivant.
	\[
	L_{1} \leftarrow \dfrac{1}{2} L_{1}
	\]
	
	\begin{center}
		\begin{tabular}{c|ccccc|c|c}
			\toprule
			Variables de base & $x_{1}$ & $x_{2}$ & $e_{1}$ & $e_{2}$ & $e_{3}$ & $-z$ & $b$ \\
			\midrule
			$x_{2}$ & $-\dfrac{3}{2}$ & 1 & $\dfrac{1}{2}$ & 0 & 0 & 0 & 1 \\
			$e_{2}$ & -1 & 2 & 0 & 1 & 0 & 0 & 4 \\
			$e_{3}$ & 1 & 1 & 0 & 0 & 1 & 0 & 5 \\
			\hline
			$-z$ & 1 & 2 & 0 & 0 & 0 & 1 & 0 \\
			\bottomrule
		\end{tabular}
	\end{center}
	
	Nous devons maintenant annuler les autres éléments de la deuxième colonne.
	\[
	\begin{aligned}
	L_{2} &\leftarrow L_{2} - 2 L_{1} \\
	L_{3} &\leftarrow L_{3} - 1 L_{1} \\
	L_{4} &\leftarrow L_{4} - 2 L_{1}
	\end{aligned}
	\]
	
	\begin{center}
		\begin{tabular}{c|ccccc|c|c}
			\toprule
			Variables de base & $x_{1}$ & $x_{2}$ & $e_{1}$ & $e_{2}$ & $e_{3}$ & $-z$ & $b$ \\
			\midrule
			$x_{2}$ & $-\dfrac{3}{2}$ & 1 & $\dfrac{1}{2}$ & 0 & 0 & 0 & 1 \\
			$e_{2}$ & 2 & 0 & -1 & 1 & 0 & 0 & 2 \\
			$e_{3}$ & $\dfrac{5}{2}$ & 0 & $-\dfrac{1}{2}$ & 0 & 1 & 0 & 4 \\
			$-z$ & 4 & 0 & -1 & 0 & 0 & 1 & -2 \\
			\bottomrule
		\end{tabular}
	\end{center}
	%==================================================================================================================================================================================================
	La solution de base correspondante n'est pas optimale car $c_{1}=-4<0$. la variable $x_{1}$ va donc entrer dans la base.\\
	
	\textbf{ÉTAPE No 3 :} Matrice du programme de base no 3 :\\
	
	\begin{center}
		\begin{tabular}{c|>{\columncolor{gray!25}}cccccc|c|c}
			\hline
			Variables de base & \textcolor{blue}{$x_{1}$} & $x_{2}$ & $e_{1}$ & $e_{2}$ & $e_{3}$ & $-z$ & $b$ & $R$ \\
			\hline
			$x_{2}$ & $\dfrac{-3}{2}$ & 1 & $\dfrac{1}{2}$ & 0 & 0 & 0 & 1 &  \\
			\rowcolor{gray!25}\textcolor{red}{$e_{2}$} & \tikz[baseline=(char.base)]{\node[draw=red,circle,very thick,inner sep=2pt] (char) {2};} & 0 & -1 & 1 & 0 & 0 & 2 & $\dfrac{2}{2}=1$ \\
			$e_{3}$ & $\dfrac{5}{2}$ & 0 & $-\dfrac{1}{2}$ & 0 & 1 & 0 & 4 & $\dfrac{4}{\dfrac{5}{2}}=\dfrac{8}{5}$ \\
			\hline
			$-z$ & -4 & 0 & 1 & 0 & 0 & 1 & 2 &  \\
			\hline
		\end{tabular}
	\end{center}
	
	Le plus petit rapport étant sur la deuxième ligne donc $e_{2}$ est une variable sortante, notre pivot sera l'élément de la deuxième ligne, première colonne. Rendons-le unitaire en le divisant par 2 :
	
	$L_{2} \leftarrow \dfrac{1}{2} L_{2} \quad$\begin{tabular}{c|cccccc|c|c}
		\hline
		Variables de base & $x_{1}$ & $x_{2}$ & $e_{1}$ & $e_{2}$ & $e_{3}$ & $-z$ & $b$ & $R$ \\
		\hline
		$x_{2}$ & $\dfrac{-3}{2}$ & 1 & $\dfrac{1}{2}$ & 0 & 0 & 0 & 1 &  \\
		$x_{1}$ & 1 & 0 & $-\dfrac{1}{2}$ & $\dfrac{1}{2}$ & 0 & 0 & 1 &  \\
		$e_{3}$ & $\dfrac{5}{2}$ & 0 & $-\dfrac{1}{2}$ & 0 & 1 & 0 & 4 &  \\
		\hline
		$-z$ & -4 & 0 & 1 & 0 & 0 & 1 & 2 &  \\
		\hline
	\end{tabular}
	
	Ainsi,
	
	\begin{center}
		\begin{tabular}{ll|llllll|l|l}
			\cline{2-10}
			\multirow{3}{*}{$L_{1} \leftarrow L_{1}+\dfrac{3}{2} L_{2}$} & Variables de base & $x_{1}$ & $x_{2}$ & $e_{1}$ & $e_{2}$ & $e_{3}$ & -z & $b$ & $R$ \\
			\cline{2-10}
			& \multirow{3}{*}{\begin{tabular}{l}
					$x_{2}$ \\
					$x_{1}$ \\
					$e_{3}$ \\
			\end{tabular}} & 0 & 1 & $-\dfrac{1}{4}$ & $\dfrac{3}{4}$ & 0 & 0 & $\dfrac{5}{2}$ & \multirow{3}{*}{} \\
			&  & 1 & 0 & $-\dfrac{1}{2}$ & $\dfrac{1}{2}$ & 0 & 0 & 1 &  \\
			$L_{3} \leftarrow L_{3}-\dfrac{5}{2} L_{2}$ &  & 0 & 0 & $\dfrac{3}{4}$ & $-\dfrac{5}{4}$ & 1 & 0 & $\dfrac{3}{2}$ &  \\
			\cline{2-10}
			$L_{4} \leftarrow L_{4}+4 L_{2}$ & $-z$ & 0 & 0 & -1 & 2 & 0 & 1 & 6 &  \\
			\cline{2-10}
		\end{tabular}
	\end{center}
	
	La solution de base obtenue a pour coût $z=-6$, cependant elle n'est pas optimale car $c_{3}=-1<0$, la variable $e_{1}$ va donc entrer dans la base.\\
	
	\textbf{ÉTAPE No 4 :} Matrice du programme de base no 4 :\\
	
	\begin{center}
		\begin{tabular}{c|cc>{\columncolor{gray!25}}cccc|c|c}
			\toprule
			Variables de base & $x_{1}$ & $x_{2}$ & \textcolor{blue}{$e_{1}$} & $e_{2}$ & $e_{3}$ & $-z$ & $b$ & $R$ \\
			\midrule
			$x_{2}$ & 0 & 1 & $-\dfrac{1}{4}$ & $\dfrac{3}{4}$ & 0 & 0 & $\dfrac{5}{2}$ &  \\
			$x_{1}$ & 1 & 0 & $-\dfrac{1}{2}$ & $\dfrac{1}{2}$ & 0 & 0 & 1 &  \\
			\rowcolor{gray!25}\textcolor{red}{$e_{3}$} & 0 & 0 & \tikz[baseline=(char.base)]{\node[draw=red,circle,very thick,inner sep=2pt] (char) {$ \dfrac{3}{4} $};} & $-\dfrac{5}{4}$ & 1 & 0 & $\dfrac{3}{2}$ & $\dfrac{\dfrac{3}{2}}{\dfrac{3}{4}}=2$ \\
			\hline
			$-z$ & 0 & 0 & -1 & 2 & 0 & 1 & 6 &  \\
			\hline
		\end{tabular}
	\end{center}
	
	Par suite, on obtient
	
	\begin{tabular}{cc|cccccc|c|c}
		\cline { 2 - 10 }
		&Variables de base & $x_{1}$ & $x_{2}$ & $e_{1}$ & $e_{2}$ & $e_{3}$ & $-z$ & $b$ & $R$ \\
		\cline { 2 - 10 }
		& $x_{2}$&0 & 1 & $-\dfrac{1}{4}$ & $\dfrac{3}{4}$ & 0 & 0 & $\dfrac{5}{2}$& \\
		& $x_{1}$& 1 & 0 & $-\dfrac{1}{2}$ & $\dfrac{1}{2}$ & 0 & 0 & 1& \\
		$L_{3} \leftarrow \dfrac{4}{3} L_{3}$&$e_{1}$ & 0 & 0 & 1 & $-\dfrac{5}{3}$ & $\dfrac{4}{3}$ & 0 & 2 &  \\
		\cline { 2 - 10 }
		&$-z$&0 & 0 & -1 & 2 & 0 & 1 & 6& \\
		\cline { 2 - 10 }
	\end{tabular}
	
	rendons la colonne associé a $e_{1}$ unitaire
	
	\begin{center}
		\begin{tabular}{cc|cccccc|c|c}
			\cline { 2 - 10 }
			& Variables de base & $x_{1}$ & $x_{2}$ & $e_{1}$ & $e_{2}$ & $e_{3}$ & $-z$ & $b$ & $R$ \\
			\cline { 2 - 10 }
			$L_{1} \leftarrow L_{1}+\dfrac{1}{4} L_{3}$& $x_{2}$ & 0 & 1 & 0 & $\dfrac{1}{3}$ & $\dfrac{1}{3}$ & 0 & 3 &  \\
			$L_{2} \leftarrow L_{2}+\dfrac{1}{2} L_{3}$ & $x_{1}$ & 1 & 0 & 0 & $-\dfrac{1}{3}$ & $\dfrac{2}{3}$ & 0 & 2 &  \\
			& $e_{1}$ & 0 & 0 & 1 & $-\dfrac{5}{3}$ & $\dfrac{4}{3}$ & 0 & 2 &  \\
			\cline { 2 - 10 }
			$L_{4} \leftarrow L_{4}+L_{3}$& $-z$ & 0 & 0 & 0 & $\dfrac{1}{3}$ & $\dfrac{4}{3}$ & 1 & 8 &  \\
			\cline { 2 - 10 }
		\end{tabular}
	\end{center}
	
	La nouvelle solution de base est
	
	$$
	\left(x_{1}, x_{2}, e_{1}, e_{2}, e_{3}\right)^{t}=(2,3,2,0,0)^{t}
	$$
	
	qui correspond à un coût minimum de $z=-8$, comme les coût réduit de toutes les variables sont positifs ou nuls, cette solution est optimale et l'algorithme du simplexe s'arrête.
	
	\section{Principes de la Méthode du Rectangle}
	
	Pour faciliter les calculs dans les tableaux du simplexe, on peut utiliser la méthode du rectangle :
	
	\begin{enumerate}
		\item Si une ligne contient un zéro (0) qui est à l'intersection avec la colonne pivot, cette ligne ne sera pas modifiée.
		
		\item Si une colonne contient un zéro (0) qui est à l'intersection avec la ligne pivot, cette colonne ne sera pas modifiée.
		
		\item Pour les cellules qui ne se trouvent pas sur la ligne et la colonne pivot, appliquer la méthode du rectangle :
		$$a' \leftarrow a - \frac{b \times c}{\text{pivot}}$$
		où $a'$ est la nouvelle valeur de $a$, $b$ l'élément situé sur la même ligne que $a$ mais dans la colonne du pivot, $c$ l'élément situé sur la même colonne que $a$ mais dans la ligne du pivot.
		
		\item Pour les éléments situés sur la ligne du pivot, ils seront divisés par le pivot.
		
		\item Pour les éléments situés sur la colonne du pivot, ceux situés en dessous et au dessus du pivot seront réduits à zéro (0).
	\end{enumerate}
	
	\subsection{Conclusion}
	L'algorithme du simplexe consiste à se déplacer d'un point extrême du polytope $K=\{x: A x= b, x \geq 0\}$ à un autre point extrême adjacent jusqu'à ce qu'on se rende compte que la solution réalisable actuelle est optimale ou que le problème n'est pas réalisable ou que le problème n'est pas borné supérieurement (ou inférieurement dans le cas de minimisation).
	
	\section{Méthodes des deux phases}
	Nous avons vu comment trouver une base de départ lorsque les contraintes du problème sont toutes de type $\leq$ : les variables d'écart constituent alors une base de départ naturelle, leurs colonnes formant a une permutation près une matrice identité.
	
	En pratique il est fréquent que le PL comporte également des contraintes de type $\geq$ ou de type $=$. Dans ce cas, nous n'avons plus de base de départ évidente, il faut en chercher une. Ce sera la première phase de la méthode du simplexe.
	
	La méthode du simplexe consiste donc en deux phases :\\
	
	\begin{description}
		\item[Phase I:] recherche d'une base de départ par la résolution par l'algorithme du simplexe du programme auxiliaire (définit plus loin) ou mise en évidence que le problème de départ n'admet aucune solution admissible.
		
		\item[Phase II :] application de l'algorithme du simplexe au problème de départ cette fois, en partant de la base mise en évidence a la fin de la phase I
	\end{description}
	
	%Exemple 3.1 
	\begin{exemple}
		Comment, par exemple trouver une solution de base réalisable initiale au problème suivant
		
		$$
		(P):\left\{\begin{aligned}
		\displaystyle\max_{x} \quad 2 x_{1}+3 x_{2} & \\
		\text { s.t. } \quad x_{1}+x_{2} & \leq 10, \\
		5 x_{1}+4 x_{2} & \geq 20 \\
		x_{1}, x_{2} & \geq 0 .
		\end{aligned}\right.
		$$
		
		On transforme ce problème sous la forme canonique $A x \leq b$ et ensuite on ajoute les variables d'écart.
		
		\begin{equation}\label{eq3.3}
		\left\{\begin{aligned}
		\max_{x} \quad 2 x_{1}+3 x_{2} & \\
		\text { s.t. } \quad x_{1}+x_{2}+e_{1} & =10, \\
		-5 x_{1}-4 x_{2}+e_{2} & =-20 \\
		x_{1}, x_{2}, e_{1}, e_{2} & \geq 0 .
		\end{aligned}\right.
		\end{equation}
		
		Clairement le point ( $0,0,10,-20$ ), n'est pas admissible.
	\end{exemple}
	
	\subsection{Recherche d'une base de départ : Phase I}
	L'objectif est de trouver une solution de base admissible qui servira de point de départ pour l'algorithme du simplexe. L'idée est de résoudre un problème intermédiaire de minimisation dont la solution fournira le point de départ de la méthode du simplexe. Ce problème intermédiaire porte le nom de Phase I du simplexe.
	
	
	\begin{colorframed}
		{\color{titleblue}\textbf{Démarche à suivre :}}
		
		\medskip
		
		On commence par écrire le {\color{keywordred}programme auxiliaire (PA)} associé au {\color{keywordred}problème de départ (P)} :
		
		\begin{enumerate}[label=\arabic*.]
			\item On s'assure que les {\color{vargreen}membres de droite $b_i$} sont tous {\color{keywordred}positifs ou nuls}, quitte à multiplier la contrainte correspondante par {\color{mathblue}$-1$} si son membre de droite est {\color{keywordred}négatif} (sinon on aboutirait à des {\color{keywordred}solutions de base non réalisables}).
			
			\item On introduit une {\color{vargreen}variable d'écart $e_i$} pour chaque contrainte $i$ de type {\color{mathblue}$\leq$}, affectée d’un signe {\color{mathblue}$+$} dans la contrainte réécrite, comme on le faisait jusqu'à présent.
			
			\item On introduit aussi une {\color{vargreen}variable d'écart $e_i$} pour chaque contrainte $i$ de type {\color{mathblue}$\geq$}, mais affectée d’un signe {\color{mathblue}$-$} dans la contrainte réécrite : l’écart qu’elle mesure est en effet {\color{keywordred}négatif ou nul}, mais la variable d’écart doit rester {\color{keywordred}positive ou nulle}.
			
			\item Le problème ne comporte plus que des contraintes d’{\color{keywordred}égalité}, mais seules les {\color{vargreen}variables d’écart associées aux contraintes initiales de type {\color{mathblue}$\leq$}} forment des colonnes de la matrice identité. Pour chacune des autres contraintes (c’est-à-dire les contraintes initiales {\color{mathblue}$\geq$} réécrites avec une variable d’écart affectée d’un signe {\color{mathblue}$-$}, et les contraintes initiales {\color{mathblue}$=$}), on introduit une {\color{vargreen}variable artificielle $t_i$} {\color{keywordred}positive ou nulle} de la même façon qu’on introduisait les variables d’écart.
			
			\item La {\color{keywordred}fonction objectif de (PA)} est alors la {\color{vargreen}somme des variables artificielles}. On cherche à la {\color{keywordred}minimiser} (en fait à l’{\color{keywordred}annuler}).
		\end{enumerate}
	\end{colorframed}
	
	\subsection{Illustration de la méthode sur un exemple}
	Dans l'exemple ci-dessus, il s'agit d'introduire deux variables d'écart $e_{1}$ et $e_{2}$ et une variable artificielle $t_{1}$ et de considérer le problème de minimisation associé
	\[(P): \left\{
	\begin{array}{ll}
	\underset{x}{\displaystyle\max} & 2x_1 + 3x_2 \\
	\text{s.t.} & x_1 + x_2 \leq 10, \\
	& 5x_1 + 4x_2 \geq 20 \\
	& x_1, x_2 \geq 0.
	\end{array}
	\right.
	\overset{\text{Var. d'écart}}{\Longrightarrow}
	\begin{array}{c}
	\left\{
	\begin{array}{ll}
	\underset{x}{\displaystyle\max} & 2x_1 + 3x_2 \\
	\text{s.t.} & x_1 + x_2 + e_1 = 10, \\
	& 5x_1 + 4x_2 - e_2 = 20 \\
	& x_1, x_2, e_1, e_2 \geq 0.
	\end{array}
	\right.
	\\
	\\
	\Big\Downarrow \substack{\textcolor{blue}{\text{Var. artificiel}} \\ \text{+ new obj}}
	\\
	\\
	(PA)\left\{
	\begin{array}{l}
	\displaystyle\min \quad z_t = \textcolor{red}{t_1(20-5x_1-4x_2+e_2)} \\
	\begin{array}{ll}
	\text{s.t.} & x_1 + x_2 + e_1 = 10, \\
	& 5x_1 + 4x_2 - e_2 + t_1 = 20 \\
	& x_1, x_2, e_1, e_2, t_1 \geq 0.
	\end{array}
	\end{array}
	\right.
	\end{array}\]
	
	
	\begin{description}
		\item[Phase I]\quad\\
		Le problème (PA) admet toujours une solution admissible, il suffit de prendre le point ( $x_{1}, x_{2}, e_{1}, e_{2}, t_{1}$ ) = ( $0,0,10,0,20$ ).
		
		On forme le tableau initial du simplexe pour ce problème.
		
		\begin{center}
			\begin{tabular}{|c|cccccc|c|c|}
				\hline
				Variables de base & $x_{1}$ & $x_{2}$ & $e_{1}$ & $e_{2}$ & $t_{1}$ & $-z_{t}$ & $b$ & $R$ \\
				\hline
				$e_{1}$ & 1 & 1 & 1 & 0 & 0 & 0 & 10 &  \\
				$t_{1}$ & 5 & 4 & 0 & -1 & 1 & 0 & 20 &  \\
				\hline
				$-z_{t}$ & 0 & 0 & 0 & 0 & 1 & 1 & 0 &  \\
				\hline
			\end{tabular}
		\end{center}
		
		Pour que $t_{1}$ joue le rôle d'une variable de base, il faut faire apparaître un 0 sur la dernière ligne de la colonne de $t_{1}$, pas suite, on transforme le tableau en soustrayant la deuxième ligne de la dernière. Nous obtenons le tableau suivant :
		
		\begin{center}
			\begin{tabular}{|c|>{\columncolor{gray!25}}cccccc|c|c|}
				\hline
				Variables de base & \textcolor{blue}{$x_{1}$} & $x_{2}$ & $e_{1}$ & $e_{2}$ & $t_{1}$ & $-z_{t}$ & $b$ & $R$ \\
				\hline
				$e_{1}$ & 1 & 1 & 1 & 0 & 0 & 0 & 10 & $\dfrac{10}{1}$ \\
				\rowcolor{gray!25}\textcolor{red}{$t_{1}$} & \tikz[baseline=(char.base)]{\node[draw=red,circle,very thick,inner sep=2pt] (char) {5};} & 4 & 0 & -1 & 1 & 0 & 20 & $\dfrac{20}{5}=4$ \\
				\hline
				$-z_{t}$ & -5 & -4 & 0 & 1 & 0 & 1 & -20 &  \\
				\hline
			\end{tabular}
		\end{center}
		
		Cette solution de base n'étant pas optimale, on poursuit avec la méthode standard du simplexe. Dans notre cas, on choisit la variable entrante $x_{1}$ et la variable sortante $t_{1}$, on obtient le tableau suivant
		
		\begin{center}
			\begin{tabular}{|c|cccccc|c|c|}
				\hline
				Variables de base & $x_{1}$ & $x_{2}$ & $e_{1}$ & $e_{2}$ & $t_{1}$ & $-z_{t}$ & $b$ & $R$ \\
				\hline
				$e_{1}$ & 0 & $\dfrac{1}{5}$ & 1 & $\dfrac{1}{5}$ & $-\dfrac{1}{5}$ & 0 & 6 &  \\
				$x_{1}$ & 1 & $\dfrac{4}{5}$ & 0 & $-\dfrac{1}{5}$ & $\dfrac{1}{5}$ & 0 & 4 &  \\
				\hline
				$-z_{t}$ & 0 & 0 & 0 & 0 & 1 & 1 & 0 &  \\
				\hline
			\end{tabular}
		\end{center}
		
		On ne peut améliorer ce résultat car la fonction objective $z=t_{1}=0$.\\
		Nous avons obtenu la solution optimale de la Phase $I(4,0,6,0,0)$ avec la base $\left\{e_{1}, x_{1}\right\}$. De plus, la variable $t_{1}=0$ devient inutile car hors-base. Ainsi, on peut enlever la colonne correspondante à $t_{1}$.
		
		Donc $(4,0,6,0,0)$ constitue une solution de base initiale du problème $(P)$, ce qui va nous permettre de construire le tableau du simplexe de (P) pour démarrer la phase II.
		
		\item[Phase II]\quad\\
		Après avoir terminer la phase I avec une solution optimale de valeur optimale nulle et où les variables artificielles sont toutes des variables hors base. Pour commencer la première itération de la phase II, il faudra remplacer les coefficients de la fonction objectif $z_{t}$ par ceux de la fonction objectif original $\mathcal{z}$, puis recalculer les coûts relatifs.
		
		\begin{center}
			\begin{tabular}{|c|ccccc|c|c|}
				\hline
				Variables de base & $x_{1}$ & $x_{2}$ & $e_{1}$ & $e_{2}$ & $z$ & $b$ & $R$ \\
				\hline
				$e_{1}$ & 0 & $\dfrac{1}{5}$ & 1 & $\dfrac{1}{5}$ & 0 & 6 &  \\
				$x_{1}$ & 1 & $\dfrac{4}{5}$ & 0 & $-\dfrac{1}{5}$ & 0 & 4 &  \\
				\hline
				$z$ & 2 & 3 & 0 & 0 & -1 & 0 &  \\
				\hline
			\end{tabular}
		\end{center}
		
		$x_{1}$ étant une variable de base nous devons remplacer 2 par 0 dans la dernière ligne du tableau par l'opération suivante $L_{3} \leftarrow L_{3}-2 L_{2}$
		
		\begin{center}
			\begin{tabular}{|c|c>{\columncolor{gray!25}}cccc|c|c|}
				\hline
				Variables de base & $x_{1}$ & \textcolor{blue}{$x_{2}$} & $e_{1}$ & $e_{2}$ & $z$ & $b$ & $R$ \\
				\hline
				$e_{1}$ & 0 & $\dfrac{1}{5}$ & 1 & $\dfrac{1}{5}$ & 0 & 6 & $6 \times 5=30$ \\
				\rowcolor{gray!25}\textcolor{red}{$x_{1}$} & 1 & \tikz[baseline=(char.base)]{\node[draw=red,circle,very thick,inner sep=2pt] (char) {$ \dfrac{4}{5} $};} & 0 & $-\dfrac{1}{5}$ & 0 & 4 & $4 \times \dfrac{5}{4}=5$ \\
				\hline
				$z$ & 0 & $\dfrac{7}{5}$ & 0 & $\dfrac{2}{5}$ & -1 & -8 &  \\
				\hline
			\end{tabular}
		\end{center}
		
		La variable $x_{2}$ entre dans la base et la variable $x_{1}$ sort de la base.
		
		\begin{center}
			\begin{tabular}{|c|c>{\columncolor{gray!25}}cccc|c|c|}
				\hline
				Variables de base & $x_{1}$ & \textcolor{blue}{$x_{2}$} & $e_{1}$ & $e_{2}$ & $z$ & $b$ & $R$ \\
				\hline
				$e_{1}$ & 0 & $\dfrac{1}{5}$ & 1 & $\dfrac{1}{5}$ & 0 & 6 & $6 \times 5=30$ \\
				\rowcolor{gray!25}\textcolor{red}{$x_{1}$} & $\dfrac{5}{4}$ & 1 & 0 & $-\dfrac{1}{5}$ & 0 & 4 & $4 \times \dfrac{5}{4}=5$ \\
				\hline
				$z$ & 0 & $\dfrac{7}{5}$ & 0 & $\dfrac{2}{5}$ & -1 & -8 &  \\
				\hline
			\end{tabular}
		\end{center}
		
		On poursuit avec la méthode standard du simplexe
		
		\begin{center}
			\begin{tabular}{|c|rrr>{\columncolor{gray!25}}rr|c|c|}
				\hline
				Variables de base & $x_{1}$ & $x_{2}$ & $e_{1}$ & \textcolor{blue}{$e_{2}$} & $z$ & $b$ & $R$ \\
				\hline
				\rowcolor{gray!25}\textcolor{red}{$e_{1}$} & $-\dfrac{1}{4}$ & 0 & 1 & $\dfrac{1}{4}$ & 0 & 5 & $5 \times 4=20$ \\
				$x_{2}$ & $\dfrac{5}{4}$ & 1 & 0 & $-\dfrac{1}{4}$ & 0 & 5 &  \\
				\hline
				$z$ & $-\dfrac{7}{5}$ & 0 & 0 & $\dfrac{3}{4}$ & -1 & -15 &  \\
				\hline
			\end{tabular}
		\end{center}
		
		La variable $e_{2}$ entre dans la base et la variable $e_{1}$ sort de la base.
		
		\begin{center}
			\begin{tabular}{|c|ccccc|c|c|}
				\hline
				Variables de base & $x_{1}$ & $x_{2}$ & $e_{1}$ & $e_{2}$ & $z$ & $b$ & $R$ \\
				\hline
				$e_{2}$ & -1 & 0 & 4 & 1 & 0 & 20 &  \\
				$x_{2}$ & 1 & 1 & 1 & 0 & 0 & 10 &  \\
				\hline
				$z$ & -1 & 0 & -3 & 0 & -1 & -30 &  \\
				\hline
			\end{tabular}
		\end{center}
		
		L'algorithme se termine à cette étape. La solution optimale est ( $0,10,0,20$ ) avec $z=30$. Ceci est bien conforme au graphique ci-dessous.
	\end{description}
	
	\section{La méthode de "The big $M$ "}
	L'idée ici est de remplacer le programme $ ( P ) $ par un autre qui a la même solution. Pratiquement il s'agit de faire deux choses :
	
	\begin{figure}[h]
		\begin{center}
			\includegraphics[scale=0.5]{2025_10_05_84c89f9c4234670b6556g-06}
			%		\captionsetup{labelformat=empty}
			\caption{Le graphe associé à l'exemple \ref{eq3.3}}
		\end{center}
	\end{figure}
	
	\begin{itemize}
		\item créer la base artificiellement en ajoutant les colonnes manquantes de la matrice identité, cela revient à introduire des variables supplémentaires, appelées "variables artificielles",
		\item modifier la fonction objective de sorte que les variables artificielles quittent la base au fur et à mesure des itérations.\\
		
	\end{itemize}

Pour le méthode du grand $M$ on pénalise les variables artificielles dans la fonction objective avec un coefficient négatif très grand noté $-M$ pour un problème de maximisation ( $+M$ pour un problème de minimisation). Ainsi, la maximisation de la fonction objective passe nécessairement par forcer les variables artificielles à quitter la base. Si à l'optimum, certaines variables artificielles sont dans la base, avec une valeur non nulle, alors le programme n'à pas de solution réalisable.\\
	
	Pour l'exemple donné plus haut, le programme associé, noté $M$, est le suivant.
	
	$$
	(P A)\left\{\begin{aligned}
	\displaystyle\max \quad z=2 x_{1}+3 x_{2}-M t_{1} & \\
	\text { s.t. } \quad x_{1}+x_{2}+e_{1} & =10, \\
	5 x_{1}+4 x_{2}-e_{2}+t_{1} & =20 \\
	x_{1}, x_{2}, e_{1}, e_{2}, t_{1} & \geq 0 .
	\end{aligned}\right.
	$$
	
	Le tableau initial du simplexe est :
	
	\begin{center}
		\begin{tabular}{|c|cccccc|c|c|}
			\hline
			Variables de base & $x_{1}$ & $x_{2}$ & $e_{1}$ & $e_{2}$ & $t_{1}$ & $z$ & $b$ & $R$ \\
			\hline
			$e_{1}$ & 1 & 1 & 1 & 0 & 0 & 0 & 10 &  \\
			$t_{1}$ & 5 & 4 & 0 & -1 & 1 & 0 & 20 &  \\
			\hline
			$z$ & 2 & 3 & 0 & 0 & $-M$ & -1 & 0 &  \\
			\hline
		\end{tabular}
	\end{center}
	
	Pour que $t_{1}$ joue le rôle d'une variable de base, il faut faire apparaître un 0 sur la dernière ligne de la colonne de $t_{1}$, pas suite, on transforme le tableau en soustrayant $M$ fois la deuxième ligne de la dernière. Nous obtenons le tableau suivant:
	
	\begin{center}
		\begin{tabular}{|c|>{\columncolor{gray!25}}cccccc|c|c|}
			\hline
			Variables de base & \textcolor{blue}{$x_{1}$} & $x_{2}$ & $e_{1}$ & $e_{2}$ & $t_{1}$ & $z$ & $b$ & $R$ \\
			\hline
			$e_{1}$ & 1 & 1 & 1 & 0 & 0 & 0 & 10 & $\dfrac{10}{1}$ \\
			\rowcolor{gray!25}\textcolor{red}{$t_{1}$} & \tikz[baseline=(char.base)]{\node[draw=red,circle,very thick,inner sep=2pt] (char) {5};} & 4 & 0 & -1 & 1 & 0 & 20 & $\dfrac{20}{5}=4$ \\
			\hline
			$z$ & $2+5 M$ & $3+4 M$ & 0 & $-M$ & 0 & -1 & $20 M$ &  \\
			\hline
		\end{tabular}
	\end{center}
	
	La variable artificiel $t_{1}$ est toujours dans la base, nous devons effectuer une autre itération
	
	\begin{center}
		\begin{tabular}{|c|cccccc|c|c|}
			\hline
			Variables de base & $x_{1}$ & $x_{2}$ & $e_{1}$ & $e_{2}$ & $t_{1}$ & $z$ & $b$ & $R$ \\
			\hline
			$e_{1}$ & 0 & $\dfrac{1}{5}$ & 1 & $\dfrac{1}{5}$ & $-\dfrac{1}{5}$ & 0 & 6 &  \\
			$x_{1}$ & 1 & $\dfrac{4}{5}$ & 0 & $-\dfrac{1}{5}$ & $\dfrac{1}{5}$ & 0 & 4 &  \\
			\hline
			$z$ & 0 & $\dfrac{7}{5}$ & 0 & $\dfrac{2}{5}$ & $-2-5 M$ & -1 & -8 &  \\
			\hline
		\end{tabular}
	\end{center}
	
	la variable artificiel $t_{1}$ est hors base avec un fort coefficient négatif et ne peut plus entrer dans la base. La phase I s'achève en éliminant la colonne de la variable $t_{1}$ du tableau. Nous retrouvons ainsi le tableau suivant avec lequel on va démarrer la phase II du simplexe comme on l'a déjà fait dans la section précédente.
	
	%Remarque 3.3 - 
	\begin{remarque}
		Il se peut que la phase I se termine avec une solution de base optimale comportant une ou plusieurs variables artificielles.
		
		\begin{itemize}
			\item Si l'une de ces variables artificielles de base prend une valeur négative, alors le problème initial est non réalisable.
			\item Si toutes les variables artificielles de base ont une valeur nulle, il faut les sortir de la base pour trouver le tableau initial du simplexe pour le problème $ ( P ) $. Pour cela, on doit remplacer chacune de ces variables artificielles par l'une des variables hors base, qui ne doit pas être une variable artificielle. Puisque la variable de sortie a une valeur nulle, alors la nouvelle variable de base aura aussi une valeur nulle.
		\end{itemize}
	\end{remarque}
	
	\section{Problèmes irréguliers}
	\subsection{Cas de solution non bornée}
	Graphiquement, ce problème est caractérisé par le fait qu'on peut déplacer la droite de la fonction objectif indéfiniment de manière à accroître la valeur, en gardant toujours une intersection non vide avec l'ensemble des solutions réalisables.
	
	Avec la méthode de simplexe, on reconnaît ce problème lorsque la variable entrante n'admet aucune limite sur sa valeur d'entrée, c'est à dire que tous les ratios $\dfrac{b_{i}}{a_{i j}}$ sont négatifs ou nuls.
	
	On considère le problème suivant :
	
	\[
	\left\{
	\begin{aligned}
	& \displaystyle\min \; -2x_1 - 3x_2 + x_3 \\
	& -x_1 - x_2 - x_3 \leq 3 \\
	& x_1 - x_2 + x_3 \leq 4 \\
	& -x_1 + x_2 + 2x_3 \leq 1 \\
	& x_1, x_2, x_3 \geq 0
	\end{aligned}
	\right.
	\overset{\text{rajout des variables d'écart}}{\Longrightarrow}
	\left\{
	\begin{aligned}
	& \displaystyle\min \; z \\
	& -x_1 - x_2 - x_3 + e_1 = 3 \\
	& x_1 - x_2 + x_3 + e_2 = 4 \\
	& -x_1 + x_2 + 2x_3 + e_3 = 1 \\
	& -2x_1 - 3x_2 + x_3 - z = 0 \\
	& x_1, x_2, x_3, e_1, e_2, e_3 \geq 0
	\end{aligned}
	\right.
	\]
	
	Le tableau initial s'écrit
	
	\begin{center}
		\begin{tabular}{|c|c>{\columncolor{gray!25}}cccccc|c|c|}
			\hline
			Variables de base & $x_{1}$ & \textcolor{blue}{$x_{2}$} & $x_{3}$ & $e_{1}$ & $e_{2}$ & $e_{3}$ & $-z$ & $b$ & $R$ \\
			\hline
			$e_{1}$ & -1 & -1 & -1 & 1 & 0 & 0 & 0 & 3 &  \\
			$e_{2}$ & 1 & -1 & 1 & 0 & 1 & 0 & 0 & 4 &  \\
			\rowcolor{gray!25}\textcolor{red}{$e_{3}$} & -1 & \tikz[baseline=(char.base)]{\node[draw=red,circle,very thick,inner sep=2pt] (char) {1};} & 2 & 0 & 0 & 1 & 0 & 1 &  \\
			\hline
			$-z$ & -2 & -3 & 1 & 0 & 0 & 3 & 1 & 0 &  \\
			\hline
		\end{tabular}
	\end{center}
	
	Variable sortante : $e_{3}$.\\
	Variable entrante : $x_{2}$.
	
	On pivote autour de de la troisième ligne et on trouve le tableau suivant
	
	\begin{center}
		\begin{tabular}{|c|>{\columncolor{gray!25}}ccccccc|c|c|}
			\toprule
			Variables de base & \textcolor{blue}{$x_{1}$} & $x_{2}$ & $x_{3}$ & $e_{1}$ & $e_{2}$ & $e_{3}$ & $-z$ & $b$ & $R$ \\
			\midrule
			$e_{1}$ & -2 & 0 & 1 & 1 & 0 & 0 & 0 & 4 &  \\
			$e_{2}$ & 0 & 0 & 3 & 0 & 1 & 0 & 0 & 5 &  \\
			$e_{3}$ & -1 & 1 & 2 & 0 & 0 & 1 & 0 & 1 &  \\
			\hline
			$-z$ & -5 & 0 & 7 & 0 & 0 & 0 & 1 & 3 &  \\
			\hline
		\end{tabular}
	\end{center}
	
	La variable entrante est $x_{1}$. Toutefois, le critère du quotient ne s'applique plus car toutes les entrées (lignes 1 à 3 ) de la colonne 1 sont négatives ou nulles.
	
	%Remarque 3.4 
	\begin{remarque}
		Selon cet exemple, on peut conclure que s'il n'est pas possible de calculer la ligne de pivot pour un choix de colonne de pivot $j$, le problème n'admet pas de solution optimale car il sera toujours non borné.
	\end{remarque}
	
	\section{Les problèmes impossibles}
	Graphiquement, on a caractérisé ces problèmes par un ensemble de solutions réalisables vide. Avec la méthode de simplexe, on reconnaît que le problème est impossible si une ou plusieurs variables artificielles sont présentes dans la base dans le tableau de simplexe optimal, ce qui signifie que la solution donnée par ce tableau n'est pas réellement réalisable.
	
	Vérifions à l'aide de la méthode de simplexe, que le problème suivant est réellement impossible:
	
	$$
	\left\{\begin{array}{c}
	\displaystyle\max \quad 4 x_{1}+3 x_{2} \\
	x_{1}+x_{2} \leq 2 \\
	3 x_{1}+x_{2} \geq 10 \\
	x_{1}, x_{2} \geq 0
	\end{array}\right.
	$$
	
	En introduisant les variables d'écarts et les variables artificielles le programme s'écrit :
	
	$$
	\left\{\begin{array}{c}
	\displaystyle\max \quad 4 x_{1}+3 x_{2}-M t \\
	x_{1}+x_{2}+e_{1}=2 \\
	3 x_{1}+x_{2}-e_{2}+t=10 \\
	x_{1}, x_{2}, e_{1}, e_{2}, t \geq 0
	\end{array}\right.
	$$
	
	La résolution par le simplexe donne les tableaux suivant
	
	\begin{center}
		\begin{tabular}{|c|cccccc|c|c|}
			\hline
			Variables de base & $x_{1}$ & $x_{2}$ & $e_{1}$ & $e_{2}$ & $t_{1}$ & $z$ & $b$ & $R$ \\
			\hline
			$e_{1}$ & 1 & 1 & 1 & 0 & 0 & 0 & 2 &  \\
			$t_{1}$ & 3 & 1 & 0 & -1 & 1 & 0 & 10 &  \\
			\hline
			$z$ & 4 & 3 & 0 & 0 & $-M$ & -1 & 0 &  \\
			\hline
		\end{tabular}
	\end{center}
	
	Ainsi,
	
	\begin{center}
		\begin{tabular}{|c|cccccc|c|c|}
			\hline
			Variables de base & $x_{1}$ & $x_{2}$ & $e_{1}$ & $e_{2}$ & $t_{1}$ & $z$ & $b$ & $R$ \\
			\hline
			$e_{1}$ & 1 & 1 & 1 & 0 & 0 & 0 & 2 &  \\
			$t_{1}$ & 3 & 1 & 0 & -1 & 1 & 0 & 10 &  \\
			\hline
			$z$ & $4+3 M$ & $3+M$ & 0 & $-M$ & 0 & -1 & $10 M$ &  \\
			\hline
		\end{tabular}
	\end{center}
	
	D'où,
	
	\begin{center}
		\begin{tabular}{|c|cccccc|c|c|}
			\hline
			Variables de base & $x_{1}$ & $x_{2}$ & $e_{1}$ & $e_{2}$ & $t_{1}$ & $z$ & $b$ & $R$ \\
			\hline
			$x_{1}$ & 1 & 1 & 1 & 0 & 0 & 0 & 2 &  \\
			$t_{1}$ & 0 & -2 & -3 & -1 & 1 & 0 & 10 &  \\
			\hline
			$z$ & 0 & $-1-2 M$ & $-4-3 M$ & 0 & $-M$ & -1 & $4 M-8$ &  \\
			\hline
		\end{tabular}
	\end{center}
	
	Le tableau de simplexe ci-dessus est optimal avec une variable artificielle dans la base.
	
	\section{Les problèmes à solutions multiples}
	Nous avons vu dans l'algorithme du simplexe que lorsque les coûts marginaux de toutes les variables hors base sont $\leq 0$, alors la solution actuelle est optimale. Mais s'ils sont tous non nuls $ (<0) $, alors la solution optimale est unique. Si par contre, l'un de ces coûts des variables hors base est nul, alors on peut avoir deux ou plusieurs solutions de base optimales. Pour illustrer, reprenons l'exercice de la section résolution graphique qui admet une infinité de solutions optimales.
	
	$$
	\left\{\begin{aligned}
	\displaystyle\min -50 x_{1}-30 x_{2} & \\
	\text { s.t } 10 x_{1}+6 x_{2} & \leq 45 \\
	4 x_{1}+6 x_{2} & \leq 36 \\
	2 x_{1}+6 x_{2} & \leq 27 \\
	x & \geq 0
	\end{aligned}\right.
	$$
	
	\begin{center}
		\begin{tabular}{|c|cccccc|c|c|}
			\hline
			Variables de base & $x_{1}$ & $x_{2}$ & $e_{1}$ & $e_{2}$ & $e_{3}$ & $-z$ & $b$ & $R$ \\
			\hline
			$e_{1}$ & 10 & 6 & 1 & 0 & 0 & 0 & 45 &  \\
			$e_{2}$ & 4 & 6 & 0 & 1 & 0 & 0 & 36 &  \\
			$e_{3}$ & 2 & 6 & 0 & 0 & 1 & 0 & 27 &  \\
			\hline
			$-z$ & -50 & -30 & 0 & 0 & 0 & 1 & 0 &  \\
			\hline
		\end{tabular}
	\end{center}
	
	La variable $x_{1}$ entre dans la base pour remplacer la variable $e_{1}$
	
	\begin{center}
		\begin{tabular}{|c|cccccc|c|c|}
			\hline
			Variables de base & $x_{1}$ & $x_{2}$ & $e_{1}$ & $e_{2}$ & $e_{3}$ & $-z$ & $b$ & $R$ \\
			\hline
			$x_{1}$ & 1 & $\dfrac{3}{5}$ & $\dfrac{1}{10}$ & 0 & 0 & 0 & $\dfrac{9}{2}$ &  \\
			$e_{2}$ & 0 & $\dfrac{18}{5}$ & $-\dfrac{2}{5}$ & 1 & 0 & 0 & 18 &  \\
			$e_{3}$ & 0 & $\dfrac{24}{5}$ & $-\dfrac{1}{5}$ & 0 & 1 & 0 & 18 &  \\
			\hline
			$-z$ & 0 & 0 & 5 & 0 & 0 & 1 & 225 &  \\
			\hline
		\end{tabular}
	\end{center}
	
	Nous obtenons une solution de base optimale dont les variables de base sont $x_{1}, e_{2}$ et $e_{3}$. Mais le coût marginal de la variable hors base $x_{2}$ est nul. Faisons entrer cette variable dans la base et appliquons la règle du plus petit rapport pour choisir la variable de sortie qui est dans ce cas $e_{3}$. Nous obtenons le tableau suivant
	
	\begin{center}
		\begin{tabular}{|c|cccccc|c|c|}
			\hline
			Variables de base & $x_{1}$ & $x_{2}$ & $e_{1}$ & $e_{2}$ & $e_{3}$ & $-z$ & $b$ & $R$ \\
			\hline
			$x_{1}$ & 1 & 0 & $\dfrac{1}{8}$ & 0 & $-\dfrac{1}{8}$ & 0 & $\dfrac{9}{4}$ &  \\
			$e_{2}$ & 0 & 0 & $-\dfrac{1}{4}$ & 1 & $-\dfrac{3}{4}$ & 0 & $\dfrac{9}{2}$ &  \\
			$x_{2}$ & 0 & 1 & $-\dfrac{1}{24}$ & 0 & $\dfrac{5}{24}$ & 0 & $\dfrac{15}{4}$ &  \\
			\hline
			$-z$ & 0 & 0 & 5 & 0 & 0 & 1 & 225 &  \\
			\hline
		\end{tabular}
	\end{center}
	
	Nous obtenons donc une autre solution de base optimale dont les variables de base sont $x_{1}, e_{2}$ et $x_{2}$.
	
	\subsection{Les problèmes à solution dégénérée}
	En appliquons la règle du plus petit rapport pour désigner la variable de sortie, il est possible que le plus petit rapport soit atteint pour deux contraintes. En effectuant le changement de base, nous obtenons alors une solution de base avec l'une des composante nulle. Nous disons qu'on a une dégénérescence. Ceci peut entraîner des difficultés d'exécution et d'efficacité de l'algorithme du simplexe. En effet, nous pouvons faire plusieurs changement de base sans améliorer la valeur de la fonction objectif. Nous pouvons même retrouver une base déjà rencontrer à une itération précédente. On dit alors qu'on a un phénomène de cyclage. Afin d'illustrer ce type de dégénérescence. considérons l'exemple suivant :
	
	$$
	\left\{\begin{array}{rl}
	\displaystyle\min & -4 x_{1}-3 x_{3} \\
	s . t & x_{1}-x_{2} \\
	& \leq 2 \\
	& 2 x_{1}+x_{3} \\
	& \leq 4 \\
	x_{1}+x_{2}+x_{3} & \leq 3 \\
	& x_{1}, x_{2}, x_{3}
	\end{array} \geq 0\right.
	$$
	
	Introduisons les variables d'écart et écrivons le premier tableau du simplexe :
	
	\begin{center}
		\begin{tabular}{|c|ccccccc|c|c|}
			\hline
			Variables de base & $x_{1}$ & $x_{2}$ & $x_{3}$ & $e_{1}$ & $e_{2}$ & $e_{3}$ & $-z$ & $b$ & $R$ \\
			\hline
			$e_{1}$ & 1 & -1 & 0 & 1 & 0 & 0 & 0 & 2 &  \\
			$e_{2}$ & 2 & 0 & 1 & 0 & 1 & 0 & 0 & 4 &  \\
			$e_{3}$ & 1 & 1 & 1 & 0 & 0 & 1 & 0 & 3 &  \\
			\hline
			$-z$ & -4 & 0 & -3 & 0 & 0 & 0 & 1 & 0 &  \\
			\hline
		\end{tabular}
	\end{center}
	
	La variable $x_{1}$ entre dans la base, mais en appliquons la règle du plus petit rapport, nous constatons que le minimum des rapports est atteint pour les deux premières contraintes, alors que seulement une des deux doit quitter la base. Choisissons arbitrairement $e_{1}$ comme variable sortante. nous obtenons alors le tableau suivant :
	
	\begin{center}
		\begin{tabular}{|c|ccccccc|c|c|}
			\hline
			Variables de base & $x_{1}$ & $x_{2}$ & $x_{3}$ & $e_{1}$ & $e_{2}$ & $e_{3}$ & $-z$ & $b$ & $R$ \\
			\hline
			$x_{1}$ & 1 & -1 & 0 & 1 & 0 & 0 & 0 & 2 &  \\
			$e_{2}$ & 0 & 2 & 1 & -2 & 1 & 0 & 0 & 0 &  \\
			$e_{3}$ & 0 & 2 & 1 & -1 & 0 & 1 & 0 & 1 &  \\
			\hline
			$-z$ & 0 & -4 & -3 & 4 & 0 & 0 & 1 & 8 &  \\
			\hline
		\end{tabular}
	\end{center}
	
	cette solution de base n'est pas optimale puisque les coûts marginaux des variables hors base ne sont pas tous positifs. Cette solution de base est dite dégénérée puisque l'une des variables de base est nulle ( $e_{2}=0$ ). La variable $x_{2}$ entre dans la base. En appliquons la règle du plus petit rapport, nous constatons que la variable $e_{2}$ quitte la base et nous obtenons le tableau suivant
	
	\begin{center}
		\begin{tabular}{|c|ccccccc|c|c|}
			\hline
			Variables de base & $x_{1}$ & $x_{2}$ & $x_{3}$ & $e_{1}$ & $e_{2}$ & $e_{3}$ & $-z$ & $b$ & $R$ \\
			\hline
			$x_{1}$ & 1 & 0 & $\dfrac{1}{2}$ & 0 & $\dfrac{1}{2}$ & 0 & 0 & 2 &  \\
			$x_{2}$ & 0 & 1 & $\dfrac{1}{2}$ & -1 & $\dfrac{1}{2}$ & 0 & 0 & 0 &  \\
			$e_{3}$ & 0 & 0 & 0 & 1 & -1 & 1 & 0 & 1 &  \\
			\hline
			$-z$ & 0 & 0 & -1 & 0 & 2 & 0 & 1 & 8 &  \\
			\hline
		\end{tabular}
	\end{center}
	
	Nous constatons que la nouvelle solution de base est non optimale et qu'elle est aussi dégénérée $\left(x_{2}=0\right)$. Ce changement de base n'a pas modifié la valeur de la fonction objectif $(z=-8)$. A l'itération suivante de l'algorithme du simplexe, cette variable $x_{2}$ quitte la base et sera remplacée par la variable $x_{3}$
	
	\begin{center}
		\begin{tabular}{|c|c|c|c|c|c|c|c|c|c|}
			\hline
			Variables de base & $x_{1}$ & $x_{2}$ & $x_{3}$ & $e_{1}$ & $e_{2}$ & $e_{3}$ & $-z$ & $b$ & $R$ \\
			\hline
			$x_{1}$ & 1 & -1 & 0 & 1 & 0 & 0 & 0 & 2 &  \\
			$x_{3}$ & 0 & 2 & 1 & -2 & 1 & 0 & 0 & 0 &  \\
			$e_{3}$ & 0 & 0 & 0 & 1 & -1 & 1 & 0 & 1 &  \\
			\hline
			$-z$ & 0 & 2 & 0 & -2 & 3 & 0 & 1 & 8 &  \\
			\hline
		\end{tabular}
	\end{center}
	
	Encore une fois, ce changement de base n'a pas modifié la valeur de la fonction objectif. En plus la nouvelle solution de base est aussi non optimale. Nous pouvons donc se poser la question de savoir si nous allons continuer ainsi sans pouvoir trouver la solution de base optimale. Heureusement ce phénomène de dégénérescence s'arrête à l'itération suivante
	
	\begin{center}
		\begin{tabular}{|c|c|c|c|c|c|c|c|c|c|}
			\hline
			Variables de base & $x_{1}$ & $x_{2}$ & $x_{3}$ & $e_{1}$ & $e_{2}$ & $e_{3}$ & $-z$ & $b$ & $R$ \\
			\hline
			$x_{1}$ & 1 & -1 & 0 & 0 & 1 & -1 & 0 & 1 &  \\
			$x_{3}$ & 0 & 2 & 1 & 0 & -1 & 2 & 0 & 2 &  \\
			$e_{1}$ & 0 & 0 & 0 & 1 & -1 & 1 & 0 & 1 &  \\
			$-z$ & 0 & 2 & 0 & 0 & 1 & 2 & 1 & 10 &  \\
			\hline
		\end{tabular}
	\end{center}
	
	%\textbf{Remarque 3.5} 
	\begin{remarque}
		Dans le cas où on a plusieurs variables candidates d'entrer en base et plusieurs variables candidates de sortir de la base, la règle du plus petit indice, découverte par R. Bland, peut nous éviter le cyclage infini. Cette règle consiste à choisir systématiquement, comme variable entrante, parmi les multiples candidats, la variable $x_{e}$ qui correspond au plus petit indice $e$ et comme variable sortante, la variable $x_{s}$ qui correspond au plus petit indice $s$.
	\end{remarque}	
	
	
	
	
	
	\chapter{Dualité}
	
	\section{Motivation et définitions}
	
	Revenons encore une fois au problème de production de l'introduction.
	
	\begin{center}
		\begin{tabular}{|c|c|c|c|}
			\hline
			& $P_{1}$ & $P_{2}$ & disponibilité \\
			\hline
			équipement & 3 & 9 & 81 \\
			main d'œuvre & 4 & 5 & 55 \\
			matière première & 2 & 1 & 20 \\
			bénéfice & 6 & 4 &  \\
			\hline
		\end{tabular}
	\end{center}
	
	Supposons à présent qu'un acheteur se présente pour acheter toutes les ressources $b_{i}$, $1 \leq i \leq m$, de l'entreprise. Il propose à l'entreprise un prix unitaire $y_{i}$, $1 \leq i \leq m$, pour chacune des ressources. L'entreprise acceptera de lui vendre toutes ses ressources uniquement si elle obtient pour chaque produit $P_{j}$ un prix de vente au moins égal au profit $c_{j}$ qu'elle ferait en vendant ses produits. On doit donc avoir
	
	\[
	\begin{aligned}
	& 3 y_{1}+4 y_{2}+2 y_{3} \geq c_{1}=6 \\
	& 9 y_{1}+5 y_{2}+1 y_{3} \geq c_{2}=4 \\
	& y_{1}, y_{2}, y_{3} \geq 0
	\end{aligned}
	\]
	
	De son côté, l'acheteur cherche à minimiser le prix total d'achat. On se demande alors quels sont les prix unitaires $y_{i}$, $1 \leq i \leq m$, que l'acheteur doit proposer à l'entreprise en question pour qu'elle accepte de vendre toutes ses ressources~? Le problème peut donc s'écrire
	
	\[
	\left\{
	\begin{array}{ll}
	\displaystyle\min_{y} & g(y)=81 y_{1}+55 y_{2}+20 y_{3} \\
	\text{s.t.} & 3 y_{1}+4 y_{2}+2 y_{3} \geq c_{1}=6 \\
	& 9 y_{1}+5 y_{2}+1 y_{3} \geq c_{2}=4 \\
	& y_{1}, y_{2}, y_{3} \geq 0
	\end{array}
	\right.
	\]
	
	\begin{definition}
		Au programme linéaire primal
		
		\[
		(PL):\left\{
		\begin{array}{ll}
		\displaystyle\max_{x \in \mathbb{R}^{n}} & f(x)=c^{t} x \\
		\text{s.t.} & A x \leq b \\
		& x \geq 0
		\end{array}
		\right.
		\]
		
		on associe le programme linéaire dual
		
		\[
		(PLD):\left\{
		\begin{array}{ll}
		\displaystyle\min_{y \in \mathbb{R}^{m}} & g(y)=b^{t} y \\
		\text{s.t.} & A^{t} y \geq c \\
		& y \geq 0
		\end{array}
		\right.
		\]
	\end{definition}
	
	\begin{exemple}
		Illustration graphique du problème primal et dual~: On considère le problème suivant~:
		
		\[
		\left\{
		\begin{array}{ll}
		\displaystyle\min_{x \in \mathbb{R}^{2}} & f(x)=-x_{1}-x_{2} \\
		\text{s.t.} & -x_{1}-3 x_{2} \geq -3 \\
		& -2 x_{1}-x_{2} \geq -2 \\
		& x_{1}, x_{2} \geq 0
		\end{array}
		\right.
		\]
		
		qui admet la solution $x^{*}=\left(\dfrac{3}{5}, \dfrac{4}{5}\right)$ et $f\left(x^{*}\right)=-\dfrac{7}{5}$.
		
		\begin{figure}[H]
			\centering
			\begin{tikzpicture}[scale=2]
			% Zone de travail
			\clip (-0.5, -0.5) rectangle (4, 3);
			
			% Région réalisable (zone cyan)
			%\fill[cyan, opacity=0.8] (0, 0) -- (0.6, 0.8) -- (0, 1) -- (0,0) -- cycle;
			\fill[purple!30, opacity=0.9] (0, 1) -- (0.6, 0.8) -- (1, 0) -- (1.2, 0) -- (0, 0) -- (0, 1.2) -- cycle;
			% Dessiner les axes
			\draw[->] (-0.3, 0) -- (3.5, 0) node[right] {$x_1$};
			\draw[->] (0, -0.3) -- (0, 2.5) node[above] {$x_2$};
			
			% Contrainte 1: x₁ + 3x₂ = 3
			\draw[thick, gray] (0, 1) -- (3, 0);
			
			% Contrainte 2: 2x₁ + x₂ = 2
			\draw[thick, gray] (0, 2) -- (1, 0);
			
			% Frontière de la région réalisable
			%\draw[thick] (0, 0) -- (0.6, 0.8) -- (0, 1);
			
			% Graduations
			\draw (1, 0) -- (1, -0.05) node[below] {$1$};
			\draw (3, 0) -- (3, -0.05) node[below] {$3$};
			\draw (0, 1) -- (-0.05, 1) node[left] {$1$};
			\draw (0, 2) -- (-0.05, 2) node[left] {$2$};
			
			\end{tikzpicture}
			\caption{domaine primal}
		\end{figure} 
		
		
		
		
		Le problème dual s'écrit sous la forme~:
		
		\[
		\left\{
		\begin{array}{ll}
		\displaystyle\max_{y \in \mathbb{R}^{2}} & g(y)=-3 y_{1}-2 y_{2} \\
		\text{s.t.} & -y_{1}-2 y_{2} \leq -1 \\
		& -3 y_{1}-y_{2} \leq -1 \\
		& y_{1}, y_{2} \geq 0
		\end{array}
		\right.
		\]
		
		qui admet la solution $y^{*}=\left(\dfrac{1}{5}, \dfrac{2}{5}\right)$ et $g\left(y^{*}\right)=-\dfrac{7}{5}$. Dans cet exemple, on observe que la valeur minimale du primal est égale à la valeur maximale du dual.
		
		\begin{figure}[H]
			\centering
			\begin{tikzpicture}[scale=4] % Ajustez le scale si nécessaire pour la taille
			% Définir une zone de travail légèrement plus grande que la région d'intérêt
			\clip (-0.5, -0.5) rectangle (1.25, 1.25);
			
			% Remplir la région réalisable (zone bleue)
			% Les sommets de la région colorée sont (0,1), (1/3, 1/3), (1,0)
			% Et on les étend vers les bords du graphique pour simuler l'infini dans le quadrant positif
			\fill[purple!30, opacity=0.9] (0, 1) -- (1/5, 4/10) -- (1, 0) -- (1.2, 0) -- (1.2, 1.2) -- (0, 1.2) -- cycle;
			
			% Dessiner les axes avec flèches uniquement à l'extrémité positive
			\draw[->] (-0.2, 0) -- (1.2, 0) node[below] {$y_1$};
			\draw[->] (0, -0.2) -- (0, 1.2) node[left] {$y_2$};
			
			% Dessiner la frontière de la région en gras
			\draw[thick] (0, 1) -- (1/5, 4/10) -- (1, 0);
			
			% Dessiner la ligne traversante, qui semble aller de (0, 1/2) à (1/3, 0)
			\draw (0, 0.5) -- (1, 0);
			
			\draw (0, 1) -- (1/3, 0);
			
			% Ajouter les graduations et les étiquettes sur les axes
			% Axe Y (y₂)
			\draw (0, 1) -- (-2pt, 1) node[left] {$1$};
			\draw (0, 0.5) -- (-2pt, 0.5) node[left] {$1/2$}; % Utilisation de 1/2 au lieu de \dfrac{1}{2} pour l'alignement
			
			% Axe X (y₁)
			\draw (1, 0) -- (1, -2pt) node[below] {$1$};
			\draw (1/3, 0) -- (1/3, -2pt) node[below] {$1/3$}; % Utilisation de 1/3 au lieu de \dfrac{1}{3} pour l'alignement
			
			\end{tikzpicture}
			\caption{domaine dual}
		\end{figure}
	\end{exemple}
	
	\subsection{Comparaison primal/dual}
	
	\begin{center}
		\begin{tabular}{|c|c|}
			\hline
			Primal & Dual \\
			\hline
			$\displaystyle\max (f)$ & $\displaystyle\min (g)$ \\
			coefficient $c$ de $f$ & second membre $c$ \\
			second membre $b$ & coefficient $b$ de $g$ \\
			$m$ contraintes inégalités $(\leq)$ & $m$ contraintes de positivité \\
			$n$ contraintes de positivité & $n$ contraintes inégalités $(\geq)$ \\
			\hline
		\end{tabular}
	\end{center}
	
	Si le problème primal est sous forme standard avec les contraintes $A x=b$ alors on passe à la forme canonique pure en écrivant les contraintes $A x=b \quad \Longleftrightarrow \quad \binom{A}{-A} x \leq \binom{b}{-b}$. De façon générale, on a la définition suivante lorsque le problème primal est sous forme canonique mixte~:
	
	\begin{center}
		\begin{tabular}{|c|c|}
			\hline
			Primal & Dual \\
			\hline
			$\displaystyle\max_{x \in \mathbb{R}^{n}} f(x)=c^{t} x$ & $\displaystyle\displaystyle\min_{y \in \mathbb{R}^{m}} g(y)=b^{t} y$ \\
			$\forall i \in I_{1}, \displaystyle\sum_{j=1}^{n} a_{ij} x_{j} \leq b_{i}$ & $\forall i \in I_{1}, y_{i} \geq 0$ \\
			$\forall i \in I_{2}, \displaystyle\sum_{j=1}^{n} a_{ij} x_{j}=b_{i}$ & $\forall i \in I_{2}, y_{i}$ de signe quelconque \\
			$\forall j \in J_{1}, x_{j} \geq 0$ & $\forall j \in J_{1}, \displaystyle\sum_{i=1}^{m} a_{ij} y_{i} \geq c_{j}$ \\
			$\forall j \in J_{2}, x_{j}$ de signe quelconque & $\forall j \in J_{2}, \displaystyle\sum_{i=1}^{m} a_{ij} y_{i}=c_{j}$ \\
			\hline
		\end{tabular}
	\end{center}
	
	\section{Propriétés -- Théorèmes de dualité}
	
	\begin{proposition}
		Le dual du dual est le primal.
	\end{proposition}
	
	Démonstration~: A partir de la définition du dual d'un PL sous forme canonique pure~:
	
	\[
	\left\{
	\begin{array}{l}
	\displaystyle\min_{y \in \mathbb{R}^{m}} g(y)=b^{t} y \\
	\text{s.t.}\\
	A^{t} y \geq c \\
	y \geq 0
	\end{array}
	\quad
	\right.
	\Longleftrightarrow
	\left\{
	\begin{array}{l}
	-\displaystyle\max_{y \in \mathbb{R}^{m}}-g(y)=(-b)^{t} y \\
	\text{s.t.}\\
	A^{t} y \leq -c \\
	y \geq 0
	\end{array}
	\right.
	\]
	Si on prend le dual du dual, on a donc
	
	\[
	\left\{
	\begin{array}{l}
	\displaystyle\min_{x \in \mathbb{R}^{n}} (-c)^{t} x \\
	(-A^{t})^{t} x \geq (-b) \\
	x \geq 0
	\end{array}
	\right.
	\Longleftrightarrow
	\left\{
	\begin{array}{l}
	\displaystyle\max_{x \in \mathbb{R}^{n}} c^{t} x \\
	A x \leq b \\
	x \geq 0
	\end{array}
	\right.
	\]
	
	et on reconnaît le problème primal du départ.
	
	\begin{theorem}[Théorème faible de dualité]
		Soit $x$ une solution réalisable d'un (PL) sous forme canonique mixte et $y$ une solution réalisable du problème dual (PLD). Alors~:
		
		\begin{enumerate}
			\item $f(x) \leq g(y)$.
			\item Si $f(x)=g(y)$ alors $x$ et $y$ sont des solutions optimales de $(PL)$ et $(PLD)$ respectivement.
		\end{enumerate}
	\end{theorem}
	
%	Démonstration~: 
	\begin{proof}
		Dans le cas où (PL) est sous forme canonique pure~:
		
		\begin{enumerate}
			\item On a d'une part $A x \leq b, x \geq 0$ et d'autre part $A^{t} y \geq c, y \geq 0$. Par conséquent,
			\[
			f(x)=c^{t} x \leq (A^{t} y)^{t} x = y^{t} \underbrace{A x}_{\leq b} \leq y^{t} b = g(y) \quad \text{car } y \geq 0
			\]
			\item Soient $x^{*}$ et $y^{*}$ des solutions réalisables de (PL) et (PLD) respectivement telles que $f(x^{*})=g(y^{*})$. D'après 1., pour toute solution réalisable $x$ de (PL), on a $f(x) \leq g(y^{*})=f(x^{*})$ donc $x^{*}$ est une solution réalisable optimale. Idem pour $y^{*}$.
		\end{enumerate}
	\end{proof}
	
	\begin{theorem}[Théorème fort de dualité]
		Si le problème primal (PL) admet une solution réalisable optimale $x^{*}$ alors le problème dual (PLD) admet lui aussi une solution réalisable optimale $y^{*}$ et on a
		\[
		f(x^{*}) = g(y^{*})
		\]
	\end{theorem}
	
	Démonstration~: On suppose que (PL) est sous forme standard. Si (PL) admet une solution réalisable optimale, alors il admet une solution de base réalisable optimale. On note $B^{*}$ la base optimale et on désigne par $x^{*}$ la solution de base réalisable optimale~: $x^{*}=A_{B^{*}}^{-1} b$. On choisit alors
	\[
	y^{*} = \left(A_{B^{*}}^{-1}\right)^{t} c_{B^{*}}
	\]
	Montrons que $y^{*}$ est une solution réalisable optimale pour le dual (PLD). On a
	\[
	\begin{aligned}
	A_{H^{*}}^{t} y^{*} & = A_{H^{*}}^{t} \left(A_{B^{*}}^{-1}\right)^{t} c_{B^{*}} \\
	& = \left(A_{B^{*}}^{-1} A_{H^{*}}\right)^{t} c_{B^{*}} \\
	& = c_{H^{*}} - L_{H^{*}}
	\end{aligned}
	\]
	Or, à l'optimum tous les coûts réduits sont négatifs ou nuls ($L_{H^{*}} \leq 0$) donc $A_{H^{*}}^{t} y^{*} \geq c_{H^{*}}$. Par définition, on a $A_{B^{*}}^{t} y^{*} = c_{B^{*}}$ et donc on obtient
	\[
	A^{t} y^{*} \geq c
	\]
	$y^{*}$ est de signe quelconque. Par conséquent, $y^{*}$ est une solution réalisable du dual (PLD). On remarquera que le problème primal (PL) étant mis sous forme standard, il n'y a pas de contrainte de positivité sur les variables $y$ du dual. Par ailleurs
	\[
	\begin{aligned}
	f(x^{*}) & = c^{t} x^{*} = c_{B^{*}}^{t} A_{B^{*}}^{-1} b \\
	& = \underbrace{\left(\left(A_{B^{*}}^{-1}\right)^{t} c_{B^{*}}\right)^{t} b}_{y^{*}} \\
	& = g(y^{*})
	\end{aligned}
	\]
	et donc par le Théorème faible de dualité, $y^{*}$ est optimale pour (PLD).
	
	On a vu qu'il y avait 3 cas possibles (et seulement 3) pour le problème primal (PL)~:
	
	\begin{enumerate}
		\item il existe (au moins) une solution optimale.
		\item l'ensemble $\mathbf{D}_{R}$ des solutions réalisables n'est pas borné et l'optimum est infini.
		\item pas de solution réalisable $\left(\mathbf{D}_{R}=\emptyset\right)$
	\end{enumerate}
	
	Les mêmes situations se retrouvent pour le problème dual. Plus précisément, le lien entre les deux problèmes en dualité est donné par le résultat suivant.
	
	\begin{theorem}
		Étant donnés un problème primal (PL) et son dual (PLD), une et une seule des trois situations suivantes peut se produire.
		\begin{enumerate}
			\item[(a)] les deux problèmes possèdent chacun des solutions optimales (à l'optimum, les coûts sont égaux).
			\item[(b)] un des problèmes possède une solution réalisable avec un optimum infini, l'autre n'a pas de solution.
			\item[(c)] aucun des deux problèmes ne possède de solution réalisable.
		\end{enumerate}
	\end{theorem}
	
	Il y a donc 3 situations (au lieu de 9) qui peuvent se résumer dans le tableau suivant~:
	
	\begin{center}
		\begin{tabular}{cc||c|c|c|}
			\multicolumn{2}{c||}{} & \multicolumn{3}{c|}{ \textbf{Dual} } \\
			\cline{3-5}
			\multicolumn{2}{c||}{} & \textit{(1) Solution optimale} & \textit{(2) Optimum infini} & \textit{(3) pas de solution} \\
			\hline \hline
			\multicolumn{1}{|c|}{\multirow{3}{*}{\rotatebox{90}{\textbf{Primal}}}} &
			\multicolumn{1}{|l||}{\textit{(1) Solution optimale}} & (a) & impossible & impossible \\
			\cline{2-5}
			\multicolumn{1}{|c|}{} &
			\multicolumn{1}{|l||}{\textit{(2) Optimum infini}} & impossible & impossible & (b) \\
			\cline{2-5}
			\multicolumn{1}{|c|}{} &
			\multicolumn{1}{|l||}{\textit{(3) pas de solution}} & impossible & (c) & impossible \\
			\hline
		\end{tabular}
	\end{center}
	
	
	
	\section{Conditions d'optimalité primal-dual (COPD)}
	
	On s'intéresse au cas (a) où les problèmes primal et dual possèdent chacun des solutions optimales (optimum fini). On peut alors calculer l'une à partir de l'autre.
	
	\begin{theorem}\label{theo4.4}
		Soient $x$ et $y$ des solutions réalisables respectivement du problème primal (PL) et du problème dual (PLD). Alors $x$ et $y$ sont des solutions réalisables optimales si et seulement si les conditions d'optimalité primal-dual (COPD) suivantes sont vérifiées~:
		\begin{itemize}
			\item Si une contrainte est satisfaite en tant qu'inégalité stricte dans (PL) (resp. dans (PLD)) alors la variable correspondante de (PLD) (resp. de (PL)) est nulle.
			\item Si la valeur d'une variable dans (PL) ou (PLD) est strictement positive alors la contrainte correspondante de l'autre programme est une égalité.
		\end{itemize}
	\end{theorem}
	
	Supposons que le problème primal s'écrive sous forme canonique mixte. Alors $x$ et $y$ sont optimales respectivement pour le problème primal et le problème dual si et seulement si on a les COPD~:
	\[
	\left\{
	\begin{array}{l}
	\forall i \in I_{1}, \displaystyle\sum_{j=1}^{n} a_{ij} x_{j} = b_{i} \ \text{ou} \ y_{i} = 0 \\
	\forall j \in J_{1}, \displaystyle\sum_{i=1}^{m} a_{ij} y_{i} = c_{j} \ \text{ou} \ x_{j} = 0
	\end{array}
	\right.
	\]
	
	%\textbf{Démonstration} ()~:  
	\begin{proof}[condition nécessaire du Théorème \ref{theo4.4}]
		Supposons pour simplifier que le problème primal (PL) soit mis sous forme canonique pure. Soient $x$ et $y$ des solutions réalisables optimales de (PL) et (PLD) respectivement. On a donc $A x \leq b, x \geq 0$ et $A^{t} y \geq c, y \geq 0$. En introduisant les variables d'écart $e$ et $\epsilon$ respectivement pour (PL) et (PLD), on a
		\[
		\begin{array}{ccc}
		A x + e = b & \\
		x, e \geq 0 & \text{et} & A^{t} y - \epsilon = c \\
		& & y, \epsilon \geq 0
		\end{array}
		\]
		Dans ces conditions~:
		\[
		\begin{aligned}
		& f(x) = c^{t} x = \left(A^{t} y - \epsilon\right)^{t} x = y^{t} A x - \epsilon^{t} x \\
		& g(y) = b^{t} y = (A x + e)^{t} y = (A x)^{t} y + e^{t} y = y^{t} A x + e^{t} y
		\end{aligned}
		\]
		Or d'après le Théorème de la dualité forte, $f(x) = g(y)$ donc on obtient~:
		\[
		\epsilon^{t} x + e^{t} y = 0
		\]
		Puisque $x \geq 0$ et $y \geq 0$, on a nécessairement~:
		\[
		\left\{
		\begin{array}{l}
		\epsilon_{i} x_{i} = 0, \ \forall i \\
		e_{j} y_{j} = 0, \ \forall j
		\end{array}
		\right.
		\]
		ce qui donne les COPD (relations d'exclusion)~:
		\[
		\left\{
		\begin{array}{l}
		\text{Si } \epsilon_{i} \neq 0 \ \text{alors} \ y_{i} = 0 \\
		\text{Si } x_{i} \neq 0 \ \text{alors} \ \epsilon_{i} = 0
		\end{array}
		\right.
		\quad
		\left\{
		\begin{array}{l}
		\text{Si } e_{j} \neq 0 \ \text{alors} \ y_{j} = 0 \\
		\text{Si } y_{j} \neq 0 \ \text{alors} \ e_{j} = 0
		\end{array}
		\right.
		\]
		La réciproque (condition suffisante) se démontre à partir du Théorème faible de dualité.
	\end{proof}
	
	\subsection{Utilisation pratique des COPD}
	
	La dualité et les COPD permettent souvent de vérifier si une solution réalisable d'un (PL) est optimale ou non, à partir de la connaissance d'une solution optimale du problème dual. Lorsque (PL) et (PLD) ont des solutions réalisables optimales $x^{*}$ et $y^{*}$ respectivement, on a~:
	\[
	\begin{aligned}
	& \displaystyle\sum_{j=1}^{n} a_{ij} x_{j}^{*} < b_{i} \Longrightarrow y_{i}^{*} = 0 \\
	& \displaystyle\sum_{i=1}^{m} a_{ij} y_{i}^{*} > c_{j} \Longrightarrow x_{j}^{*} = 0
	\end{aligned}
	\]
	et
	\[
	\begin{aligned}
	& y_{i}^{*} > 0 \Longrightarrow \displaystyle\sum_{j=1}^{n} a_{ij} x_{j}^{*} = b_{i} \\
	& x_{j}^{*} > 0 \Longrightarrow \displaystyle\sum_{i=1}^{m} a_{ij} y_{i}^{*} = c_{j}
	\end{aligned}
	\]
	
	\begin{exemple}
		Le problème dual du problème de production s'écrit~:
		\[
		\left\{
		\begin{array}{ll}
		\displaystyle\min_{y} & g(y) = 81 y_{1} + 55 y_{2} + 20 y_{3} \\
		\text{s.t.} & 3 y_{1} + 4 y_{2} + 2 y_{3} \geq c_{1} = 6 \\
		& 9 y_{1} + 5 y_{2} + 1 y_{3} \geq c_{2} = 4 \\
		& y_{1}, y_{2}, y_{3} \geq 0
		\end{array}
		\right.
		\]
		Le problème primal (PL) admet la solution optimale~:
		\[
		\begin{aligned}
		& e_{1}^{*} = \dfrac{27}{2} > 0 \stackrel{\text{COPD}}{\Longrightarrow} y_{1}^{*} = 0 \\
		& x_{1}^{*} = \dfrac{15}{2} > 0 \stackrel{\text{COPD}}{\Longrightarrow} 3 y_{1}^{*} + 4 y_{2}^{*} + 2 y_{3}^{*} = 6 \\
		& x_{2}^{*} = 5 > 0 \stackrel{\text{COPD}}{\Longrightarrow} 9 y_{1}^{*} + 5 y_{2}^{*} + y_{3}^{*} = 4 \\
		& e_{2}^{*} = e_{3}^{*} = 0
		\end{aligned}
		\]
		En résolvant le système pour $y^{*}$, on obtient~:
		\[
		y_{1}^{*} = 0, \quad y_{2}^{*} = \dfrac{1}{3}, \quad y_{3}^{*} = \dfrac{7}{3}
		\]
	\end{exemple}
	
	\subsection{Calcul de la solution du problème dual à partir du tableau du simplexe du problème primal}
	
	On considère le problème de production~:
	\[
	\left\{
	\begin{array}{l}
	\displaystyle\max_{(x_{1}, x_{2})} \ f(x_{1}, x_{2}) = 6 x_{1} + 4 x_{2} \\
	\text{s.t.} \ 3 x_{1} + 9 x_{2} \leq 81 \\
	\quad 4 x_{1} + 5 x_{2} \leq 55 \\
	\quad 2 x_{1} + x_{2} \leq 20 \\
	x_{1}, x_{2} \geq 0
	\end{array}
	\right.
	\]
	Le tableau final du simplexe est~:
	\begin{center}
		\begin{tabular}{|c|rrrrrr|c|}
			\toprule
			Variables de base & $x_{1}$ & $x_{2}$ & $e_{1}$ & $e_{2}$ & $e_{3}$ & $z$ & $b$ \\
			\hline
			$e_{1}$ & 0 & 0 & 1 & $-\dfrac{15}{6}$ & $\dfrac{7}{2}$ & 0 & $\dfrac{27}{2}$ \\
			$x_{2}$ & 0 & 1 & 0 & $\dfrac{1}{3}$ & $-\dfrac{2}{3}$ & 0 & 5 \\
			$x_{1}$ & 1 & 0 & 0 & $-\dfrac{1}{6}$ & $\dfrac{5}{6}$ & 0 & $\dfrac{15}{2}$ \\
			\hline
			$z$ & 0 & 0 & 0 & $-\dfrac{1}{3}$ & $-\dfrac{7}{3}$ & -1 & -65 \\
			\hline
		\end{tabular}
	\end{center}
	La solution optimale est donc~:
	\[
	(x_{1}, x_{2}, e_{1}, e_{2}, e_{3})^{t} = \left( \dfrac{15}{2}, 5, \dfrac{27}{2}, 0, 0 \right)^{t}
	\]
	qui correspond à un bénéfice maximum de $z = 65$.
	
	Le dual s'écrit~:
	\[
	\left\{
	\begin{array}{ll}
	\displaystyle\min_{y} & g(y) = 81 y_{1} + 55 y_{2} + 20 y_{3} \\
	\text{s.t.} & 3 y_{1} + 4 y_{2} + 2 y_{3} \geq 6 \\
	& 9 y_{1} + 5 y_{2} + 1 y_{3} \geq 4 \\
	& y_{1}, y_{2}, y_{3} \geq 0
	\end{array}
	\right.
	\]
	Après application du simplexe au dual, on obtient~:
	\begin{center}
		\begin{tabular}{|c|rrrrrr|c|}
			\toprule
			Variables de base & $y_{1}$ & $y_{2}$ & $y_{3}$ & $\epsilon_{1}$ & $\epsilon_{2}$ & $-z$ & $b$ \\
			\midrule
			$y_{3}$ & $-\dfrac{7}{2}$ & 0 & 1 & $-\dfrac{5}{6}$ & $\dfrac{2}{3}$ & 0 & $\dfrac{7}{3}$ \\
			
			$y_{2}$ & $\dfrac{5}{2}$ & 1 & 0 & $\dfrac{1}{6}$ & $-\dfrac{1}{3}$ & 0 & $\dfrac{1}{3}$ \\
			\midrule
			$-z$ & $\dfrac{27}{2}$ & 0 & 0 & $\dfrac{15}{2}$ & 5 & 1 & -65 \\
			\bottomrule
		\end{tabular}
	\end{center}
	La solution optimale est~:
	\[
	(y_{1}, y_{2}, y_{3}, \epsilon_{1}, \epsilon_{2})^{t} = \left(0, \dfrac{1}{3}, \dfrac{7}{3}, 0, 0\right)^{t}
	\]
	On observe que la solution $y$ se trouve à la dernière ligne du tableau final du simplexe du problème primal. Les coefficients $c_{i}$ de la dernière ligne du tableau correspondent à~:
	\[
	\left\{
	\begin{array}{l}
	y_{j} = -c_{n+j}, \quad j = 1, \ldots, m \\
	\epsilon_{i} = -c_{i}, \quad i = 1, \ldots, n
	\end{array}
	\right.
	\]
	En appliquant sur le tableau primal, on obtient~:
	\[
	c_{3} = y_{1} = 0, \quad c_{4} = y_{2} = \dfrac{1}{3}, \quad c_{5} = y_{3} = \dfrac{7}{3}, \quad c_{1} = \epsilon_{1} = 0, \quad c_{2} = \epsilon_{2} = 0
	\]
	En parfaite dualité, on a les mêmes relations par rapport au tableau du problème dual~:
	\[
	\left\{
	\begin{array}{l}
	x_{i} = c_{m+i}, \quad i = 1, \ldots, n \\
	e_{j} = c_{j}, \quad j = 1, \ldots, m
	\end{array}
	\right.
	\]
	En appliquant sur le tableau dual, on obtient~:
	\[
	c_{4} = x_{1} = \dfrac{15}{2}, \quad c_{5} = x_{2} = 5, \quad c_{1} = e_{1} = \dfrac{27}{2}, \quad c_{2} = e_{2} = 0, \quad c_{3} = e_{3} = 0
	\]
\end{document}
